\chapter{Complete examples and synthesis}
\index{BFJ theory!examples}\index{Green function!non-Archimedean}
\label{chap:bfj-examples}

We close the BFJ part by carrying the main constructions through four
families: curves, toric varieties, atomic measures at divisorial points, and
simple degenerations.  The examples also locate the exact overlap between the
global BFJ theory of \Cref{part:bfj} and the local Chambert--Loir--Ducros
calculus of \Cref{part:cld}.  They are calculations, not analogies: every
normalization can be tested against total mass, constants, and the
two-component curve.

\section{Orientation}
\label{sec:ch31-orientation}

\subsection{Why a chapter of examples, and why here}
\label{sec:ch31-why}

\Crefrange{chap:bfj-model-forms}{chap:bfj-functoriality} construct an
apparatus of considerable size.  There is a space of closed model forms
(\Cref{def:ch20-closed-model-form}) with an exact sequence
(\Cref{thm:ch20-ddc-exact}); a class \(\PSH(X,\theta)\)
(\Cref{def:ch21-bfj-psh}) with a compactness theorem
(\Cref{thm:bfj-psh-compactness}) and a decreasing regularization theorem
(\Cref{thm:bfj-decreasing-regularization}); envelopes that are continuous
(\Cref{thm:bfj-envelope-regularity}); a Monge--Amp\`ere operator in four
incarnations (\Cref{sec:ch24-four-products}); a capacity
(\Cref{def:bfj-capacity}) with quasicontinuity
(\Cref{thm:bounded-quasicontinuity}); an energy
(\Cref{def:model-energy,def:extended-energy}) with a finite-energy class
(\Cref{def:e1}); a differentiability theorem
(\Cref{thm:energy-envelope-differentiability}) coupled to orthogonality
(\Cref{def:orthogonality-property}); a variational solution of the
Monge--Amp\`ere equation (\Cref{thm:bfj-supported-cy}); a continuity theorem
for its solutions (\Cref{cor:bfj-original-continuous-solution}); and a
functoriality package (\Cref{thm:base-change-ma}).

None of this is worth anything if the constants are wrong, and none of it can
be used with confidence until one has seen the operators evaluated on
objects one can hold in the hand.  This chapter is where that happens.  Its
role in \Cref{part:bfj} is the role of \Cref{chap:model-calculations} in
\Cref{part:foundations} and of \Cref{chap:cld-examples} in \Cref{part:cld}: not to
introduce anything, but to \emph{compute}.  Accordingly the chapter cites
heavily and restates almost nothing.  A reader who wants the statement of a
theorem used below will find it at the label; what is added here is the
number that comes out when the theorem is applied.

There is a second role, and it is the reason the chapter closes the part
rather than opening it.  The BFJ theory and the CLD theory are two different
constructions that produce, on an overlap, the same measures.  The overlap
was determined once, in \Cref{thm:ch19-cld-bfj-comparison}, and its
complement was tabulated once, in \Cref{sec:ch19-limitations}.  Every
theorem of \Cref{part:bfj} now has to be placed relative to that overlap:
does it have a CLD counterpart, on what locus, and with which implication
running in which direction?  \Cref{sec:ch31-table} answers that question
theorem by theorem.  It is the chapter's distinctive deliverable, and it is
the object that \Cref{chap:master-dictionary} in \Cref{part:comparisons}
will extend to the remaining four theories.

\subsection{The four families, and what each one tests}
\label{sec:ch31-families}

The families are not a random sample.  Each is chosen because a different
part of the apparatus becomes computable on it, and because a different
hypothesis becomes visible.

\begin{center}
\small
\begin{tabular}{p{.17\textwidth}|p{.25\textwidth}|p{.24\textwidth}|p{.24\textwidth}}
\toprule
\textbf{Family}&\textbf{What collapses}&\textbf{What stays hard}
 &\textbf{Hypothesis made visible}\\
\midrule
curves,
 \Cref{sec:ch31-curves}
 &the operator becomes affine: \(\MA_\omega(u)=\omega+\Delta_\Gamma u\)
 &nothing on the skeleton; everything off it
 &retraction to a skeleton is a psh operation only in dimension one\\
toric varieties,
 \Cref{sec:ch31-toric}
 &positivity becomes convexity and \(\MA\) becomes \(n!\MA_{\mathbf R}\)
 &no equation is solved without an invariance hypothesis
 &torus invariance, \Cref{ex:ch31-noninvariant}\\
divisorial supports,
 \Cref{sec:ch31-divisorial}
 &the equation becomes a finite-dimensional concave programme
 &modelness of the solution, tied to finite generation
 &orthogonality, \Cref{def:orthogonality-property}\\
degenerations,
 \Cref{sec:ch31-degenerations}
 &canonical metrics are computed in closed form
 &nothing is canonical without a group action
 &normalization of the ground valuation under base change\\
\bottomrule
\end{tabular}
\end{center}

The first and second families overlap on \(\mathbf P^1\), the second and
fourth on totally degenerate abelian varieties, and the third meets all of
them.  Those overlaps are used as consistency checks: a quantity computed in
two ways must come out the same, and in \Cref{sec:ch31-mass-ledger} every
such double computation is recorded.

\subsection{Six checks performed on every example}
\label{sec:ch31-checks}

The following list is applied, silently, to every computation below.  It is
the same list as \Cref{sec:appJ-unit-tests}, specialized to
\Cref{part:bfj}.

\begin{enumerate}
\item\label{itm:ch31-check:mass}
 \emph{Total mass.}  The Monge--Amp\`ere measure of an \(\omega\)-psh
 potential has mass \(V=(\{\omega\}^n)\), not one
 (\Cref{prop:model-ma-basic}, \Cref{thm:ch24-mass}).  Every example below
 computes \(V\) twice: once as an intersection number on \(X\), once as the
 sum of the masses of the atoms or the integral of the density.
\item\label{itm:ch31-check:constant}
 \emph{Constants.}  Adding \(c\in\mathbf R\) to a potential leaves the
 measure unchanged and adds \(cV\) to the energy
 (\eqref{eq:bfj-energy-translation}).  This is the fastest way to detect a
 missing factor of \(V\).
\item\label{itm:ch31-check:sign}
 \emph{Sign of the Dirichlet term.}  The quadratic part of the energy is
 nonpositive (\eqref{eq:model-dirichlet-negativity}); on a curve it is
 \(-\tfrac12\sum_e c_e\int_e|u'|^2\,\mathrm ds\) (\Cref{prop:ch26-graph-energy}).  A positive
 quadratic term means the sign of \(\ddc\) has been taken backwards.
\item\label{itm:ch31-check:vertical}
 \emph{Vertical divisor test.}  An effective vertical divisor gives a
 nonnegative model function and \(\varphi_{\cX_s}=1\)
 (\eqref{eq:appJ-model-divisor}, \eqref{eq:appJ-special-fibre-test}).
\item\label{itm:ch31-check:factorial}
 \emph{The factor \(n!\).}  In every toric or tropical comparison, the
 non-Archimedean measure is \(n!\) times the real Alexandrov measure
 (\eqref{eq:appJ-alexandrov}); the factor is invisible for \(n=1\) and is
 the commonest error for \(n\geq2\).
\item\label{itm:ch31-check:normalization}
 \emph{Ground valuation.}  Skeletal lengths and potential values scale with
 the normalization \(v(\pi)=1\); intersection numbers do not
 (\Cref{sec:ch31-base-change}).
\end{enumerate}

\begin{warning}{Two masses, two normalizations}
Nothing below is a probability measure unless the word \emph{probability}
appears.  In the toric family the reader will meet
\(n!\vol_M(P)=(L^n)\) and in the abelian family a Haar measure of total mass
one multiplied by \((L^g)\); these are the same convention, differently
displayed.  Dividing by \(V\) is legitimate only because \(V>0\), which is
where ampleness enters; see \Cref{ex:ch31-nef-not-big}.
\end{warning}

\subsection{Standing hypotheses and normalizations}
\label{sec:ch31-hypotheses}

Except where a departure is announced, we work under the original BFJ
hypotheses \eqref{eq:ch19-bfj-hypotheses} recalled in
\Cref{chap:bfj-model-forms}: \(K\) is complete and discretely valued with
\(v(\pi)=1\) and residue characteristic zero, \(X\) is smooth, connected and
projective over \(K\) of dimension \(n\), and
\(\theta\in\mathcal Z^{1,1}(X)\) has ample generic class
\(\{\theta\}\in N^1(X)\).  When a specific line bundle is in play we write
\(\omega\) for a closed model form representing \(c_1(L)\) and
\begin{equation}
              \|s\|_\varphi=\rme^{-\varphi}\|s\|_0,
              \qquad \MA_\omega(\varphi)(X^{\An})
                       =V=(\{\omega\}^n).
\label{eq:examples-global-normalization}
\end{equation}
The first half of \eqref{eq:examples-global-normalization} is the additive
weight convention of \Cref{sec:appJ-metrics}; the second is the mass
convention of \Cref{sec:appJ-model-ma}.  Both are binding for the whole
chapter.

Three of the four families require a temporary enlargement of the
hypotheses, and each enlargement is flagged where it occurs.  Curves over an
arbitrary complete nontrivially valued field are treated in
\Cref{sec:ch31-curves} because the continuity theorem
\Cref{thm:curves-arbitrary-field-continuity} is available for ample
metric envelopes and skeletal-supported measures there; the toric
comparison \Cref{thm:toric-bfj-alexandrov} is stated for a projective toric
variety, which is smooth only for a smooth fan; and totally degenerate
abelian varieties in \Cref{sec:ch31-abelian} are treated over an
algebraically closed field so that Raynaud uniformization is split.  Nothing
in \Cref{sec:ch31-table} depends on these enlargements: the table is stated
under \eqref{eq:ch19-bfj-hypotheses} throughout.

\subsection{What is imported}
\label{sec:ch31-imported}

The chapter proves no result of general theory.  Its imports are of three
kinds.

\emph{From \Cref{part:bfj}.}  The model measure
\eqref{eq:model-mixed-measure}, its mass \eqref{eq:mixed-total-mass}, the
monotone extension \Cref{thm:bounded-ma-extension}, the energy
\eqref{eq:model-energy} and its cocycle \eqref{eq:energy-cocycle},
orthogonality \Cref{def:orthogonality-property} and its provenance
\Cref{thm:orthogonality-provenance}, the variational solution
\Cref{thm:bfj-supported-cy}, and the continuity theorems of
\Cref{chap:bfj-continuity}.

\emph{From \Cref{part:foundations} and \Cref{part:cld}.}  The graph Laplacian and
its Green kernel (\Cref{prop:graph-green}), the slope formula
(\Cref{thm:ch08-slope-formula}), the toric dictionary
(\Cref{sec:toric-convex}), the Chambert--Loir measure of a model metric on a
curve (\Cref{prop:ch19-curve-model-curvature}), the toric Monge--Amp\`ere
comparison (\Cref{thm:ch19-toric-ma}), the canonical metric on an abelian
variety (\Cref{thm:ch19-canonical-ma}), and above all the comparison theorem
\Cref{thm:ch19-cld-bfj-comparison} with its scope clause and the limitations
table \Cref{sec:ch19-limitations-table}.

\emph{From the appendices.}  Convexity, Legendre duality, mixed volumes and
the Alexandrov measure (\Cref{app:convexity}); intersection numbers on
models (\Cref{app:intersection-bdivisors}); and the sign and mass
normalizations of \Cref{app:normalizations}, which are binding.

\subsection{Plan of the chapter}
\label{sec:ch31-plan}

\Cref{sec:ch31-curves} treats curves, where the operator is affine and the
equation is a Poisson equation on a metric graph; it ends with an atom at a
divisorial point that is not on any toric skeleton, an example reused twice
later.  \Cref{sec:ch31-toric} treats toric varieties, computing Alexandrov
measures in coordinates for polyhedral and for smooth potentials, and
exhibiting a semipositive model metric on \(\mathbf P^1\) that shows the
invariance hypothesis cannot be dropped.  \Cref{sec:ch31-divisorial} treats
measures supported at finitely many divisorial points, where the equation
becomes a finite-dimensional concave programme whose multipliers are the
masses.  \Cref{sec:ch31-degenerations} treats good reduction, nodal
degenerations, the Tate curve and totally degenerate abelian varieties.
\Cref{sec:ch31-table} is the theorem-by-theorem BFJ--CLD compatibility
table, \Cref{sec:ch31-diagnostic} the diagnostic list, and
\Cref{sec:ch31-synthesis} the synthesis of \Cref{part:bfj}.

\section{Curves: the nonlinear equation becomes linear}
\label{sec:ch31-curves}

% ISOLATED REVISED2 CANDIDATE. See GATE.md; all new bytes remain review-gated.

\subsection{The skeleton, its lengths, and the Laplacian}
\label{sec:ch31-curve-setup}

In this subsection and through \cref{rem:ch31-curve-solution-scope},
\(K\) is any complete nontrivially valued non-Archimedean field.
Let \(X\) be a smooth projective geometrically connected \(K\)-curve,
and suppose a proper strictly semistable formal model \(\cX\) over
\(K^\circ\) is given.  Write \(\Gamma\) for its finite connected
skeleton and \(p_{\cX}:X^{\An}\to\Gamma\) for its retraction.
No assertion that every curve has such a model over the original
\(K\) is implicit in this hypothesis.

An edge \(e\), corresponding to a node \(p_e\) of the special fibre,
has its analytic modulus length \(\ell_e>0\) and conductance
\[
                    c_e=[\kappa(p_e):\widetilde K].
\]
The conductance is the same at both ends, including both incidences
of a loop, and is retained under subdivision
\cite[\S2.2.3, p.~41]{Thu05}.
Lengths are measured by the chosen absolute value, using the natural
logarithm: an annulus \(r<|t|<R\) has length \(\log(R/r)\).
Where an \'{e}tale node chart has equation \(xy=\alpha\), this is
\(-\log|\alpha|\).  A node need not be rational over \(\widetilde K\);
one must use its \'{e}tale chart, not silently replace its completed
local ring by a rational-node equation.  In the discretely valued
examples with \(-\log|\pi|=1\), \(xy=\pi^n\) gives length \(n\).
Regularity of that total-space node forces \(n=1\); general
semistable formal models need not be regular before resolution.

For a real continuous piecewise affine function \(F\) on \(\Gamma\),
define
\begin{equation}
\label{eq:ch31-graph-laplacian}
 \Delta_\Gamma F
   :=\sum_{x\in\Gamma}
       \left(\sum_{\nu\in T_x\Gamma}
              c_{e(\nu)}\partial_\nu F(x)\right)\delta_x,
 \qquad
 \mathrm d\mathrm d^c(F\circ p_{\cX})=\iota_*\Delta_\Gamma F ,
\end{equation}
where \(\iota:\Gamma\hookrightarrow X^{\An}\), and
\(\partial_\nu\) is the outgoing derivative in the analytic length
coordinate.  The second equality is
\cite[Theorem~3.2.10(ii)]{Thu05}; properness removes its boundary
exception.  It also says that \(F\circ p_{\cX}\) is Thuillier-lisse.
It does not say that this function is a rational model function.
That stronger assertion requires the slope, breakpoint, and value-group
compatibilities of \cite[Proposition~3.5]{GJKM19}.

Edgewise integration by parts, counting each edge once, gives
\begin{equation}
\label{eq:graph-dirichlet-sign}
 \Delta_\Gamma F(\Gamma)=0,\qquad
 \int_\Gamma F\,\mathrm d\Delta_\Gamma F
       =-\sum_e c_e\int_e|F'|^2\,\mathrm ds\leq0.
\end{equation}
The two endpoint terms of every edge cancel in the total mass.
For the second identity they combine with the interior slope jumps
to give the displayed integral; equivalently use the bilinear
identity proved in \cref{lem:ch31-green-kernel}.
The weights are essential even when \(X\) is geometrically connected.
For example, if \(F\) is affine on an edge from \(v\) to \(w\), its
slope is \((F(w)-F(v))/\ell_e\).  Its energy contribution is
\[
             \frac{c_e}{\ell_e}\bigl(F(w)-F(v)\bigr)^2.
\]
Thus \(c_e\) is the weight in the differential energy, while
\(c_e/\ell_e\) is the conductance of the entire edge in the finite
vertex matrix.  Its reciprocal is the resistance length
\(\ell_e/c_e\).

\begin{remark}[Lengths, subdivisions, and attached branches]
\label{rem:ch31-lengths}
The combinatorial graph depends on the model.  An admissible
semistable refinement can subdivide edges and can also attach new
branches.  More precisely, the old skeleton embeds in the new one,
their retractions commute, each new connected piece attaches at one
point, and the weighted Laplacian pushes forward under the induced
retraction \cite[Proposition~2.2.26]{Thu05}.
A subdivision preserves the analytic length coordinate.  The lengths
of the new pieces sum to the original edge length, and every new
piece inherits the original conductance.  The resistance length
\(\ell_e/c_e\) used for the estimates below is an auxiliary metric;
it does not replace the analytic modulus length in the slope formula.

A model function determined on a dominating model can vary on an
attached branch.  It need not factor through the original \(p_{\cX}\).
One must either impose that factorization explicitly or compute on
a skeleton that actually determines the function and the reference.
A common semistable model can be used after the required scalar
extension; its existence over the original \(K\) is not presumed.
\end{remark}

\subsection{The Monge--Amp\`ere operator is affine}
\label{sec:ch31-curve-affine}

Let \(L\) have degree \(V>0\).  In this subsection assume an actual
relatively nef formal line bundle \(\cL\) on the given \(\cX\)
extends \(L^{\otimes m}\), for an integer \(m>0\).
Let \(\lVert\cdot\rVert_0\) be its rational formal metric and
\(\omega=c_1(L,\lVert\cdot\rVert_0)\).
This is additional data, not a consequence of the existence of
\(\cX\).  Put \(\mu_0=\MA_\omega(0)\).

\begin{lemma}
\label{lem:ch31-curve-ma-affine}
The reference measure is
\[
 \mu_0=\sum_v d_v\delta_v,\qquad
 d_v=\frac1m\deg_{\widetilde K}(\cL|_{C_v}),\qquad
 \sum_vd_v=V.
\]
It is positive and supported on the original vertex set.
For every rational model function of the form
\(u=F\circ p_{\cX}\), where \(F\) is real piecewise affine,
\begin{equation}
\label{eq:curve-ma-laplacian}
                 \MA_\omega(u)=\mu_0+\iota_*\Delta_\Gamma F.
\end{equation}
Here for a model metric that is not nef the left side means its signed
curvature measure.  The function \(u\) is \(\omega\)-psh in the
model sense if and only if \(\mu_0+\Delta_\Gamma F\geq0\).
For a measure \(\mu\) supported on \(\Gamma\), the equation becomes
\begin{equation}
\label{eq:curve-poisson}
                      \Delta_\Gamma F=\mu-\mu_0 .
\end{equation}
Measures on \(\Gamma\) are identified with their pushforwards on
\(X^{\An}\) when no confusion is possible.

More generally, let \(\theta\) be a real closed model form determined
on \(\cX\), let \(d_x(\theta)\) be its component degree over
\(\widetilde K\) at an original vertex, and put \(d_x(\theta)=0\)
elsewhere.  For an actual rational model function
\(F\circ p_{\cX}\), the numerical criterion is
\begin{equation}
\label{eq:curve-slope-criterion}
 \theta+\mathrm d\mathrm d^c(F\circ p_{\cX})\text{ is nef}
 \quad\Longleftrightarrow\quad
 d_x(\theta)+\sum_{\nu\in T_x\Gamma}
       c_{e(\nu)}\partial_\nu F(x)\geq0
 \quad\text{for every }x\in\Gamma.
\end{equation}
Only the original vertices and slope breaks impose nonzero conditions.
\end{lemma}

\begin{proof}
For the reference model the curvature mass at \(v\) is its degree
over the base residue field divided by \(m\).
In \cite[Theorem~18.5.5 and Proposition~18.5.6(a)--(c)]{CLD25},
the generic smoothness and geometrically reduced strictly semistable
special fibre make the additional multiplicity, defect, and graded
residue factors equal to one.  Total mass is \(V\) by
\cite[Theorem~18.3.11]{CLD25}.
If a component has constant field \(k_v\), its degree over
\(\widetilde K\) is \([k_v:\widetilde K]\) times its degree over
\(k_v\); the latter cannot be substituted without this factor.

The formal weights are Thuillier-lisse by the tensor-power and
analytic-unit argument in the proof of
\cref{thm:curves-arbitrary-field-continuity}.
The original-field model comparison there identifies Thuillier and
CLD curvature, so additivity and \eqref{eq:ch31-graph-laplacian}
prove \eqref{eq:curve-ma-laplacian}.
For the positivity criterion, a nef model metric has positive
curvature.  Conversely, positive curvature makes each lisse local
weight Thuillier-subharmonic by \cite[Theorem~3.4.12(ii)]{Thu05}.
This remains true after extension to \(\widehat{K^{\mathrm a}}\)
by \cite[Corollary~3.4.5]{Thu05}.  There one may choose a common
semistable model actually determining the pulled-back metric; the
slope and component-degree formula
\cite[Proposition~3.7]{GJKM19} says all its component degrees are
nonnegative.  Thus that model metric is nef.
Model-independence and reflection of nefness under scalar extension
\cite[Proposition~3.5 and Lemma~3.3]{GM19} give nefness over \(K\).
Finally \eqref{eq:curve-poisson} is
\eqref{eq:curve-ma-laplacian} rearranged.

We spell out the reverse implication for the real-form criterion
\eqref{eq:curve-slope-criterion}, since positivity cannot be recovered
from the positivity of a scalar pushforward.
Choose a representative of \(\theta\) on its actual determining
model \(\cX\) as a finite real linear combination
\[
       \theta=\sum_{i=1}^N r_i
                   c_1(M_i,\lVert\cdot\rVert_i),
       \qquad r_i\in\mathbf R,
\]
where the \(\lVert\cdot\rVert_i\) are rational formal metrics
determined on \(\cX\).  Equivalently, the representative is a finite
real combination of model line-bundle classes; tensor-power
denominators may be absorbed into the \(r_i\).
This representation does not turn \(\theta\) into the Chern class
of one line bundle.

On an analytic open set \(W\) where all \(M_i\) have invertible frames \(e_i\),
put \(h_i=-\log\lVert e_i\rVert_i\) and
\[
                  H=\sum_{i=1}^N r_i h_i+F\circ p_{\cX}.
\]
Each \(h_i\) is Thuillier-lisse by the tensor-power and analytic-unit
argument already used, and \(F\circ p_{\cX}\) is lisse by
\eqref{eq:ch31-graph-laplacian}.  Thus \(H\) is a real lisse function,
even when some \(r_i\) are negative or irrational.
The component-degree and outgoing-slope identities extend
real-linearly, giving on this open set
\begin{equation}
\label{eq:ch31-real-local-curvature}
 \mathrm d\mathrm d^cH
  =\left(\sum_{v\text{ an original vertex}}
                 d_v(\theta)\delta_v
                       +\iota_*\Delta_\Gamma F\right)\bigg|_W.
\end{equation}
The real linearity proves this equality of signed measures;
it is not by itself an argument about positivity.

These local calculations are independent of the choices.
Changing a frame \(e_i\) adds a real multiple of the logarithm
of an analytic unit to \(H\); that logarithm is harmonic by the
Poincar\'e--Lelong formula, so its
\(\mathrm d\mathrm d^c\) vanishes.
Changing the line-bundle representative of the same numerical class
on \(\cX\) leaves every component degree \(d_v(\theta)\) unchanged.
Equation \eqref{eq:ch31-real-local-curvature} therefore gives the
same local curvature current.  On overlaps the resulting weights
differ by a harmonic function, by \cite[Lemma~3.3.12]{Thu05}.
In particular subharmonicity is unaffected.  We do not need to
glue the \(H\) into a global function or an actual metrized line bundle.

Now suppose the weighted inequalities in
\eqref{eq:curve-slope-criterion} hold.  They make the right side of
\eqref{eq:ch31-real-local-curvature} positive, hence make \(H\)
subharmonic by \cite[Theorem~3.4.12(ii)]{Thu05}.
Put \(\Omega=\widehat{K^{\mathrm a}}\) and let
\(\rho:X_\Omega^{\An}\to X^{\An}\) be the scalar projection.
The functions \(H\circ\rho\) are subharmonic by
\cite[Corollary~3.4.5]{Thu05}.
Over \(\Omega\), choose a common strictly semistable model actually
determining the finitely many pulled-back model metrics and the
rational model function \(F\circ p_{\cX}\).
On this model the curvature of \(H\circ\rho\), computed real-linearly,
has as its vertex coefficients the individual component degrees of
\(\theta_\Omega+\mathrm d\mathrm d^c(F\circ p_{\cX}\circ\rho)\);
this is \cite[Proposition~3.7]{GJKM19}.
Its positivity implies that each of those component degrees is
nonnegative.  The pulled-back real numerical model class is therefore
nef.  The degree and projection-formula argument in the proof of
\cref{thm:ch29-real-class-curve-model-envelope} reflects nefness
of real model classes through scalar extension, so the original
class is nef over \(K\).

Conversely, if the original real model class is nef, its pullback
is nef.  The upstairs component-degree formula gives positive
local curvature measures.  Their ordinary scalar pushforwards
are \(\mathrm d\mathrm d^cH\) by
\cite[Theorem~3.2.10(i)]{Thu05}, so
\eqref{eq:ch31-real-local-curvature} is positive and yields the
weighted inequalities.  This direction uses pushforward only to
preserve positivity, not to reflect it.
The reverse implication above instead uses pullback of
subharmonic local weights before reading the individual upstairs
degrees.  No positivity of individual terms is inferred from a
positive sum of lifted degrees.
Throughout, \(F\circ p_{\cX}\) is an actual rational model function
with the stated factorization; domination alone never supplies it.
\end{proof}

\begin{example}[A nonsplit node and its two degree factors]
\label{ex:ch31-nonsplit-weight-normalization}
Let \(R\) be a complete discrete valuation ring with fraction field \(K\)
and residue field \(k\), normalized by \(-\log|\pi|=1\).
Choose a closed point \(p\in\mathbf P^1_k\) whose residue field
\(k'/k\) is finite separable of degree \(r>1\), and blow up
\(\mathbf P^1_R\) at \(p\).  Complete this model formally along its
special fibre.  Its generic fibre is still \(\mathbf P^1_K\), so it
is geometrically connected.

The special fibre has two multiplicity-one components: the strict
transform \(D\simeq\mathbf P^1_k\) and the exceptional component
\(E\simeq\mathbf P^1_{k'}\).  They meet transversely at a node of residue
field \(k'\).  To see the local thickness, choose an \'{e}tale
parameter \(t\) at \(p\).  The blowup chart containing their
intersection has equation \(\pi=t z\); its two special-fibre
branches are \(t=0\) and \(z=0\).  Thus the skeleton is a segment
of analytic length one with endpoints \(v_D,v_E\), and its single
edge has weight \(c_e=r\).

Consider \(\mathcal O(E)\), which is trivial on the generic fibre.
Since \(E\) cuts out the degree-\(r\) point \(p\) on \(D\),
\[
                 \mathcal O(E)|_D=\mathcal O_D(p),
                 \qquad \deg_k\bigl(\mathcal O(E)|_D\bigr)=r.
\]
The tautological description of the exceptional divisor in a blowup
gives \(\mathcal O(-E)|_E=\mathcal O_{\mathbf P^1_{k'}}(1)\).
Consequently
\[
 \mathcal O(E)|_E=\mathcal O_{\mathbf P^1_{k'}}(-1),
 \qquad
 \deg_{k'}\bigl(\mathcal O(E)|_E\bigr)=-1,
 \qquad
 \deg_k\bigl(\mathcal O(E)|_E\bigr)=-r.
\]
The model metric on the trivial generic line bundle has potential
\(F=-\log\lVert1\rVert\).  On the displayed chart a frame of
\(\mathcal O(E)\) is \(t^{-1}\), so
\(1=t(t^{-1})\) gives \(F=-\log|t|\).
It follows that \(F(v_D)=0\), \(F(v_E)=1\), and \(F\) is affine
with slope \(+1\) from \(v_D\) to \(v_E\).  The outgoing slope
at the other endpoint is \(-1\).  Therefore
\begin{equation}
\label{eq:ch31-nonsplit-normalization-test}
       \Delta_\Gamma F=r\delta_{v_D}-r\delta_{v_E},
\end{equation}
exactly the two component degrees over \(k\) computed above.
The constant field of \(D\) is \(k\), while that of \(E\) is \(k'\);
nevertheless both slope terms use the same node residue degree \(r\).
The Dirac measures in \eqref{eq:ch31-nonsplit-normalization-test}
are unit Diracs, and the total mass is zero, as required by the
trivial generic line bundle.

After an unramified scalar extension whose residue field splits
\(k'/k\), the exceptional component splits into \(r\) components,
each meeting the lifted \(D\) in a weight-one edge.  Their ordinary
pushforward recovers the coefficient \(r\) of the original edge,
in agreement with \eqref{eq:ch29-curve-scalar-weights}.
This is a signed-curvature normalization test:
\(\mathcal O(E)\) is not nef because its degree on \(E\) is negative.
No semipositive degree-zero conclusion is asserted.
\end{example}

\begin{warning}{Affine, not linear}
Equation \eqref{eq:curve-ma-laplacian} is affine in the potential,
with linear part \(\Delta_\Gamma\) and constant term \(\mu_0\).
The set of continuous semipositive metric weights relative to
\(\lVert\cdot\rVert_0\) is a convex set, not in general a cone or a
vector space: the constraint is
\(\mu_0+\Delta_\Gamma F\geq0\).
Within the rational model class, taking an arbitrary real convex
combination need not remain a rational model function.
\end{warning}

\subsection{Solving the equation: the graph Green kernel}
\label{sec:ch31-curve-green}

The following graph lemma uses arbitrary positive conductances and
does not depend on any analytic existence or envelope theorem.

\begin{lemma}
\label{lem:ch31-green-kernel}
Let \(\Gamma\) be a finite connected metric graph.  Give each edge \(e\)
a positive real constant \(c_e\), retained on every subdivision of that
edge.  Use the weighted outgoing-slope Laplacian
\eqref{eq:ch31-graph-laplacian}.
Let \(q\in \Gamma\).  There is a unique symmetric continuous kernel
\(g_q:\Gamma\times \Gamma\to\mathbf R\) for which
\begin{equation}
\label{eq:ch31-green-kernel}
 \Delta_{\Gamma,x}g_q(x,y)=\delta_y-\delta_q,
 \qquad g_q(q,y)=0,
\end{equation}
and \(g_q(\cdot,y)\) is affine outside the graph vertices and
\(\{q,y\}\).  If \(d_R\) is the path metric assigning length
\(\ell_e/c_e\) to \(e\), and \(D_R=\operatorname{diam}_{d_R}\Gamma\), then
\begin{equation}
\label{eq:ch31-green-kernel-estimates}
 |g_q(x,y)|\leq D_R,
 \qquad
 |g_q(x,y)-g_q(x',y)|\leq\sqrt{D_Rd_R(x,x')}.
\end{equation}
For a finite signed Radon measure \(\nu\) of mass zero,
\begin{equation}
\label{eq:ch31-green-potential}
 U_\nu(x)=\int_\Gamma g_q(x,y)\,\mathrm d\nu(y)
\end{equation}
is the unique continuous solution of \(\Delta_\Gamma U_\nu=\nu\) with
\(U_\nu(q)=0\).  Distributional equality here means
\(\int_\Gamma U_\nu\,\mathrm d\Delta_\Gamma F=\int_\Gamma F\,\mathrm d\nu\)
for every continuous piecewise affine \(F\).
If \(\nu_j\to\nu\) weakly and their total variations are uniformly
bounded, then \(U_{\nu_j}\to U_\nu\) uniformly.
\end{lemma}


\begin{proof}
Subdivide to include all breakpoints.  Edgewise integration by parts gives
\begin{equation}
\label{eq:ch31-weighted-green-energy}
 \int_\Gamma F\,\mathrm d\Delta_\Gamma H
   =-\sum_e c_e\int_e F'H'\,\mathrm ds
   =\int_\Gamma H\,\mathrm d\Delta_\Gamma F.
\end{equation}
In particular \(\Delta_\Gamma F(\Gamma)=0\).  An edge is counted once in the
energy and both of its tangent directions are counted in the vertex
sum, including for loops.  This is the weighted version of
\cite[Proposition~1.2.13]{Thu05}; it also follows immediately by
cancelling the two endpoint terms of each edge.

Fix \(y\), and insert \(q,y\) as vertices.  On functions affine on
each resulting edge, minus the Laplacian has quadratic form
\[
                 \sum_{e=(v,w)}\frac{c_e}{\ell_e}
                                      (a_v-a_w)^2.
\]
The associated matrix is symmetric; its kernel consists exactly of
constants by connectedness.  Its image is consequently the vectors
whose coordinate sum is zero.  Solve for \(\delta_y-\delta_q\),
and normalize the solution to vanish at \(q\).  Further subdivisions
do not change it, since the affine extension has zero mass at each
new point.  If \(\Gamma\) is a point, the solution is identically zero.

Suppose now that \(q\ne y\).  The maximum principle for this finite
weighted matrix says that the maximum of \(h=g_q(\cdot,y)\) occurs
at \(q\) and its minimum at \(y\).  Indeed at a maximum every
outgoing slope is nonpositive.  If the Laplacian there is zero,
positivity of every conductance forces all those slopes to vanish;
propagate the maximum along edges until meeting a nonzero mass.
A positive mass cannot occur at a maximum.  The argument for a
minimum is the same with signs reversed.  Thus \(h(q)=0\),
\(h\leq0\), and \eqref{eq:ch31-weighted-green-energy} gives
\[
 E(h):=\sum_e c_e\int_e|h'|^2\,\mathrm ds=-h(y)
                                         =\operatorname{osc}(h).
\]
For any edge-simple path \(\gamma\) from \(a\) to \(b\), after
subdivision at its endpoints, Cauchy--Schwarz gives
\[
 |h(a)-h(b)|^2
 \leq\left(\int_\gamma\frac{\mathrm ds}{c_e}\right)
       \left(\int_\gamma c_e|h'|^2\,\mathrm ds\right)
 \leq\operatorname{length}_R(\gamma)E(h).
\]
The second inequality uses that no edge segment is traversed twice.
Shortest resistance-length paths may be chosen edge-simple, since all
edge lengths are positive.  Take such a path from \(q\) to \(y\) to
obtain \(\operatorname{osc}(h)\leq d_R(q,y)\leq D_R\), then use
a shortest path from \(x\) to \(x'\).  This proves
\eqref{eq:ch31-green-kernel-estimates}.

Apply \eqref{eq:ch31-weighted-green-energy} to
\(g_q(\cdot,y)\) and \(g_q(\cdot,z)\), on a subdivision containing
both \(y,z\).  The normalizing values at \(q\) vanish, giving
\(g_q(z,y)=g_q(y,z)\).  Symmetry and the uniform estimate already
proved give the same modulus of continuity in the second variable.
Hence the kernel is jointly continuous.  This avoids assuming
continuous dependence of a moving-vertex matrix at a degenerating edge.

The integral in \eqref{eq:ch31-green-potential} is therefore
continuous.  For piecewise affine \(F\), Fubini and
\eqref{eq:ch31-weighted-green-energy} yield
\[
 \int_\Gamma U_\nu\,\mathrm d\Delta_\Gamma F
  =\int_\Gamma\left(F(y)-F(q)\right)\mathrm d\nu(y)
  =\int_\Gamma F\,\mathrm d\nu.
\]
Both integrations are against finite signed measures and a bounded
continuous kernel, so Fubini applies.

For uniqueness in the asserted \emph{continuous} class, do not assume
an energy integral before proving piecewise affineness.  If \(w\) is
continuous and distributionally harmonic, test with the piecewise
affine tent on a closed subinterval \([a,b]\) of an open edge, of
height one at \(x\in(a,b)\) and zero outside \((a,b)\).
The equality \(\int w\,\mathrm d\Delta_\Gamma F=0\) becomes
\[
 \frac{w(a)}{x-a}+\frac{w(b)}{b-x}
       -w(x)\left(\frac1{x-a}+\frac1{b-x}\right)=0.
\]
It follows that \(w\) is affine on every open edge, and continuity
extends this to its endpoints.  Vertex tent tests then give the
weighted zero-slope sum at the vertices.  The energy identity now
applies and forces \(w\) to be constant.

Finally, the family \(\{g_q(x,\cdot):x\in \Gamma\}\) is compact in
\(C(\Gamma)\), being the continuous image of the compact graph.  Weak
convergence of measures with a uniform total-variation bound is uniform
on this compact family: approximate it by a finite sup-norm net and
bound the error by that total-variation bound.  This proves the last
assertion.
\end{proof}



\begin{example}[Rational masses need not give rational slopes]
\label{ex:ch31-rational-mass-irrational-slope}
Consider the finite metric graph with two vertices \(q,y\) joined by
two edges of lengths \(1\) and \(\sqrt{2}\), each of weight one.
This is a graph calculation; no realization over a particular
ground field is assumed.  Orient both edges from \(q\) to \(y\),
and seek the normalized potential
\(\Delta_\Gamma U=\delta_y-\delta_q\), \(U(q)=0\).
It must be affine on both open edges by
\cref{lem:ch31-green-kernel}.  If \(A=U(y)\), continuity requires
the two slopes to be \(A\) and \(A/\sqrt{2}\).
The mass at \(q\) is their sum, hence
\[
        A+\frac{A}{\sqrt{2}}=-1,
        \qquad
        A=\sqrt{2}-2,\qquad
        \frac{A}{\sqrt{2}}=1-\sqrt{2}.
\]
The outgoing slopes at \(y\) have the opposite signs and sum to \(1\),
so these formulas satisfy the equation at both vertices.
Both slopes are irrational, although the two prescribed atomic
masses are integers.  Thus rational masses alone cannot justify
rational slopes; value-group compatibility of vertex values is
a further issue.  This example does not obstruct the continuous
Green solution or the separate rational formal approximation
argument for curve metrics.
\end{example}

\begin{proposition}
\label{prop:graph-ma-solution}
With the given model reference above, fix \(q\in\Gamma\), and let
\(\mu\) be any positive Radon measure on \(\Gamma\) of mass \(V\).  The function
\(F=U_{\mu-\mu_0}\) of \cref{lem:ch31-green-kernel} is the unique
continuous solution of \eqref{eq:curve-poisson} normalized by \(F(q)=0\).
The metric
\[
             \lVert\cdot\rVert_0\mathrm e^{-F\circ p_{\cX}}
\]
is continuous and semipositive, and its Thuillier, Chambert--Loir,
and CLD curvature measures all equal \(\mu\).
The solution is unique up to an additive constant in the weight,
or a positive constant scale in the metric.
\end{proposition}

\begin{proof}
The measure \(\nu=\mu-\mu_0\) has mass zero.  The graph lemma gives
the asserted continuous potential and its distributional uniqueness.
To check its analytic curvature without assuming rational slopes,
choose finite positive atomic measures \(\mu_j\) of mass exactly
\(V\) converging weakly to \(\mu\).  Partition the finitely many
edges into intervals of mesh tending to zero, place each interval's
mass at one point, and retain vertex masses; this gives such measures
with real, not necessarily rational, masses.
Then \(F_j=U_{\mu_j-\mu_0}\) are real piecewise affine and converge
uniformly to \(F\) by \cref{lem:ch31-green-kernel}.

The reference weights and \(F_j\circ p_{\cX}\) are Thuillier-lisse.
The slope formula gives curvature \(\mu_j\) for the corresponding
metrics.  Their weights are subharmonic, and the uniform limits
remain subharmonic by the harmonic comparison definition.
For a compactly supported lisse test function \(a\), the error in
curvature pairing is bounded by
\(\lVert F_j-F\rVert_\infty
  \lVert\mathrm d\mathrm d^ca\rVert_{\mathrm{TV}}\).
Thus the limiting metric has Thuillier curvature \(\mu\);
density of lisse tests \cite[Proposition~3.2.7]{Thu05} identifies
the measures.

The optional model-reference formula
\eqref{eq:ch29-curve-model-reference-solution} in
\cref{thm:curves-arbitrary-field-continuity} now applies to exactly
this metric.  Its proof promotes the continuous positive-curvature
metric to a uniform limit of semipositive rational formal metrics
by comparison over the completed algebraic closure and finite
normal descent, and identifies the three curvature measures.
This use is not circular: that proof uses only the graph lemma above,
not the present proposition.
Uniqueness of analytic metrics follows because the logarithm of
their ratio is a continuous harmonic function on the proper connected
curve, hence constant, as proved in
\cref{prop:ch29-curve-point-reference}.

Neither the \(F_j\) nor their pullbacks have been declared rational
model functions.  Rational atomic masses alone would not imply
rational slopes on a graph with arbitrary modulus lengths.
For a nonatomic measure the limiting potential is only asserted
continuous; it need not be Thuillier-lisse.
\end{proof}

\begin{remark}[What this proposition is, and is not]
\label{rem:ch31-curve-solution-scope}
The proposition solves the equation for every positive measure of
the right mass supported on this given finite skeleton.
The quantitative bound is
\(\sup|g_q|\leq D_R\), where \(D_R\) is the diameter for resistance
lengths \(\ell_e/c_e\); the analytic lengths remain \(\ell_e\).
Measures not supported on a finite skeleton are outside this
continuous-solution assertion.
No degree-zero rational-model existence claim, unrestricted envelope
identity, or higher-dimensional reduction is inferred.
\end{remark}


\subsection{The two-component curve, completely}
\label{sec:ch31-two-component}

For the explicit examples from here to \cref{rem:ch31-retraction},
we return to the discrete normalization \(-\log|\pi|=1\) and use
split node charts, so every displayed edge has conductance \(c_e=1\).
These are assumptions of the examples only, not restrictions on
\cref{prop:graph-ma-solution}.
In graph formulas \(\omega\) also denotes its reference measure \(\mu_0\).

\begin{example}[Two components, solved completely]
\label{ex:two-component-complete}
Suppose \(\Gamma=[x_0,x_1]\) has length one and
\(\omega=m_0\delta_{x_0}+m_1\delta_{x_1}\), with \(m_i>0\).  For vertex
values \(a=u(x_0)\), \(b=u(x_1)\),
\begin{equation}
 \MA_\omega(u)
 =(m_0+b-a)\delta_{x_0}+(m_1+a-b)\delta_{x_1}.
\label{eq:two-component-ma}
\end{equation}
To solve \(\MA_\omega(u)=(m_0+m_1)\delta_{x_0}\), take
\(a=-m_1\), \(b=0\).  The energy is
\begin{equation}
 E_\omega(u)
 =-m_0m_1-\frac12m_1^2.
\label{eq:two-component-energy}
\end{equation}
The coefficient at \(x_1\) vanishes, the coefficient at \(x_0\) is
\(V=m_0+m_1\), and adding a constant changes the energy by that constant
times \(V\).  These three checks fix the sign and mass normalization.
\end{example}

The example deserves to be worked out in full, because every later
normalization is tested against it.

\subsubsection*{The geometry behind the graph}

Let \(\cX\to\Spec R\) be regular semistable with
\(\cX_s=E_0+E_1\), the two components smooth and meeting transversely at a
single node whose completed local ring is \(R[[x,y]]/(xy-\pi)\).  Then
\(\Gamma\) is a segment of length one with endpoints \(x_{E_0},x_{E_1}\),
which we write \(x_0,x_1\).  Let \(\cL\) be a model of \(L\) with
\(\deg(\cL|_{E_i})=m_i>0\).  By \Cref{lem:ch31-curve-ma-affine},
\(\mu_0=m_0\delta_{x_0}+m_1\delta_{x_1}\) and
\(V=\deg L=m_0+m_1\), the last equality being the flatness computation
\(\deg(\cL|_{\cX_s})=\deg(\cL|_{E_0})+\deg(\cL|_{E_1})\).  This is the
first appearance of check~\cref{itm:ch31-check:mass}: the mass \(m_0+m_1\)
of \eqref{eq:two-component-ma} is an intersection number on \(X\).

\subsubsection*{The affine formula and the admissible region}

Identify \(\Gamma\) with \([0,1]\), with \(x_0\) at \(0\).  A model function
on \(\Gamma\) is a piecewise affine \(u\), and by
\eqref{eq:curve-ma-laplacian} its measure is
\begin{equation}
\label{eq:ch31-segment-general-ma}
 \MA_\omega(u)
 =\bigl(m_0+u'(0^+)\bigr)\delta_{x_0}
  +\bigl(m_1-u'(1^-)\bigr)\delta_{x_1}
  +\sum_{t\in(0,1)}\bigl(u'(t^+)-u'(t^-)\bigr)\delta_t ,
\end{equation}
because the outgoing slope at \(0\) is \(u'(0^+)\), the outgoing slope at
\(1\) is \(-u'(1^-)\), and at an interior point the two outgoing slopes sum
to the jump of \(u'\).  Hence \(u\) is \(\omega\)-psh precisely when
\begin{equation}
\label{eq:ch31-segment-psh}
 u\ \text{is convex on }(0,1),
 \qquad u'(0^+)\geq-m_0,
 \qquad u'(1^-)\leq m_1 .
\end{equation}
The three conditions are independent, and the total mass of
\eqref{eq:ch31-segment-general-ma} is
\(m_0+m_1+\bigl(u'(0^+)-u'(1^-)+\sum\text{jumps}\bigr)=m_0+m_1=V\), because
the bracket telescopes to zero.  That is
check~\cref{itm:ch31-check:mass} again, now for an arbitrary potential.

We now specialize to \(u\) affine, which by
\eqref{eq:ch31-segment-general-ma} is the case in which the measure is
carried by the two endpoints.  For \(u\) affine with
\(u(x_0)=a\), \(u(x_1)=b\) on an edge of length one,
the outgoing slope at \(x_0\) is \(b-a\) and at \(x_1\) is \(a-b\), so
\(\Delta_\Gamma u=(b-a)\delta_{x_0}+(a-b)\delta_{x_1}\) and
\eqref{eq:two-component-ma} follows from
\eqref{eq:curve-ma-laplacian}.  Positivity of both coefficients is
\begin{equation}
\label{eq:ch31-two-component-window}
 -m_0\leq b-a\leq m_1,
 \qquad\text{equivalently}\qquad
 |a-b|\leq\min\{m_0,m_1\}\ \text{when }m_0=m_1 .
\end{equation}
The general window \eqref{eq:ch31-two-component-window} is
\Cref{ex:two-vertex-ma}; the symmetric form is the one used in
\Cref{ex:ch26-two-vertex-energy}.

\subsubsection*{The solution and its three checks}

Solving \(\MA_\omega(u)=V\delta_{x_0}\) forces \(m_1+a-b=0\), so
\(b-a=m_1\), which is the right endpoint of the window
\eqref{eq:ch31-two-component-window}.  Normalizing \(b=0\) gives
\(a=-m_1\), as in \Cref{ex:two-component-complete}.  The three checks are:
\begin{enumerate}
\item the coefficient at \(x_1\) is \(m_1+a-b=m_1-m_1=0\);
\item the coefficient at \(x_0\) is \(m_0+b-a=m_0+m_1=V\);
\item replacing \((a,b)\) by \((a+c,b+c)\) changes neither coefficient and
 changes \(E_\omega\) by \(cV\).
\end{enumerate}
For the energy, \Cref{prop:ch26-graph-energy} gives
\(E_\omega(u)=\int u\,\rd\mu_0-\tfrac12\int|u'|^2\rd s\); with
\(u(x_0)=a\), \(u(x_1)=b\) and edge length one this is
\(m_0a+m_1b-\tfrac12(a-b)^2\), which is
\eqref{eq:ch26-two-vertex-energy}.  At \((a,b)=(-m_1,0)\) it equals
\(-m_0m_1-\tfrac12m_1^2\), which is \eqref{eq:two-component-energy}.

\begin{example}[The same computation with numbers]
\label{ex:ch31-two-component-numbers}
Take \(m_0=2\), \(m_1=3\), so \(V=\deg L=5\).  The admissible window
\eqref{eq:ch31-two-component-window} is \(-2\leq b-a\leq3\), a segment of
length \(V=5\) in the variable \(b-a\); this is not an accident, since the
two coefficients of \eqref{eq:two-component-ma} are nonnegative and sum to
\(5\).  Normalizing \(b=0\):
\begin{center}
\small
\begin{tabular}{c|c|c|c|c}
\toprule
\(a\)&\(\MA_\omega(u)(\{x_0\})\)&\(\MA_\omega(u)(\{x_1\})\)
 &\(E_\omega(u)\)&remark\\
\midrule
\(0\)&\(2\)&\(3\)&\(0\)&reference metric\\
\(-3\)&\(5\)&\(0\)&\(-\tfrac{21}2\)&all mass at \(x_0\)\\
\(+2\)&\(0\)&\(5\)&\(2\)&all mass at \(x_1\)\\
\(-1\)&\(3\)&\(2\)&\(-\tfrac52\)&interior point\\
\(-4\)&\(6\)&\(-1\)&---&not \(\omega\)-psh\\
\bottomrule
\end{tabular}
\end{center}
Every row has masses summing to \(5\).  The energy entries use
\(E_\omega(a,0)=2a-\tfrac12a^2\); at \(a=-3\) this is
\(-6-\tfrac92=-\tfrac{21}2\), matching \eqref{eq:two-component-energy} with
\(m_0=2,m_1=3\): \(-6-\tfrac92\).  At \(a=+2\) it is \(4-2=2\), matching the
mirror-image formula \(-m_0m_1-\tfrac12m_0^2+\) the normalization shift; the
direct check is easier: \(E_\omega(2,0)=2\cdot2-\tfrac12\cdot4=2\).  The
maximum of \(a\mapsto2a-\tfrac12a^2\) over the admissible window
\([-3,2]\) is at \(a=2\), the right endpoint, with value \(2\); the
unconstrained maximum of the quadratic is also at \(a=2\), where the
derivative \(2-a\) vanishes, and indeed \(\rd E/\rd a=\MA(\{x_0\})\)
vanishes there.  That is
\eqref{eq:energy-first-variation} in the smallest possible instance.
\end{example}

\begin{remark}[Reading the derivative of the energy]
\label{rem:ch31-two-component-derivative}
In \Cref{ex:ch31-two-component-numbers} the derivative of the energy in the
direction \(\partial/\partial a\) is the mass of the Monge--Amp\`ere measure
at \(x_0\).  This is the whole content of
\Cref{thm:e1-energy-differentiability} in dimension one and with a
one-parameter family of model directions: the direction is the model
function \(\chi\) with \(\chi(x_0)=1\), \(\chi(x_1)=0\), and
\(\frac{\rd}{\rd t}E_\omega(u+t\chi)=\int\chi\,\MA_\omega(u+t\chi)\).
The reader who wants the general statement with all its hypotheses will
find it in \Cref{sec:ch27-e1-segment}; the reader who wants to remember it
should remember this table.
\end{remark}

\subsection{A chain of four components, solved for four measures}
\label{sec:ch31-chain}

The two-component curve is too small to distinguish an atomic measure from a
diffuse one.  The next example is the smallest one that does.

\begin{example}[Inverting the Monge--Amp\`ere operator on a chain]
\label{ex:ch31-chain-solve}
Return to \Cref{ex:ch20-chain-three}: \(\cX_s\) is a chain
\(E_0,E_1,E_2,E_3\) of four rational curves, all of multiplicity one,
\(\Gamma=[0,3]\) subdivided at the integers with all edges of length one,
the degrees of \(\cL\) are \((1,0,0,1)\), so
\[
 \mu_0=\delta_0+\delta_3,\qquad V=\deg L=2 .
\]
\Cref{ex:ch26-chain-energy} computes the measure and the energy of three
given potentials.  Here we go in the opposite direction: we prescribe
\(\mu\) and solve \eqref{eq:curve-poisson}.

Write \(u\) for the piecewise affine solution with vertex values
\((u_0,u_1,u_2,u_3)\) and edge slopes \(s_i=u_i-u_{i-1}\).  By
\eqref{eq:ch31-graph-laplacian},
\begin{equation}
\label{eq:ch31-chain-laplacian}
 \Delta_\Gamma u
 =s_1\delta_0+(s_2-s_1)\delta_1+(s_3-s_2)\delta_2-s_3\delta_3 ,
\end{equation}
so \eqref{eq:curve-poisson} is a triangular linear system: the slopes are
determined by \(\mu\) through
\(s_1=\mu(\{0\})-1\), \(s_2=s_1+\mu(\{1\})\), \(s_3=s_2+\mu(\{2\})\), and
the last equation \(-s_3=\mu(\{3\})-1\) is automatic because
\(\mu(\Gamma)=2\).  The following table records four measures, the solution
normalized by \(\max_iu_i=0\), and the energy computed from
\eqref{eq:ch26-chain-energy}.
\begin{center}
\small
\begin{tabular}{l|c|c|c}
\toprule
\(\mu\)&\((s_1,s_2,s_3)\)&\((u_0,u_1,u_2,u_3)\)&\(E_\omega(u)\)\\
\midrule
\(\delta_0+\delta_3\) (the reference)&\((0,0,0)\)&\((0,0,0,0)\)&\(0\)\\
\(\delta_1+\delta_2\) (the tent)&\((-1,0,1)\)&\((0,-1,-1,0)\)&\(-1\)\\
\(\tfrac12(\delta_0+\delta_1+\delta_2+\delta_3)\)
 &\((-\tfrac12,0,\tfrac12)\)&\((0,-\tfrac12,-\tfrac12,0)\)
 &\(-\tfrac14\)\\
\(2\delta_1\)&\((-1,1,1)\)&\((-1,-2,-1,0)\)&\(-\tfrac52\)\\
\(2\delta_0\)&\((1,1,1)\)&\((-3,-2,-1,0)\)&\(-\tfrac92\)\\
\bottomrule
\end{tabular}
\end{center}
Each row is checked three ways.  The slopes reproduce \(\mu\) through
\eqref{eq:ch31-chain-laplacian}; the four masses are the four left-hand
sides of \eqref{eq:ch20-chain-constraints}, hence nonnegative, so each
\(u\) is \(\omega\)-psh; and the total mass is \(2\) in every row.  For the
energies, \eqref{eq:ch26-chain-energy} gives, for instance in the last row,
\(u_0+u_3-\tfrac12(1+1+1)=-3-\tfrac32=-\tfrac92\).  The second row agrees
with \Cref{ex:ch26-chain-energy}, which is the only overlap between the two
tables; that agreement is a normalization test, not a coincidence.

Two features are worth noticing.  First, the energy decreases as the measure
is concentrated: moving from the balanced measure
\(\tfrac12\sum_i\delta_i\) to \(2\delta_0\) drops \(E_\omega\) from
\(-\tfrac14\) to \(-\tfrac92\).  This is the concavity of \(E_\omega\)
together with the fact that the reference potential \(0\) maximizes it among
potentials normalized by \(\sup u=0\), which is
\Cref{cor:energy-monotone-concave} in dimension one.  Second, the solution
for \(2\delta_0\) is the potential whose measure is a single atom at the
\emph{end} of the chain, and it is the extreme point of the admissible
polytope in the direction \(-e_0\); this is the finite-dimensional shadow of
the envelope description of \Cref{sec:ch31-divisorial}.
\end{example}

\begin{example}[A diffuse measure on the same chain]
\label{ex:ch31-chain-lebesgue}
Keep the data of \Cref{ex:ch31-chain-solve} and take for \(\mu\) the
normalized length measure
\[
 \mu=\tfrac23\,\rd s\quad\text{on }\Gamma=[0,3],
 \qquad \mu(\Gamma)=\tfrac23\cdot3=2=V .
\]
Now \eqref{eq:curve-poisson} reads \(u''=\tfrac23\) on \((0,3)\) together
with the endpoint conditions \(1+u'(0^+)=0\) and \(1-u'(3^-)=0\) coming from
\eqref{eq:ch31-segment-general-ma}.  Hence \(u'(t)=-1+\tfrac23t\) and,
normalizing \(u(0)=0\),
\begin{equation}
\label{eq:ch31-chain-lebesgue-potential}
 u(t)=\tfrac13t^2-t,
 \qquad u(0)=u(3)=0,
 \qquad \min_\Gamma u=u(\tfrac32)=-\tfrac34 .
\end{equation}
The endpoint conditions are consistent: \(u'(3^-)=-1+2=1\) as required.  The
energy is
\[
 E_\omega(u)=\int_\Gamma u\,\rd\mu_0-\tfrac12\int_0^3|u'|^2\rd s
 =\bigl(u(0)+u(3)\bigr)-\tfrac12\int_0^3\Bigl(\tfrac23t-1\Bigr)^2\rd t
 =0-\tfrac12\cdot1=-\tfrac12 ,
\]
the integral being \(\tfrac32\int_{-1}^{1}w^2\,\rd w=1\) after the
substitution \(w=\tfrac23t-1\).

This potential is continuous and \(\omega\)-psh but is \emph{not} a model
function: it is strictly convex on \((0,3)\), whereas a model function is
piecewise affine.  It is therefore an element of the second row of the table
in \Cref{sec:ch24-four-products} and not of the first, and it is the
smallest example in this chapter of a potential to which
\Cref{thm:ch19-cld-bfj-comparison}\,(B) applies but not (A).  A model
approximation is obtained by piecewise affine interpolation at
\(t=0,\tfrac1N,\ldots,3\), whose measure is atomic of mass \(2\) and
converges weakly to \(\tfrac23\rd s\); the uniform error is
\(\tfrac13\cdot\tfrac{(3/N)^2}{4}\cdot\tfrac{1}{1}=\tfrac{3}{4N^2}\), since
the interpolation error of a function with \(u''=\tfrac23\) on an interval
of length \(h\) is \(\tfrac{1}{8}u''h^2\).
\end{example}

\subsection{The circle: the Tate curve in the BFJ language}
\label{sec:ch31-circle}

\begin{example}[The canonical metric on a Tate curve, computed twice]
\label{ex:ch31-tate-bfj}
\index{Tate curve!BFJ Monge--Ampere}
Let \(q\in K^\times\) with \(\ell=v(q)>0\), let
\(E_q=\mathbf G_{m,K}^{\An}/q^{\mathbf Z}\) be the corresponding Tate curve
and \(L=\cO_{E_q}(d\cdot(O))\), of degree \(d>0\).  The minimal regular
semistable model has special fibre a cycle of \(\ell\) rational curves when
\(\ell\in\mathbf Z_{>0}\), and in all cases the skeleton is the circle
\[
 \Gamma=\mathbf R/\ell\mathbf Z ,
\]
of circumference \(\ell\) (\Cref{ex:ch08-tate-curve},
\Cref{ex:ch19-tate-canonical}).  Choose the model in which all the degree of
\(\cL\) sits on one component, so that
\[
 \mu_0=d\,\delta_{t=0},\qquad V=\deg L=d .
\]

\emph{Solving for the Haar measure.}  Prescribe
\(\mu=\tfrac d\ell\,\rd t\), of mass \(d=V\).  Equation
\eqref{eq:curve-poisson} becomes \(u''=\tfrac d\ell\) on \((0,\ell)\)
together with the single condition that \(u\) descend to the circle, that
is, \(u(0)=u(\ell)\).  Thus \(u'(t)=c+\tfrac d\ell t\) and
\(\int_0^\ell u'=c\ell+\tfrac{d\ell}{2}=0\), whence \(c=-\tfrac d2\) and,
normalizing \(u(0)=0\),
\begin{equation}
\label{eq:ch31-tate-potential}
 u(t)=\frac{d}{2\ell}\,t\,(t-\ell),
 \qquad 0\leq t\leq\ell,
 \qquad \min_\Gamma u=u(\tfrac\ell2)=-\frac{d\ell}{8} .
\end{equation}
The slope jump at \(t=0\) is
\(u'(0^+)-u'(\ell^-)=-\tfrac d2-\tfrac d2=-d\), which cancels the atom
\(d\delta_0\) of \(\mu_0\) exactly, as it must.

\emph{Comparison with \Cref{part:cld}.}  Formula
\eqref{eq:ch31-tate-potential} is verbatim \eqref{eq:ch19-tate-weight}, the
weight of the \emph{canonical} metric of \(L\) computed there from the
functional equation \eqref{eq:ch19-tate-cocycle} of the theta cocycle, and
the measure \(\tfrac d\ell\rd t\) is \eqref{eq:ch19-tate-measure}.  Two
constructions with nothing in common---a Poisson equation on a circle, and
the unique solution of a cocycle relation on \(\mathbf G_m\)---produce the
same function.  That they must is
\Cref{thm:ch19-cld-bfj-comparison}\,(B) together with
\Cref{cor:ch19-transfer}\,\cref{itm:ch19-tr:abelian}; that they do is a
check on the sign of \(\ddc\), the normalization \(v(\pi)=1\), and the edge
lengths.

\emph{The energy in closed form.}  By \Cref{prop:ch26-graph-energy} and
\(u(0)=0\),
\begin{equation}
\label{eq:ch31-tate-energy}
 E_\omega(u)=\int_\Gamma u\,\rd\mu_0-\frac12\int_0^\ell|u'|^2\rd t
 =0-\frac12\cdot\frac{\ell d^2}{12}
 =-\frac{\ell d^2}{24},
\end{equation}
the integral being
\(\int_0^\ell\bigl(\tfrac d\ell t-\tfrac d2\bigr)^2\rd t
 =\tfrac\ell d\int_{-d/2}^{d/2}w^2\rd w=\tfrac{\ell d^2}{12}\).
For \(d=\ell=3\) this is \(-\tfrac98\).

\emph{The model approximations, exactly.}  Let \(u_N\) be the piecewise
affine interpolation of \(u\) at the \(N\) points \(t_i=i\ell/N\).  Since
\(u\) is a quadratic polynomial, the slope of \(u_N\) on
\([t_i,t_{i+1}]\) is \(u'\) evaluated at the midpoint, namely
\(a_i=\tfrac dN\bigl(i+\tfrac12-\tfrac N2\bigr)\), and
\[
 \int_\Gamma|u_N'|^2\rd t
 =\frac\ell N\sum_{i=0}^{N-1}a_i^2
 =\frac{\ell d^2}{N^3}\cdot\frac{N^3-N}{12}
 =\frac{\ell d^2}{12}\Bigl(1-\frac1{N^2}\Bigr),
\]
using \(\sum_{i=0}^{N-1}(i+\tfrac12-\tfrac N2)^2=\tfrac{N^3-N}{12}\).
Hence
\begin{equation}
\label{eq:ch31-tate-energy-approx}
 E_\omega(u_N)=-\frac{\ell d^2}{24}\Bigl(1-\frac1{N^2}\Bigr)
 \;\searrow\;E_\omega(u)=-\frac{\ell d^2}{24},
\end{equation}
with \(u_N\geq u\) and \(\|u_N-u\|_\infty=\tfrac{d\ell}{8N^2}\).  Both the
monotonicity \(E_\omega(u_N)\geq E_\omega(u)\) and the Lipschitz bound
\(|E_\omega(u_N)-E_\omega(u)|\leq V\|u_N-u\|_\infty=\tfrac{d^2\ell}{8N^2}\)
of \Cref{prop:ch26-energy-lipschitz} are visible in
\eqref{eq:ch31-tate-energy-approx}, the true error \(\tfrac{\ell d^2}{24N^2}\)
being three times smaller than the Lipschitz bound permits.  The measures
\(\MA_\omega(u_N)=\tfrac dN\sum_i\delta_{t_i}\) are
\eqref{eq:ch19-tate-model-approx}.  The energy values in
\eqref{eq:ch31-tate-energy-approx} decrease for all integers \(N\), but the
interpolants themselves are ordered only when the subdivisions are nested.
In particular convexity gives \(u_{2^{j+1}}\leq u_{2^j}\), so the dyadic
subsequence decreases uniformly to \(u\).
\end{example}

\begin{remark}[Where the Green kernel of the circle appears]
\label{rem:ch31-circle-green}
\eqref{eq:ch31-tate-potential} can also be produced from the kernel of
\Cref{lem:ch31-green-kernel} by integrating the explicit circle Green
function \eqref{eq:ch08-green-circle} against \(\mu-\mu_0\).  Doing so is a
useful bookkeeping check and gives the same answer; the differential
route above is shorter because \(\mu-\mu_0\) has a density on the
complement of one point.  The Green function of the circle is the
non-Archimedean N\'eron function of a degree-zero divisor, and
\eqref{eq:ch31-tate-potential} is its second primitive, as recorded after
\eqref{eq:ch19-tate-weight}.
\end{remark}

\subsection{An atom off every toric skeleton}
\label{sec:ch31-p1-blowup}

The examples so far live on a skeleton.  The next one does not, and it will
be used twice: in \Cref{sec:ch31-toric} to show that the invariance
hypothesis cannot be dropped, and in \Cref{sec:ch31-divisorial} as the
smallest instance of an envelope with prescribed divisorial support.

\begin{example}[A semipositive model metric on \(\mathbf P^1\) whose measure
sits off the toric skeleton]
\label{ex:ch31-p1-blowup}
Let \(X=\mathbf P^1_K\), \(L=\cO(1)\), so \(V=\deg L=1\).  Let
\(\cX_0=\mathbf P^1_R\) with \(\cL_0=\cO(1)\), whose model metric is the
standard one of \eqref{eq:p1-standard-weight}; its curvature is
\begin{equation}
\label{eq:ch31-p1-reference}
 \mu_0=\MA_\omega(0)=\delta_{x_{\mathrm G}},
\end{equation}
a unit atom at the Gauss point, which is \eqref{eq:p1-standard-measure}.

Let \(\xi\in\mathbf P^1_k\) be the reduction of the \(K\)-point \(z=1\), and
let \(\rho:\cY\to\cX_0\) be the blowup at \(\xi\), with exceptional divisor
\(E\).  Then \(\cY\) is regular, \(\cY_s=\widetilde{\mathbf P^1_k}+E\) is
reduced with two components meeting transversely at one point, so \(\cY\) is
semistable, and \(b_E=1\).  Its skeleton is a segment of length one,
\[
 \Delta_{\cY}=[x_{\mathrm G},x_E],
\]
where \(x_E\) is the divisorial point attached to \(E\); concretely
\(x_E\) is the Gauss point of the closed disc
\(D(1,|\pi|)\), a point of type \(2\) at tree distance \(1\) from
\(x_{\mathrm G}\) in the direction of \(\xi\)
(\Cref{sec:ch08-p1-tree}, \eqref{eq:gauss-disc}).

\emph{The two model functions.}  Let \(\varphi_E\) be the model function
determined on \(\cY\) by the vertical divisor \(E\); by
\eqref{eq:appJ-model-divisor} it is the pullback under the retraction
\(p_{\cY}\) of the affine function on \([x_{\mathrm G},x_E]\) with
\[
 \varphi_E(x_{\mathrm G})=b_{\widetilde{\mathbf P^1}}^{-1}
   \ord_{\widetilde{\mathbf P^1}}(E)=0,
 \qquad
 \varphi_E(x_E)=b_E^{-1}\ord_E(E)=1 .
\]
Thus \(\varphi_E\) rises with slope \(1\) from \(0\) at \(x_{\mathrm G}\) to
\(1\) at \(x_E\) and is constant elsewhere.  Set \(\psi:=-\varphi_E\).

\emph{Their Laplacians and measures.}  On the segment
\([x_{\mathrm G},x_E]\), \eqref{eq:ch31-graph-laplacian} gives
\(\ddc\varphi_E=\delta_{x_{\mathrm G}}-\delta_{x_E}\), so
\begin{equation}
\label{eq:ch31-p1-two-measures}
 \MA_\omega(\varphi_E)=2\delta_{x_{\mathrm G}}-\delta_{x_E},
 \qquad
 \MA_\omega(\psi)=\delta_{x_E}.
\end{equation}
Hence \(\varphi_E\) is \emph{not} \(\omega\)-psh, while \(\psi\) is, and its
measure is a single unit atom at \(x_E\).

\emph{The same two facts read off the model.}  Let
\(\cL_\psi=\rho^*\cL_0\otimes\cO_{\cY}(-E)\).  Since \(E\simeq\mathbf P^1_k\)
with \((E^2)=-1\) and \((E\cdot\widetilde{\mathbf P^1})=1\),
\[
 \deg\bigl(\cL_\psi|_E\bigr)=-(E^2)=1,
 \qquad
 \deg\bigl(\cL_\psi|_{\widetilde{\mathbf P^1}}\bigr)
  =1-(E\cdot\widetilde{\mathbf P^1})=0 .
\]
Both are nonnegative, so \(\cL_\psi\) is relatively nef and \(\psi\) is
\(\omega\)-psh by \Cref{def:ch20-semipositive}; and the model formula
\eqref{eq:model-mixed-measure} returns
\(1\cdot1\cdot\delta_{x_E}+1\cdot0\cdot\delta_{x_{\mathrm G}}=\delta_{x_E}\),
confirming \eqref{eq:ch31-p1-two-measures}.  For \(\cL_{\varphi_E}
=\rho^*\cL_0\otimes\cO_{\cY}(E)\) the two degrees are \(-1\) and \(2\), the
first negative; this is the algebraic form of the failure of positivity.
The total mass is \(1\) in both cases, as \eqref{eq:mixed-total-mass}
requires even for a signed measure.

\emph{Where \(x_E\) is.}  The toric skeleton of \(\mathbf P^1\) for the
standard torus \(\mathbf G_m=\Spec K[z^{\pm1}]\) is the path from \(0\) to
\(\infty\) through \(x_{\mathrm G}\), namely the set of monomial points
\(\eta_t\) with \(-\log|z|=t\).  Every point of the disc \(D(1,|\pi|)\)
satisfies \(|z|=1\), so
\begin{equation}
\label{eq:ch31-p1-trop-collapse}
 \trop(x_E)=0=\trop(x_{\mathrm G}),
 \qquad x_E\notin\Sk(\mathbf G_m),
\end{equation}
and \(x_E\) lies in the open tube over \(\xi\), which meets the toric
skeleton not at all.  Moreover, since \(\varphi_E\) and \(\psi\) factor
through \(p_{\cY}\) and the whole toric skeleton retracts to
\(x_{\mathrm G}\),
\begin{equation}
\label{eq:ch31-p1-restrictions-agree}
 \varphi_E\big|_{\Sk(\mathbf G_m)}
 =\psi\big|_{\Sk(\mathbf G_m)}=0 .
\end{equation}
Two potentials with the same restriction to the toric skeleton, one
\(\omega\)-psh and one not, with measures \(\delta_{x_E}\) and
\(2\delta_{x_{\mathrm G}}-\delta_{x_E}\).  This is the content of the
example, and it is used in \Cref{ex:ch31-noninvariant}.

\emph{Energy.}  By \Cref{prop:ch26-graph-energy} applied on
\(\Delta_{\cY}\), with \(\mu_0=\delta_{x_{\mathrm G}}\) and
\(\psi(x_{\mathrm G})=0\), \(\psi(x_E)=-1\),
\[
 E_\omega(\psi)=\int\psi\,\rd\mu_0-\tfrac12\int|\psi'|^2\rd s
 =0-\tfrac12=-\tfrac12 .
\]
The reader may check \eqref{eq:two-component-energy} with
\(m_0=1,m_1=0\) gives the same number: the present example is the
degenerate case of \Cref{ex:two-component-complete} in which one of the two
components carries no degree.
\end{example}

\begin{remark}[What changes off the skeleton]
\label{rem:ch31-retraction}
Let the reference form \(\omega\) be determined on the semistable
model defining \(p:X^{\An}\to\Gamma\).
For an \(\omega\)-psh rational model function \(u\),
\cite[Proposition~3.8]{GJKM19} gives
\(u\leq u\circ p\), with \(u\circ p\) again a rational model
\(\omega\)-psh function.  The psh hypothesis is essential for this
inequality; it is not asserted for arbitrary model functions.
The weighted refinement identity
\cite[Proposition~2.2.26(iv)]{Thu05}, computed on a skeleton that
actually determines \(u\) and the reference (after the required
scalar extension when necessary), gives
\[
 p_*\bigl(\mathrm d\mathrm d^c u\bigr)
              =\Delta_\Gamma(u|_\Gamma).
\]
Scalar projection uses ordinary measure pushforward, as in
\eqref{eq:ch29-curve-scalar-weights}.
Thus pushing the two measures of
\eqref{eq:ch31-p1-two-measures} to the toric skeleton gives
\(\delta_{x_{\mathrm G}}\) in both cases, consistent with
\eqref{eq:ch31-p1-restrictions-agree}.
The information lost is the position of the mass, not its amount.
The analytic solution of \cref{prop:graph-ma-solution} follows
from its current and approximation proof, not from a presumed
factorization of every model function through \(\Gamma\).

The corresponding retraction statement is false in higher dimension.  No
proof later in the book treats a higher-dimensional skeleton as though it
were automatically a deformation-invariant psh quotient, and the reader
should treat any such argument with suspicion; the correct higher-dimensional
statement is the equicontinuity of \Cref{thm:bfj-skeletal-equicontinuity}
together with the skeletal characterization
\Cref{thm:bfj-skeletal-characterization}, both of which are statements about
a \emph{family} of skeleta and not about a single one.
\end{remark}

\section{Toric varieties and Alexandrov measures}
\label{sec:ch31-toric}

\subsection{Coordinates, conventions, and the recession function}
\label{sec:ch31-toric-setup}

Let \(X=X_\Sigma\) be a projective toric variety of dimension \(n\) over
\(K\) with dense torus \(T=\Spec K[M]\), cocharacter lattice
\(N=\operatorname{Hom}(M,\mathbf Z)\), and let \(L=L_P\) be an ample toric
line bundle with lattice polytope \(P\subset M_\mathbf R\).  We use the
convention of \Cref{chap:cld-examples}, recorded in
\eqref{eq:toric-trop}:
\begin{equation}
\label{eq:ch31-trop-convention}
 \operatorname{trop}(x)(m)=-\log|\chi^m(x)|=v_x(\chi^m),
 \qquad
 \trop:X_\Sigma^{\An}\to N_\Sigma ,
\end{equation}
and metric weights are convex functions \(F\) on \(N_\mathbf R\) with the
prescribed recession function
\begin{equation}
\label{eq:ch31-recession}
 \check h_P(u)=\sup_{m\in P}\langle-m,u\rangle=h_{-P}(u),
\end{equation}
that is, \(F-\check h_P\) is bounded.  Under that convention the subgradient
image \(\overline{\bigcup_u\partial F(u)}\) is \(-P\)
(\Cref{lem:ch19-toric-recession}); switching to concave roof functions on
\(P\) changes the signs but not the measure, as recorded in
\Cref{def:ch26-roof}.  We write \(\iota:N_\mathbf R\to T^{\An}\) for the
Gauss section \eqref{eq:toric-gauss-section} and \(\Sk(T)=\iota(N_\mathbf R)\)
for the toric skeleton, so that \(\trop\circ\iota=\operatorname{id}\)
(\Cref{lem:ch19-toric-chart}).

Three normalizations are in force and are the source of every sign error in
this subject.

\begin{enumerate}
\item\label{itm:ch31-toric:trop}
 \emph{Valuation, not logarithm.}  \eqref{eq:ch31-trop-convention} uses
 \(v=-\log|\cdot|\), so a point close to the toric boundary where
 \(|\chi^m|\) is small has \emph{large} \(\trop\).  This is why the
 recession function \eqref{eq:ch31-recession} is \(h_{-P}\) and not
 \(h_P\).
\item\label{itm:ch31-toric:lattice}
 \emph{Lattice Lebesgue measure.}  \(\vol_M\) normalizes the covolume of
 \(M\) to one; \(\MA_{\mathbf R}\) is defined by
 \eqref{eq:appD-alexandrov-measure} using \(\operatorname{Leb}_M\).  A
 sublattice of index \(r\) multiplies every mass by \(r\).
\item\label{itm:ch31-toric:factorial}
 \emph{The factor \(n!\).}  The non-Archimedean measure is \(n!\) times the
 Alexandrov measure, \eqref{eq:appJ-alexandrov}.  For \(n=1\) it is
 invisible; for \(n=2\) it is the factor \(2\) that makes
 \(2\vol_M(P)=(L^2)\).
\end{enumerate}

\subsection{The comparison theorem}
\label{sec:ch31-toric-theorem}

\begin{theorem}
\label{thm:toric-bfj-alexandrov}
For an invariant semipositive metric with convex potential \(F\),
\begin{equation}
        \operatorname{trop}_*\MA_L(F)=n!\,\MA_{\mathbf R}(F),
\label{eq:toric-bfj-alexandrov}
\end{equation}
with lattice Lebesgue measure on \(M_\mathbf R\).  Hence
\begin{equation}
        \MA_L(F)(X^{\An})=n!\vol_M(P)=(L^n).
\label{eq:toric-total-mass}
\end{equation}
For polyhedral \(F\), the mass at \(u\in N_\mathbf R\) is
\begin{equation}
                  n!\vol_M(\partial F(u)),
\label{eq:polyhedral-alexandrov-atom}
\end{equation}
where \(\partial F(u)\) is the subgradient polytope.
\end{theorem}

\begin{proof}[Model-to-continuous proof]
\emph{Step 1: rational polyhedral \(F\) and its toric model.}  Suppose
first that \(F\) is a convex piecewise affine function with slopes in
\(-P\cap M\) and rational breaks, so that after replacing \(K\) by a finite
extension if necessary all its values lie in \(v(K^\times)\).  Write
\(F=\max_{1\leq j\leq r}\{\langle -m_j,u\rangle+c_j\}\) with
\(m_j\in P\cap M\) and \(c_j\in v(K^\times)\).  The domains of affinity of
\(F\) form a polyhedral complex \(\Pi_F\) subdividing \(N_\mathbf R\), and
by the toric dictionary of \Cref{sec:toric-convex} the pair
\((\Pi_F,\{c_j\})\) determines a toric model \((\cX_F,\cL_F)\) of
\((X,L)\) over \(R\) whose special-fibre components correspond bijectively
to the \emph{vertices} of \(\Pi_F\), the divisorial point attached to the
component over a vertex \(u\) being \(\iota(u)\).  The associated model
metric has weight \(F\); this is \Cref{rem:ch08-toric-model}.

\emph{Step 2: the intersection number at a vertex.}  Let \(u\) be a vertex
of \(\Pi_F\) and \(E_u\) the corresponding component of \((\cX_F)_s\), with
multiplicity \(b_{E_u}\) equal to the denominator of \(u\), so that
\(b_{E_u}=1\) when \(u\) is a lattice point.  Then \(E_u\) is the
projective toric variety over \(k\) attached to the normal fan of the
subgradient polytope \(\partial F(u)\subseteq-P\), and the restriction
\(\cL_F|_{E_u}\) is the associated toric line bundle.  By the toric degree
formula \eqref{eq:appJ-toric-degree} the weighted degree is
\begin{equation}
\label{eq:ch31-toric-vertex-degree}
 b_{E_u}\bigl(\cL_F|_{E_u}\bigr)^n=n!\vol_M\bigl(\partial F(u)\bigr),
\end{equation}
the multiplicity compensating exactly for the lattice index introduced by a
non-integral vertex; for a lattice vertex both factors are read off
directly.  Substituting into the model formula
\eqref{eq:model-mixed-measure} gives
\begin{equation}
\label{eq:ch31-toric-model-measure}
 \MA_L(F)=\sum_{u\ \text{vertex of}\ \Pi_F}
      n!\vol_M\bigl(\partial F(u)\bigr)\,\delta_{\iota(u)},
\end{equation}
which is \eqref{eq:polyhedral-alexandrov-atom}, including the factor
\(n!\); and since \(\MA_{\mathbf R}(F)
 =\sum_u\vol_M(\partial F(u))\delta_u\) for a polyhedral \(F\)
(\Cref{lem:ch08-alexandrov-polyhedral}, \eqref{eq:appD-polyhedral-atom}),
formula \eqref{eq:ch31-toric-model-measure} is
\eqref{eq:toric-bfj-alexandrov} in the polyhedral case, the pushforward
\(\trop_*\) being trivial because \(\trop(\iota(u))=u\).

\emph{Step 3: total mass in the polyhedral case.}  The subgradient
polytopes \(\partial F(u)\), for \(u\) ranging over the vertices of
\(\Pi_F\), cover \(-P\) and have pairwise disjoint interiors
(\Cref{lem:ch08-alexandrov-polyhedral}\,(2)), so
\(\sum_u\vol_M(\partial F(u))=\vol_M(P)\), and the total mass of
\eqref{eq:ch31-toric-model-measure} is \(n!\vol_M(P)\).  On the other hand
that mass is \((L^n)\) by \eqref{eq:mixed-total-mass}.  This proves
\eqref{eq:toric-total-mass} for polyhedral \(F\), and in particular
reproves the toric degree formula.

\emph{Step 4: approximation.}  Let now \(F\) be an arbitrary convex function
with \(F-\check h_P\) bounded, so that \(\|\cdot\|_F\) is a continuous
invariant semipositive metric.  Rational polyhedral convex functions with
the same recession function are uniformly dense among such \(F\): one takes
for \(F_\varepsilon\) the supremum of finitely many affine minorants of
\(F\) with slopes in \(-P\cap\tfrac1kM\) and rational constants, chosen so
that the supremum is within \(\varepsilon\) of \(F\) on a large ball and has
the same recession function \(\check h_P\), which by
\Cref{lem:ch19-toric-recession} then forces
\(\|F_\varepsilon-F\|_\infty\leq\varepsilon\).  This is Step~3 of the proof
of \Cref{thm:ch19-toric-ma}, to which we refer for the details.  On the
non-Archimedean side,
uniform convergence of weights gives weak convergence of the measures by
\Cref{prop:ch24-continuous-uniform}.  On the convex side, Alexandrov
measures are weakly continuous under locally uniform convergence of finite
convex functions (\Cref{thm:appD-alexandrov-measure}).  Passing to the limit
in Step~2 proves \eqref{eq:toric-bfj-alexandrov} in general, and
\eqref{eq:toric-total-mass} follows since both sides of
\eqref{eq:toric-bfj-alexandrov} have the constant total mass computed in
Step~3.

\emph{Step 5: consistency with \Cref{part:cld}.}  The measure produced in
Steps~1--4 is the BFJ measure of \Cref{def:mixed-model-measure} and its
uniform extension.  That it coincides with the Chambert--Loir measure and
with the CLD Bedford--Taylor product is
\Cref{thm:ch19-cld-bfj-comparison}\,(A) and (B), and the resulting identity
is \Cref{thm:ch19-toric-ma}, proved there by the tropical route
\eqref{eq:cld-toric-ma-comparison} of \cite[Theorem~6.1.1]{BGJKK25}.  The
two proofs are genuinely different: the present one computes an
intersection number on a toric special fibre, whereas the one in
\Cref{chap:cld-examples} computes a superform product on a tropical chart.
Local facewise comparison is also treated in \cite{Vil20}.
\end{proof}

\begin{remark}[Which hypotheses are used where]
\label{rem:ch31-toric-hypotheses}
Step~1 uses the rationality of \(F\) and the existence of the toric model,
hence a discretely valued \(K\) with \(v(\pi)=1\) after a finite extension.
Step~2 uses only toric intersection theory over the residue field.  Step~4
uses neither discreteness nor residue characteristic, so
\eqref{eq:toric-bfj-alexandrov} in the form
\(\trop_*c_1(L,\|\cdot\|_F)^n=n!\MA_{\mathbf R}(F)\) is available over any
complete nontrivially valued field, which is how it is stated in
\Cref{thm:ch19-toric-ma}.  What requires \eqref{eq:ch19-bfj-hypotheses} is
the identification of the left-hand side with the \emph{BFJ}
Monge--Amp\`ere measure, since the BFJ theory has no models over a
non-discrete valuation ring; this is the first item of
\Cref{rem:ch19-cld-bfj-hypotheses}.
\end{remark}

\subsection{Polyhedral potentials: atoms and lattice volumes}
\label{sec:ch31-toric-polyhedral}

We begin with the smallest possible checks and work upwards in dimension.

\begin{example}[Two atoms on \(\mathbf P^1\)]
\label{ex:ch31-toric-p1-two-atoms}
Let \(X=\mathbf P^1\), \(N=M=\mathbf Z\), \(L=\cO(2)\), so
\(P=[0,2]\), \(\vol_M(P)=2\) and \((L)=1!\cdot2=2=\deg\cO(2)\).  The
recession function is \(\check h_P(u)=\max\{0,-2u\}\).  Take
\[
 F(u)=\max\{0,\;-u+1,\;-2u+1\}.
\]
The three slopes \(0,-1,-2\) lie in \(-P=[-2,0]\), and
\(0\leq F-\check h_P\leq1\), so \(F\) is a model metric weight for \(L\).
Its breaks are at \(u=0\), where the pieces of slope \(-1\) and \(-2\)
meet, and at \(u=1\), where the pieces of slope \(0\) and \(-1\) meet.
Hence
\[
 \partial F(0)=[-2,-1],\qquad \partial F(1)=[-1,0],
\]
of lattice length \(1\) each, and by
\eqref{eq:polyhedral-alexandrov-atom} with \(n!=1\),
\[
 \MA_L(F)=\delta_{\iota(0)}+\delta_{\iota(1)},
 \qquad\text{mass }2=\deg\cO(2).
\]
Here \(\iota(0)=x_{\mathrm G}\) is the Gauss point and \(\iota(1)\) is the
monomial point with \(|z|=|\pi|\).  The two subgradient intervals tile
\(-P=[-2,0]\), which is Step~3 of the proof above in the smallest case.
As a curve computation this is \Cref{lem:ch31-curve-ma-affine} applied to
the toric model whose special fibre is a chain of two rational curves of
degrees \(1,1\); it is the two-component curve of
\Cref{ex:two-component-complete} with \(m_0=m_1=1\), and one checks
\(E_\omega(F)\) against \eqref{eq:two-component-energy} after normalizing
\(F\) to vanish at one of the two points.
\end{example}

\begin{example}[A polyhedral atom on \(\mathbf P^n\)]
\label{ex:ch31-toric-atom-pn}
Let \(P=d\Delta_n\), so \(L=\cO_{\mathbf P^n}(d)\).  Choose a
polyhedral potential whose subgradient at the origin is all of \(-P\) and
whose subgradient is lower-dimensional elsewhere; the reflected support
function \(F=\check h_P=\max\{0,-du_1,\ldots,-du_n\}\) is such a potential.
Then
\[
             \MA_L(F)=n!\vol(P)\,\delta_{\xi_{\mathrm G}}
                      =d^n\delta_{\xi_{\mathrm G}},
\]
where \(\xi_{\mathrm G}\) is the toric Gauss point.  This is the toric
realization of the divisorial Green envelope in
\eqref{eq:green-envelope-divisorial}.  The atom is not normalized to
probability unless it is divided by \(d^n\).
\end{example}

The verification is worth doing once.  Since
\(\vol_M(d\Delta_n)=d^n/n!\), the mass is
\(n!\cdot d^n/n!=d^n=(\cO(d)^n)\), which is
check~\cref{itm:ch31-check:mass}, and the factor \(n!\) cancels exactly
against the factor \(1/n!\) in the volume of a simplex, which is
check~\cref{itm:ch31-check:factorial}.  For \(n=3\), \(d=1\) the numbers are
\(\vol_M(\Delta_3)=\tfrac16\) and \(3!\cdot\tfrac16=1=(\cO(1)^3)\); the
reader who has never carried the factor \(6\) through a computation should
do so here.  The associated model is \(\mathbf P^n_R\) with the standard
model metric, whose special fibre is the single component
\(\mathbf P^n_k\) of degree \(d^n\); so the intersection-theoretic and the
convex computations agree atom for atom.

\begin{example}[Two atoms on \(\mathbf P^2\), with a tiling check]
\label{ex:ch31-toric-p2-two-atoms}
Let \(X=\mathbf P^2\), \(L=\cO(2)\), \(P=2\Delta_2\), so
\(\vol_M(P)=2\) and \((L^2)=2!\cdot2=4=(2H)^2\).  The recession function is
\(\check h_P(u)=\max\{0,-2u_1,-2u_2\}\).  Take
\[
 F(u)=\max\{\,\underbrace{0}_{f_0},\;
             \underbrace{-2u_1}_{f_1},\;
             \underbrace{-2u_2}_{f_2},\;
             \underbrace{-u_1-u_2+1}_{f_3}\,\},
\]
whose slopes \((0,0)\), \((-2,0)\), \((0,-2)\), \((-1,-1)\) all lie in
\(-P\); the last is the midpoint of an edge of \(-P\) and is a lattice
point, so \(F\) is an integral polyhedral function and
\(0\leq F-\check h_P\leq1\).

\emph{Locating the vertices.}  A point of \(N_\mathbf R\) carries an atom
exactly when at least three of the four affine pieces are simultaneously
active and maximal (\Cref{lem:ch08-alexandrov-polyhedral}\,(3)).  The four
triples give
\[
 \begin{array}{llll}
 f_0=f_1=f_2:&u=(0,0),&f_3=1>0:&\text{not maximal};\\
 f_0=f_1=f_3:&u=(0,1),&f_2=-2<0:&\text{maximal};\\
 f_0=f_2=f_3:&u=(1,0),&f_1=-2<0:&\text{maximal};\\
 f_1=f_2=f_3:&\text{no solution},&&
 \end{array}
\]
the last because \(f_1=f_2\) forces \(u_1=u_2=t\) and then \(f_3=f_1\)
gives \(-2t+1=-2t\).

\emph{The masses.}  At \(u=(0,1)\) the active slopes are
\((0,0),(-2,0),(-1,-1)\), so
\(\partial F(0,1)=\operatorname{conv}\{(0,0),(-2,0),(-1,-1)\}\), of lattice
area \(\tfrac12|{\det}\bigl(\begin{smallmatrix}-2&0\\-1&-1\end{smallmatrix}\bigr)|
 =\tfrac12\cdot2=1\).  Symmetrically
\(\partial F(1,0)=\operatorname{conv}\{(0,0),(0,-2),(-1,-1)\}\), also of
area \(1\).  Hence
\begin{equation}
\label{eq:ch31-p2-two-atoms}
 \MA_{\mathbf R}(F)=\delta_{(0,1)}+\delta_{(1,0)},
 \qquad
 \MA_L(F)=2\delta_{\iota(0,1)}+2\delta_{\iota(1,0)},
\end{equation}
of total mass \(2\) and \(4\) respectively.

\emph{Three checks.}  First, the two subgradient triangles tile the triangle
\(-P\) with vertices \((0,0)\), \((-2,0)\), \((0,-2)\): they share the edge
from \((0,0)\) to \((-1,-1)\), and \(1+1=2=\vol_M(P)\).  Second, the
non-Archimedean mass \(4\) equals \((\cO(2)^2)\).  Third, the associated
toric model has special fibre with two components, of residual
self-intersection \(2\) each; the model is obtained from
\(\mathbf P^2_R\) by the weighted blowup encoded by \(\Pi_F\), and the
degrees \(2,2\) are the atoms of \eqref{eq:ch31-p2-two-atoms}.  Compare
\Cref{ex:ch19-toric-two-atoms}, which performs the same computation on
\(\mathbf P^1\times\mathbf P^1\) with two atoms of mass \(1\); the
structural difference is that there the subgradient polytopes are half-unit
squares and here they are triangles of lattice area \(1\).
\end{example}

\subsection{Smooth potentials: densities in coordinates}
\label{sec:ch31-toric-smooth}

\begin{example}[The logistic potential on \(\mathbf P^1\)]
\label{ex:ch31-toric-logistic}
Let \(X=\mathbf P^1\), \(L=\cO(1)\), \(P=[0,1]\), so
\(\check h_P(u)=\max\{0,-u\}\) and \(V=(L)=1\).  Put
\[
 F(u)=\log\bigl(1+\rme^{-u}\bigr).
\]
Then \(F\) is smooth and strictly convex, with
\[
 F'(u)=-\frac{\rme^{-u}}{1+\rme^{-u}}\in(-1,0),
 \qquad
 F''(u)=\frac{\rme^{-u}}{(1+\rme^{-u})^2}>0,
\]
and \(0\leq F-\check h_P\leq\log2\), so \(\|\cdot\|_F\) is a continuous
invariant semipositive metric on \(\cO(1)\).  Writing \(p=-F'(u)\), one has
\(F''=p(1-p)\) and
\begin{equation}
\label{eq:ch31-logistic-measure}
 \MA_L(F)=\iota_*\bigl(p(1-p)\,\rd u\bigr),
 \qquad
 \int_{\mathbf R}p(1-p)\,\rd u
  =-\int_1^0\rd p=1=\deg\cO(1),
\end{equation}
the substitution being legitimate because \(u\mapsto p\) is a
diffeomorphism of \(\mathbf R\) onto \((0,1)\), with reversed orientation
and \(\rd p=-F''\rd u\).  Here
\(\iota:\mathbf R\hookrightarrow X^{\An}\) is the toric skeleton section;
there is no pullback of measures denoted by \(\trop^*\).
The measure is absolutely continuous on the toric skeleton, with density
maximal \(\tfrac14\) at \(u=0\) and decaying like \(\rme^{-|u|}\).

This is the one-dimensional companion of \Cref{ex:ch19-toric-p2-smooth}, and
it is the smallest example in the book of a Monge--Amp\`ere measure that is
neither atomic nor supported on a proper subset of a skeleton.  Under
\Cref{prop:graph-ma-solution} it is also the solution of a Poisson equation:
on the skeleton of \(\mathbf P^1\), with reference measure
\(\mu_0=\delta_{\iota(0)}\) attached to \(\check h_P\), the potential
\(F-\check h_P\) satisfies \(\Delta(F-\check h_P)=p(1-p)\rd u-\delta_0\), and
one checks the atom: the slope of \(F-\check h_P\) jumps by
\(F'(0^+)-F'(0^-)-1=-1\) across \(u=0\) once the corner of
\(\check h_P\) is subtracted.
\end{example}

\begin{example}[Smooth real data need not give a model metric]
\label{ex:ch31-toric-smooth-not-model}
If \(F\) is strictly convex and smooth on the dense tropical stratum, its
Alexandrov measure has density \(\det(D^2F)\).  The associated invariant
semipositive metric is continuous, but it is not a model metric unless
\(F\) is rational polyhedral.  Thus the equality of BFJ and real measures is
not restricted to atomic model measures.
\end{example}

The logistic potential of \Cref{ex:ch31-toric-logistic} makes the point
concrete: its measure has no atoms, whereas by
\eqref{eq:ch31-toric-model-measure} the measure of any model metric is a
finite sum of Dirac masses.  Hence \(F\) is not a model function, and the
metric \(\|\cdot\|_F\) belongs to the second row and not the first of the
table in \Cref{sec:ch24-four-products}.  The same applies to the family
\(\phi_\varepsilon\) of \Cref{ex:ch19-toric-p2-smooth}, whose densities
\eqref{eq:ch19-p2-density} concentrate as \(\varepsilon\downarrow0\) on the
single atom of \Cref{ex:ch31-toric-atom-pn} with \(d=1\).

\begin{example}[The log-sum-exp potential on \(\mathbf P^n\), mass computed
in coordinates]
\label{ex:ch31-toric-logsumexp}
Let \(X=\mathbf P^n\), \(L=\cO(d)\), \(P=d\Delta_n\).  Put
\[
 F(u)=\log\Bigl(1+\sum_{i=1}^n\rme^{-du_i}\Bigr),
 \qquad
 p_i=\frac{\rme^{-du_i}}{1+\sum_j\rme^{-du_j}},
 \qquad
 p_0=1-\sum_{i=1}^np_i .
\]
Then \(\partial F/\partial u_i=-dp_i\), and
\[
 D^2F=d^2\bigl(\operatorname{diag}(p)-pp^{\mathsf T}\bigr),
 \qquad
 \det D^2F=d^{2n}p_0p_1\cdots p_n ,
\]
the determinant being the classical identity
\(\det(\operatorname{diag}(p)-pp^{\mathsf T})=p_1\cdots p_n(1-\sum_ip_i)\).
The gradient map \(u\mapsto(-dp_1,\ldots,-dp_n)\) is a diffeomorphism onto
the interior of \(-d\Delta_n\), with Jacobian determinant of absolute value
\(\det D^2F\); hence
\begin{equation}
\label{eq:ch31-logsumexp-mass}
 \begin{aligned}
 \int_{N_\mathbf R}\det D^2F\,\rd u
   &=\vol_M(d\Delta_n)=\frac{d^n}{n!},\\
 \MA_L(F)(X^{\An})
   &=n!\cdot\frac{d^n}{n!}=d^n=(\cO(d)^n).
 \end{aligned}
\end{equation}
For \(n=d=1\) this is \eqref{eq:ch31-logistic-measure}; for \(n=2\) and
\(d=1\) it is the member \(\varepsilon=1\) of the family
\eqref{eq:ch19-p2-density}.  The computation is recorded here because it is
the cleanest place in the book where the factor \(n!\) is seen to cancel
against the volume of a simplex, and because it exhibits, for every \(n\),
a continuous semipositive metric whose Monge--Amp\`ere measure is
absolutely continuous on an \(n\)-dimensional skeleton.
\end{example}

\subsection{Mixed measures and mixed volumes}
\label{sec:ch31-toric-mixed}

\begin{proposition}
\label{prop:ch31-toric-mixed}
Let \(L_1,\ldots,L_n\) be nef toric line bundles on \(X_\Sigma\) with
polytopes \(P_1,\ldots,P_n\), and assume that
\(L_1+\cdots+L_n\) is ample.  Give them invariant continuous semipositive
metrics with convex weights \(F_1,\ldots,F_n\).  Then
\begin{equation}
\label{eq:ch31-toric-mixed-measure}
 \trop_*\MA(F_1,\ldots,F_n)=n!\,\MA_{\mathbf R}(F_1,\ldots,F_n),
\end{equation}
where the right-hand side is the polarization
\eqref{eq:appD-mixed-ma-polarization} of the Alexandrov measure.  If all
\(F_i\) are polyhedral, then the mass at a point
\(u\in N_{\mathbf R}\) is
\begin{equation}
\label{eq:ch31-toric-mixed-atom}
 n!\,V\bigl(\partial F_1(u),\ldots,\partial F_n(u)\bigr),
\end{equation}
the mixed volume normalized by \(V(K,\ldots,K)=\vol_M(K)\); and the total
mass is
\begin{equation}
\label{eq:ch31-toric-mixed-mass}
 n!\,V(P_1,\ldots,P_n)=(L_1\cdots L_n).
\end{equation}
\end{proposition}

\begin{proof}
For \(\varepsilon_i>0\), the toric line bundle
\(L_\varepsilon:=\sum_i\varepsilon_iL_i\) is ample: it is a positive
multiple of the ample bundle \(\sum_iL_i\), plus a nef class.  The diagonal
toric identity \eqref{eq:toric-bfj-alexandrov} therefore applies to the
weight \(F_\varepsilon:=\sum_i\varepsilon_iF_i\).  Both sides are
homogeneous polynomials of degree \(n\) in \(\varepsilon\): on the
non-Archimedean side this is the multilinearity of Chambert--Loir mixed
measures, first for model metrics and then by uniform approximation, and on
the real side it is the defining polarization.  Thus their coefficients
agree.  Concretely,
\(\MA_{\mathbf R}(\sum_i\varepsilon_iF_i)\) is a homogeneous polynomial of
degree \(n\) in \(\varepsilon\) whose coefficients are, up to multinomial
factors, the mixed measures; the corresponding statement on the
non-Archimedean side is \(\MA(\sum_i\varepsilon_i\varphi_i)\) expanded by
multilinearity, and the two polynomials agree coefficientwise because they
agree for every \(\varepsilon\) in the positive orthant.

For \eqref{eq:ch31-toric-mixed-atom}, apply the same polarization to
\eqref{eq:polyhedral-alexandrov-atom}: for polyhedral \(F_i\) the
subgradient of \(\sum\varepsilon_iF_i\) at \(u\) is the Minkowski sum
\(\sum\varepsilon_i\partial F_i(u)\) (\Cref{lem:appD-support-calculus}), and
\(\vol_M(\sum\varepsilon_i\partial F_i(u))\) expands by
\Cref{thm:appD-minkowski-polynomial} with coefficients the mixed volumes.
Matching coefficients of \(\varepsilon_1\cdots\varepsilon_n\) gives
\eqref{eq:ch31-toric-mixed-atom}.  The mass statement
\eqref{eq:ch31-toric-mixed-mass} follows by the same expansion applied to
Step~3 of \Cref{thm:toric-bfj-alexandrov}, together with the toric mixed
degree formula, which is \eqref{eq:appJ-toric-degree} polarized; it also
follows from \eqref{eq:ch24-mass}, which gives the mass
\((\{\theta_1\}\cdots\{\theta_n\})\) directly.
\end{proof}

\begin{example}[A mixed measure on \(\mathbf P^1\times\mathbf P^1\), with
numbers]
\label{ex:ch31-toric-mixed-numbers}
Let \(X=\mathbf P^1\times\mathbf P^1\), \(L_1=\cO(1,0)\),
\(L_2=\cO(0,1)\), with polytopes \(P_1=[0,1]\times\{0\}\) and
\(P_2=\{0\}\times[0,1]\).  Neither is ample, but the pair is a legitimate
input to \eqref{eq:ch31-toric-mixed-mass} because
\(L_1+L_2=\cO(1,1)\) is ample, and
\[
 V(P_1,P_2)=\tfrac12\bigl(\vol_M(P_1+P_2)-\vol_M(P_1)-\vol_M(P_2)\bigr)
 =\tfrac12(1-0-0)=\tfrac12 ,
\]
so \(2!\,V(P_1,P_2)=1=(L_1\cdot L_2)=(H_1\cdot H_2)\).  Taking the standard
model metrics, with weights \(F_1(u)=\max\{0,-u_1\}\) and
\(F_2(u)=\max\{0,-u_2\}\), the subgradients at the origin are the two unit
segments \([-1,0]\times\{0\}\) and \(\{0\}\times[-1,0]\), whose mixed volume
is \(\tfrac12\); elsewhere at least one subgradient is a point and the mixed
volume vanishes.  Hence
\[
 \MA(F_1,F_2)=2\cdot\tfrac12\,\delta_{\iota(0,0)}=\delta_{x_{\mathrm G}},
 \qquad\text{mass }1 .
\]
The same number is obtained on the model \(\mathbf P^1_R\times\mathbf P^1_R\)
from \eqref{eq:model-mixed-measure}: the special fibre has one component
\(\mathbf P^1_k\times\mathbf P^1_k\) of multiplicity one, and
\((\cO(1,0)\cdot\cO(0,1))=1\).  For a genuinely nontrivial mixed
computation, in which the two metrics are determined on different models,
see \Cref{ex:ch24-toric-surface-mixed}.
\end{example}

\begin{example}[Nef and not big: the mass vanishes]
\label{ex:ch31-nef-not-big}
Keep \(X=\mathbf P^1\times\mathbf P^1\) and take
\(L=\cO(1,0)=\operatorname{pr}_1^*\cO(1)\),
whose polytope is the segment \(P=[0,1]\times\{0\}\subset M_\mathbf R\).
Then \(L\) is nef but not big:
\[
 \vol_M(P)=0,\qquad (L^2)=2!\cdot0=0 .
\]
Consequently \(V=0\) and every statement of \Cref{part:bfj} that divides by
\(V\), or that asserts the mass of a Monge--Amp\`ere measure to be positive,
fails.  Concretely, for any convex \(F\) with \(F-\check h_P\) bounded, the
subgradient image is contained in the segment \(-P\), which has empty
interior, so \(\MA_{\mathbf R}(F)=0\) by
\eqref{eq:appD-alexandrov-measure} and hence \(\MA_L(F)=0\).  The
Monge--Amp\`ere equation \(\MA_L(F)=\mu\) is unsolvable for every
\(\mu\neq0\), the energy \(E_\omega\) is identically zero on every potential
by \eqref{eq:model-energy}, and its differentiability
(\Cref{thm:energy-envelope-differentiability}) is vacuous rather than
informative.

This is the promised example in which a standing hypothesis cannot be
dropped.  The hypothesis is the ampleness of \(\{\theta\}\) in
\eqref{eq:ch19-bfj-hypotheses}, and its role is not technical: it is what
makes \(V>0\), hence what makes the normalization
\eqref{eq:examples-global-normalization} meaningful, the capacity
\eqref{eq:capacity-definition} nondegenerate, and the variational problem of
\Cref{chap:bfj-finite-energy-measures} nontrivial.  Weakening ampleness to
nefness is therefore not a matter of a more careful proof.  What
\emph{can} be done, and is done in \Cref{part:bj} for the trivially valued
theory, is to replace \((\{\theta\}^n)\) by \(\vol(\{\theta\})\) and restrict
to big classes; the hypothesis that cannot be dropped is bigness, and on
\(\cO(1,0)\) it fails.
\end{example}

\begin{warning}{Nef is not enough, and neither is big alone}
\Cref{ex:ch31-nef-not-big} shows that nef and not big is fatal.  Big and not
nef is a different matter: there the volume is positive but the model
measure of \eqref{eq:model-mixed-measure} need not be positive, and one must
pass to an envelope.  \Cref{ex:ch20-nef-not-ample} separates nef from ample,
and the systematic treatment of big classes is \Cref{part:bj}, not
\Cref{part:bfj}.  No statement of the present chapter is asserted outside
the ample case.
\end{warning}

\subsection{Invariance cannot be dropped}
\label{sec:ch31-toric-invariance}

\Cref{thm:toric-bfj-alexandrov} and its CLD counterparts
\Cref{thm:toric-positivity} and \Cref{thm:ch19-toric-ma} all begin with the
words \emph{invariant metric}.  The next example shows what those words are
doing, using nothing beyond the computation already performed in
\Cref{ex:ch31-p1-blowup}.

\begin{example}[Two non-invariant metrics with the same tropical restriction]
\label{ex:ch31-noninvariant}
Let \(X=\mathbf P^1\), \(L=\cO(1)\), \(P=[0,1]\), and let
\(\varphi_E,\psi=-\varphi_E\) be the two model functions of
\Cref{ex:ch31-p1-blowup}, attached to the blowup of \(\mathbf P^1_R\) at the
reduction \(\xi\) of the non-torus-fixed point \(z=1\).  Both are model
functions on \(\mathbf P^{1,\An}\); neither is torus invariant, since
\(\xi\) is not a torus-fixed point of the special fibre.  By
\eqref{eq:ch31-p1-restrictions-agree} their restrictions to the toric
skeleton coincide:
\[
 \varphi_E|_{\Sk(T)}=\psi|_{\Sk(T)}=0,
\]
so the two metric weights restrict to the \emph{same} convex function
\(\check h_P\) on \(N_\mathbf R\).  Nevertheless, by
\eqref{eq:ch31-p1-two-measures},
\[
 \MA_L(\psi)=\delta_{x_E},
 \qquad
 \MA_L(\varphi_E)=2\delta_{x_{\mathrm G}}-\delta_{x_E},
\]
and \(\psi\) is \(\omega\)-psh while \(\varphi_E\) is not.  Three
consequences.

\emph{(i) The positivity criterion fails.}
\Cref{thm:toric-positivity} asserts that an invariant metric is
semipositive if and only if its weight is convex.  The weight of
\(\|\cdot\|_{\varphi_E}\) restricts to the convex function
\(\check h_P\) on \(N_\mathbf R\), yet the metric is not semipositive.  Thus
``convexity of the restriction to the toric skeleton'' is not a sufficient
criterion for positivity of a non-invariant metric, and the invariance
hypothesis in \Cref{thm:toric-positivity} cannot be dropped.

\emph{(ii) The measure is not determined by the restriction.}  The two
metrics have equal restrictions and different Monge--Amp\`ere measures, one
of which, \(\delta_{x_E}\), is supported entirely off the toric skeleton.
Hence in \eqref{eq:toric-bfj-alexandrov} the pushforward \(\trop_*\) cannot
be removed: it is false that \(\MA_L(F)=n!\iota_*\MA_{\mathbf R}(F)\) for a
non-invariant metric whose restriction is \(F\).

\emph{(iii) What does survive.}  Pushing forward, one finds
\(\trop_*\MA_L(\psi)=\delta_0=\trop_*\MA_L(\varphi_E)\), and both equal
\(1!\MA_{\mathbf R}(\check h_P)\).  This is not an accident: in dimension one
the retraction to a skeleton pushes curvature forward
(\Cref{rem:ch31-retraction}), so the identity
\(\trop_*\MA=n!\MA_{\mathbf R}(\text{restriction})\) survives the loss of
invariance for \(n=1\).  It is the \emph{location} of the mass, and with it
positivity, that is lost.  In dimension \(\geq2\) even the pushforward
statement is not available for non-invariant metrics, because there is no
retraction theorem; this is the content of the second half of
\Cref{rem:ch31-retraction} and of the boundary clause attached to the toric
row of \Cref{sec:ch31-table-main}.
\end{example}

\begin{remark}[Invariance as a hypothesis, not a normalization]
\label{rem:ch31-invariance-is-a-hypothesis}
It is tempting to read \Cref{thm:toric-bfj-alexandrov} as saying that toric
Monge--Amp\`ere theory \emph{is} convex analysis.  What it says is that the
restriction of the theory to the invariant category is convex analysis.  The
invariant category is a closed subspace of the space of all semipositive
metrics, not a retract of it: there is no natural projection from metrics to
invariant metrics compatible with \(\MA\), and
\Cref{ex:ch31-noninvariant} shows that restriction to the skeleton is not
such a projection.  The same warning is item \cref{itm:ch19-toric-inv} of
\Cref{sec:ch19-toric-continuous}, and it is why the toric rows of every
table in this book carry the word \emph{invariant}.
\end{remark}

\subsection{Energy and orthogonality, checked against Part III}
\label{sec:ch31-toric-energy}

The toric family also computes two functionals whose general theory occupies
\Cref{chap:bfj-energy,chap:bfj-differentiability}.

\emph{Energy.}  By \Cref{thm:ch26-toric-energy}, the Monge--Amp\`ere energy
of an invariant potential is, up to the normalization there, an integral of
the roof function of \Cref{def:ch26-roof} over \(P\).  On
\Cref{ex:ch31-toric-atom-pn} the energy of \(\check h_P\) relative to itself
is zero, and on \Cref{ex:ch31-toric-p1-two-atoms} the energy relative to the
standard model metric can be computed either from the toric formula or from
the two-vertex graph formula \eqref{eq:ch26-two-vertex-energy}; the two
agree, which is the check.

\emph{Orthogonality.}  \Cref{thm:ch27-toric-orthogonality} proves the
orthogonality property \Cref{def:orthogonality-property} for invariant
data by Legendre duality, the mechanism being the elementary convex
statement \Cref{thm:ch27-legendre-orthogonality}.  On the toric side the
contact locus of an envelope is the set where the convex envelope of the
obstacle touches it, and orthogonality is the statement that the Alexandrov
measure of the envelope is carried by that set.  The reader who wants to see
this in one variable may take \(X=\mathbf P^1\), \(P=[0,1]\), obstacle
\(g(u)=\min\{\check h_P(u)+1,\,\check h_P(u)+\max(0,1-|u|)\}\), and compute
both sides; the general statement, with the hypotheses that make it true, is
\Cref{cor:ch27-toric-energy}.

\begin{warning}{The toric case proves orthogonality, it does not prove
envelope continuity in general}
\Cref{thm:ch27-toric-orthogonality} is a theorem about invariant data, and
its proof is convex duality.  It does not supply the continuity of envelopes
in \Cref{thm:bfj-envelope-regularity}, whose proof needs multiplier ideals
and residue characteristic zero, nor its positive-characteristic substitute
\Cref{thm:ch23-gjkm-positive-characteristic}.  A reader who has verified
orthogonality toricly has verified one hypothesis of
\Cref{thm:energy-envelope-differentiability} in one family, not the theorem.
\end{warning}

\section{Measures at divisorial points}
\label{sec:ch31-divisorial}

\subsection{The finite-pole envelope}
\label{sec:ch31-finite-pole}

Let \(S=\{x_1,\ldots,x_N\}\) be a finite set of divisorial points contained
in an SNC dual complex.  For \(t=(t_1,\ldots,t_N)\in\mathbf R^N\), define
\begin{equation}
 \Phi_{S,t}:=\sup\{\varphi\in\PSH(X,\omega)\mid
                         \varphi(x_i)\leq t_i\text{ for all }i\}.
\label{eq:finite-pole-envelope}
\end{equation}

The construction is a global obstacle problem with a \emph{finite} obstacle:
instead of prescribing an upper bound everywhere, one prescribes it at
finitely many valuations.  The point of the next proposition is that this
finite amount of data already determines a potential whose Monge--Amp\`ere
measure is atomic and carried by \(S\), and that every atomic measure
carried by \(S\) arises this way.

\begin{proposition}
\label{prop:finite-pole-envelope}
Assume orthogonality.  The upper-semicontinuous regularization of
\(\Phi_{S,t}\) is continuous and \(\omega\)-psh, and
\begin{equation}
                         \Supp\MA_\omega(\Phi_{S,t})\subset S.
\label{eq:finite-pole-support}
\end{equation}
Conversely, if
\(\mu=\sum_i a_i\delta_{x_i}\), with \(a_i\geq0\) and
\(\sum_i a_i=V\), there is a vector \(t\) for which
\(\MA_\omega(\Phi_{S,t})=\mu\).  The resulting solution is unique modulo
constants; no uniqueness assertion is made for the constraint vector
\(t\).
\end{proposition}

\begin{proof}
Choose an SNC model \(\mathcal X\) on which every \(x_i\) is a vertex.
The compactness estimate for normalized psh functions supplies \(M>0\)
such that \(\sup_X\varphi\leq M\) whenever
\(\varphi(x_1)\leq t_1\).  Adding one constant to all the \(t_i\) translates
the envelope by the same constant.  We may therefore normalize so that
\(t_i\leq-1\) for every \(i\), and so that every admissible \(\varphi\) is
at most \(-1\) on \(X^{\An}\).

Let \(f_{\mathcal X}\) be the model function that factors through the
retraction \(p_{\mathcal X}\), is affine on every face of
\(\Delta_{\mathcal X}\), takes the value \(t_i\) at \(x_i\), and is zero
at every other vertex.  Convexity on faces and
\(\varphi\leq\varphi\circ p_{\mathcal X}\) then give the crucial
equivalence
\[
 \varphi(x_i)\leq t_i\ \text{for every }i
 \quad\Longleftrightarrow\quad
 \varphi\leq f_{\mathcal X}.
\]
Thus \(\Phi_{S,t}=P_\omega(f_{\mathcal X})\), so it is continuous and
psh.  Orthogonality puts its Monge--Amp\`ere measure on the contact locus;
the normalization above places that locus in \(\{f_{\mathcal X}<0\}\).
Doing this on every dominating SNC model yields
\[
 \Supp\MA_\omega(\Phi_{S,t})
 \subseteq\bigcap_{\mathcal X}\{f_{\mathcal X}<0\}.
\]
The intersection is exactly \(S\): in terms of the center map
\(c_{\mathcal X}:X^{\An}\to\mathcal X_0\), the set
\(\{f_{\mathcal X}<0\}\) is the union of the tubes over the generic points
\(c_{\mathcal X}(x_i)\), and a point lying in all these tubes is one of the
divisorial valuations \(x_i\).  This proves
\eqref{eq:finite-pole-support}.

For the converse, the supported Calabi--Yau theorem gives a continuous
solution \(\varphi\) for \(\mu\).  Put \(t_i=\varphi(x_i)\).  Then
\(\varphi\leq\Phi_{S,t}\), while the two functions agree on
\(\Supp\MA(\varphi)\subseteq S\).  The domination lemma gives the reverse
inequality and hence \(\Phi_{S,t}=\varphi\).  This is precisely
\cite[Lemma~8.5 and Proposition~8.6]{BFJ15}.
\end{proof}

Each step of that proof quotes a theorem of \Cref{part:bfj}, and it is worth
saying which, because the list is exactly the list of hypotheses that the
statement carries.

\begin{enumerate}
\item\label{itm:ch31-pole:compact}
 The compactness estimate is
 \Cref{thm:ch22-proper-normalization}: on a fixed dual complex, the
 difference \(\sup_X\varphi-\varphi(x)\) is bounded uniformly over
 \(\PSH(X,\omega)\).  It fails for \(x\) outside every dual complex, and
 \Cref{ex:ch31-type-one-fails} shows that the failure is not repairable.
\item\label{itm:ch31-pole:envelope}
 Continuity of \(P_\omega(f_{\cX})\) is
 \Cref{thm:bfj-envelope-regularity}, which is where residue characteristic
 zero enters; over a field of positive residue characteristic one must
 invoke \Cref{thm:ch23-gjkm-positive-characteristic} instead, with its own
 hypotheses.
\item\label{itm:ch31-pole:orth}
 Orthogonality is \Cref{def:orthogonality-property}, available by
 \Cref{thm:orthogonality-provenance}; the statement that the measure of an
 envelope lives on the contact locus is
 \Cref{prop:ch27-supergradients-on-contact}.
\item\label{itm:ch31-pole:cy}
 The converse uses the supported Calabi--Yau theorem
 \Cref{thm:bfj-supported-cy} and the domination principle
 \Cref{cor:ch25-domination-uniqueness}.
\end{enumerate}

\begin{remark}[Why the supremum is needed]
\label{rem:ch31-why-supremum}
\eqref{eq:finite-pole-support} is a statement about the envelope and not
about an arbitrary competitor.  On the two-component curve of
\Cref{ex:two-component-complete}, with \(S=\{x_0\}\) and \(t_0=0\), the
affine potential with values \((0,0)\) satisfies the constraint and has
\(\MA_\omega=m_0\delta_{x_0}+m_1\delta_{x_1}\), whose support is not
contained in \(S\).  It is the maximality of \(\Phi_{S,t}\), through
orthogonality, that forces the mass onto \(S\).  The reader who
misremembers \eqref{eq:finite-pole-support} as a property of constrained
potentials rather than of the envelope will produce false statements
immediately.
\end{remark}

\subsection{One pole: the divisorial Green function}
\label{sec:ch31-one-pole}

The case \(N=1\) is the non-Archimedean Green function of a divisorial
valuation, already met as \eqref{eq:green-envelope-divisorial}.

\begin{proposition}
\label{prop:ch31-one-pole}
Let \(x\) be a divisorial point lying on some SNC dual complex, and set
\(\varphi_x=\Phi_{\{x\},0}\) as in \eqref{eq:finite-pole-envelope}.  Then
\(\varphi_x\) is continuous, \(\omega\)-psh, \(\varphi_x(x)=0\), and
\begin{equation}
\label{eq:ch31-one-pole-measure}
 \MA_\omega(\varphi_x)=V\,\delta_x .
\end{equation}
Moreover \(\varphi_x\geq0\), with \(\varphi_x=0\) precisely on a set
containing \(x\), and \(\sup_X\varphi_x<\infty\).
\end{proposition}

\begin{proof}
Continuity, pshness and \(\Supp\MA_\omega(\varphi_x)\subseteq\{x\}\) are
\Cref{prop:finite-pole-envelope} with \(N=1\), \(t=0\).  A positive measure
supported at a single point is a multiple of \(\delta_x\), and its mass is
\(V\) by \Cref{thm:ch24-mass}; this gives
\eqref{eq:ch31-one-pole-measure}.  That \(\varphi_x(x)=0\) is the assertion
that the constraint is active; if \(\varphi_x(x)<0\) then
\(\varphi_x-\varphi_x(x)\) would be a strictly larger competitor.  Since
the constant \(0\) is admissible, \(\varphi_x\geq0\).  Finiteness of the
supremum is item \cref{itm:ch31-pole:compact} above.
\end{proof}

\begin{example}[The Green function of a divisorial point on
\(\mathbf P^1\), explicitly]
\label{ex:ch31-p1-green}
Continue \Cref{ex:ch31-p1-blowup}: \(X=\mathbf P^1\), \(L=\cO(1)\),
\(V=1\), \(\mu_0=\delta_{x_{\mathrm G}}\), and \(x_E\) the divisorial point
at tree distance \(1\) from \(x_{\mathrm G}\) in the direction of
\(\xi=\overline{(z=1)}\).  Write \(t(y)\in[0,\infty]\) for the distance from
\(x_{\mathrm G}\) to \(y\) measured along the branch towards \(\xi\), with
\(t(y)=0\) if \(y\) does not lie in that branch.  Then
\begin{equation}
\label{eq:ch31-p1-green-formula}
 \varphi_{x_E}(y)=1-\min\{t(y),1\}
 =\begin{cases}
   1,&y\ \text{outside the branch},\\
   1-t(y),&0\leq t(y)\leq1,\\
   0,&t(y)\geq1 .
  \end{cases}
\end{equation}
Indeed \(\varphi_{x_E}=1-\varphi_E\) in the notation of
\Cref{ex:ch31-p1-blowup}, so
\(\MA_\omega(\varphi_{x_E})=\MA_\omega(-\varphi_E)=\delta_{x_E}=V\delta_{x_E}\)
by \eqref{eq:ch31-p1-two-measures}, and \(\varphi_{x_E}(x_E)=0\).

\emph{That it is the envelope.}  By \Cref{prop:ch31-one-pole} the envelope
\(\Phi_{\{x_E\},0}\) is a continuous \(\omega\)-psh function with
\[
 \MA_\omega\bigl(\Phi_{\{x_E\},0}\bigr)=V\delta_{x_E}=\delta_{x_E},
 \qquad
 \Phi_{\{x_E\},0}(x_E)=0 .
\]
The function \(1-\varphi_E\) has the same two properties, by
\eqref{eq:ch31-p1-two-measures} and
\eqref{eq:ch31-p1-green-formula}.  The measure \(\delta_{x_E}\) is supported
on the dual complex \(\Delta_{\cY}\) of the SNC model \(\cY\), so the
uniqueness clause of \Cref{thm:bfj-supported-cy} applies: two continuous
\(\omega\)-psh functions with this Monge--Amp\`ere measure differ by a
constant.  Both take the value \(0\) at \(x_E\), so the constant is zero and
\(\Phi_{\{x_E\},0}=1-\varphi_E\).

The reader who prefers a direct verification may argue as follows.  By
\Cref{cor:ch23-envelope-by-model-functions} it is enough to bound the
\emph{model} competitors, so let \(\psi\) be an \(\omega\)-psh model
function with \(\psi(x_E)\leq0\), determined on some model dominating
\(\cY\), and write \(t\) for the coordinate along
\([x_{\mathrm G},x_E]\).

Off \(x_{\mathrm G}\) the reference measure \(\mu_0=\delta_{x_{\mathrm G}}\)
vanishes, so \(\Delta\psi=\MA_\omega(\psi)\geq0\) on
\(X^{\An}\smallsetminus\{x_{\mathrm G}\}\).  Consequently, on every branch
\(B\) of the tree hanging off a point \(y\neq x_{\mathrm G}\), the function
\(\psi|_B\) has nonnegative Laplacian on the interior and therefore attains
its maximum at \(y\) (\Cref{prop:ch08-graph-maximum}); in particular the
outgoing slope of \(\psi\) at \(y\) into \(B\) is \(\leq0\), and
\(\psi\leq\psi(x_E)\leq0\) beyond \(x_E\).  Applying this at each interior
point of \((x_{\mathrm G},x_E)\) shows that the two slopes of \(\psi\) along
the segment sum to a nonnegative number there, that is, \(\psi\) is convex
in \(t\) on \([0,1]\).

Next, \(\psi'(0^+)\geq-1\).  Indeed at \(x_{\mathrm G}\),
\[
 0\leq\MA_\omega(\psi)(\{x_{\mathrm G}\})
   =1+\sum_{\nu}d_\nu\psi(x_{\mathrm G})
   \leq1+\psi'(0^+),
\]
because among the outgoing slopes at \(x_{\mathrm G}\) the one along the
segment is \(\psi'(0^+)\) and all the others are \(\leq0\) by the previous
paragraph.  Convexity now gives
\(\psi(x_E)\geq\psi(x_{\mathrm G})+\psi'(0^+)\), so
\[
 \psi(x_{\mathrm G})\leq\psi(x_E)-\psi'(0^+)\leq0+1=1 ,
\]
and, again by convexity and the chord inequality,
\(\psi(t)\leq(1-t)\psi(x_{\mathrm G})+t\psi(x_E)\leq1-t\) for
\(0\leq t\leq1\).  Off the branch, \(\psi\leq\psi(x_{\mathrm G})\leq1\) by
the maximum principle.  Comparing with
\eqref{eq:ch31-p1-green-formula} gives \(\psi\leq\varphi_{x_E}\)
everywhere.

\emph{Numbers.}  \(\sup_X\varphi_{x_E}=1\); \(\varphi_{x_E}(x_E)=0\); the
measure is \(\delta_{x_E}\) of mass \(1=\deg\cO(1)\); and by
\Cref{prop:ch26-graph-energy} on the segment
\(\Delta_{\cY}=[x_{\mathrm G},x_E]\),
\[
 E_\omega(\varphi_{x_E})
 =\int\varphi_{x_E}\,\rd\mu_0-\tfrac12\int|\varphi_{x_E}'|^2\rd s
 =1-\tfrac12=\tfrac12 .
\]
Consistently, \(E_\omega(\varphi_{x_E})=E_\omega(-\varphi_E)+1\cdot V
 =-\tfrac12+1=\tfrac12\), which is
check~\cref{itm:ch31-check:constant}.
\end{example}

\begin{example}[Green functions of the toric divisorial points]
\label{ex:ch31-toric-green}
Let \(X=\mathbf P^n\), \(L=\cO(d)\), \(P=d\Delta_n\), and let
\(x=\xi_{\mathrm G}=\iota(0)\) be the toric Gauss point, the monomial point
over the origin of \(N_{\mathbf R}\).  Take for \(\omega\) the closed model
form of the standard model \(\mathbf P^n_R\) with \(\cO(d)\), whose weight
is \(\check h_P\) and whose curvature is \(d^n\delta_{\xi_{\mathrm G}}\) by
\Cref{ex:ch31-toric-atom-pn}.  Then
\begin{equation}
\label{eq:ch31-toric-green}
 \Phi_{\{\xi_{\mathrm G}\},0}=0,
 \qquad
 \MA_L\bigl(\Phi_{\{\xi_{\mathrm G}\},0}\bigr)
  =d^n\delta_{\xi_{\mathrm G}}=V\delta_{\xi_{\mathrm G}} .
\end{equation}
Indeed the constant \(0\) is \(\omega\)-psh, satisfies the constraint, and
has the right measure; by \Cref{prop:ch31-one-pole} the envelope has measure
\(V\delta_{\xi_{\mathrm G}}\) and vanishes at \(\xi_{\mathrm G}\), and since
\(\delta_{\xi_{\mathrm G}}\) is carried by the dual complex of
\(\mathbf P^n_R\), the uniqueness clause of \Cref{thm:bfj-supported-cy}
forces the two to agree.  In the language of weights, the Green function of
the toric Gauss point \emph{is} the reflected support function
\(\check h_P\), whose subgradient at the origin is all of \(-P\), so that
the whole mass \(n!\vol_M(P)=d^n\) sits at one point.

For a lattice point \(u_0\in N=\mathbf Z^n\) the same computation applies
after translating by any \(t\in T(K)\) with \(\trop(t)=u_0\), which exists
because \(v(K^\times)=\mathbf Z\); the envelope becomes
\(\check h_P(\cdot-u_0)\) relative to the translated reference.  The example
is worth pairing with \Cref{ex:ch31-p1-green}: in the toric case the Green
function is invariant, polyhedral, and equal to the reference weight, while
in the non-toric case of \Cref{ex:ch31-p1-green} it is none of these, yet
both are model functions with a single atom of mass \(V\).
\end{example}

\begin{example}[The hypothesis that the poles be divisorial cannot be
dropped]
\label{ex:ch31-type-one-fails}
Keep \(X=\mathbf P^1\), \(L=\cO(1)\), \(\mu_0=\delta_{x_{\mathrm G}}\), and
let \(x^{(1)}\in\mathbf P^{1,\An}\) be the type-\(1\) point corresponding to
the \(K\)-rational point \(z=1\).  It is not divisorial and lies on no dual
complex.  Consider the analogue of \eqref{eq:finite-pole-envelope} with
\(S=\{x^{(1)}\}\) and \(t=0\):
\[
 \Phi:=\sup\{\varphi\in\PSH(X,\omega)\mid\varphi(x^{(1)})\leq0\} .
\]
For \(m\in\mathbf Z_{>0}\), let \(x_m\) be the divisorial point at tree
distance \(m\) from \(x_{\mathrm G}\) along the branch towards
\(\xi=\overline{(z=1)}\), namely the Gauss point of the disc
\(D(1,|\pi|^m)\).  By \Cref{ex:ch31-p1-green} with \(x_E\) replaced by
\(x_m\),
\[
 \varphi_{x_m}(y)=m-\min\{t(y),m\},
 \qquad
 \MA_\omega(\varphi_{x_m})=\delta_{x_m},
 \qquad
 \sup_X\varphi_{x_m}=m .
\]
Since \(t(x^{(1)})=+\infty\), one has \(\varphi_{x_m}(x^{(1)})=0\) for every
\(m\), so every \(\varphi_{x_m}\) is admissible and
\[
 \Phi\geq\sup_m\varphi_{x_m},
 \qquad
 \Phi(x_{\mathrm G})\geq\sup_m m=+\infty .
\]
Hence \(\Phi\equiv+\infty\) off the branch and
\eqref{eq:finite-pole-envelope} defines nothing.  The failure is exactly the
failure of item \cref{itm:ch31-pole:compact}: the estimate of
\Cref{thm:ch22-proper-normalization} bounds \(\sup_X\varphi-\varphi(x)\)
uniformly for \(x\) in a \emph{fixed} dual complex, and the constant
degenerates as the complex is enlarged.  Numerically, the constant for the
dual complex \([x_{\mathrm G},x_m]\) is \(m\), and it tends to infinity.

Two further remarks.  First, the same computation shows that the
one-parameter family \(\varphi_{x_m}\) has energies
\(E_\omega(\varphi_{x_m})=m-\tfrac m2\cdot1=\tfrac m2\) --- more precisely,
by \Cref{prop:ch26-graph-energy} on \([x_{\mathrm G},x_m]\) with
\(\mu_0=\delta_{x_{\mathrm G}}\) and slope \(-1\) on a segment of length
\(m\), \(E_\omega(\varphi_{x_m})=m-\tfrac12m=\tfrac m2\) --- so the family is
not bounded in energy either, and \Cref{prop:e1-compact-superlevels} does
not apply.  Second, normalizing by the supremum instead produces
\(\varphi_{x_m}-m=-\min\{t,m\}\), a \emph{decreasing} sequence of
\(\omega\)-psh model functions converging pointwise to
\[
 u(y):=-t(y)\in[-\infty,0],
\]
which is \(\omega\)-psh by \Cref{thm:ch21-decreasing-stability} and takes
the value \(-\infty\) exactly at the single type-\(1\) point \(x^{(1)}\).
Its energies are
\(E_\omega(\varphi_{x_m}-m)=E_\omega(\varphi_{x_m})-mV=\tfrac m2-m
 =-\tfrac m2\to-\infty\), so \(u\notin\cE^1(X,\omega)\).  Meanwhile
\(\MA_\omega(-\min\{t,m\})=\delta_{x_m}\to\delta_{x^{(1)}}\) weakly, since
\(x_m\to x^{(1)}\) in \(X^{\An}\): all the mass migrates to the polar point.
The nonpluripolar convention \Cref{def:nonpluripolar-ma} discards it: the
truncation \(\max\{u,-m\}=-\min\{t,m\}\) has measure \(\delta_{x_m}\), and
\(\{u>-m\}=\{t<m\}\) does not contain \(x_m\), so every layer contributes
zero and \(\MA_\omega(u)=0\) has mass \(0<V\).  This is the same mechanism as
\Cref{ex:ch26-infinite-energy}, and it is what
\Cref{prop:bounded-ma-no-polar-mass} excludes as long as the potentials
stay bounded.
\end{example}

\begin{warning}{Divisorial, quasimonomial, and arbitrary points}
\Cref{prop:finite-pole-envelope} is stated for divisorial points lying on an
SNC dual complex.  The proof works verbatim for quasimonomial points of such
a complex, because \Cref{thm:ch22-proper-normalization} is a statement about
the complex and not about rationality.  It fails, by
\Cref{ex:ch31-type-one-fails}, for points of type~\(1\), and it has no
content for points whose center is not a single generic point.  This is the
same trichotomy as in \Cref{sec:ch08-p1-tree}, and it is the reason that
\Cref{part:bfj} normalizes divisorial points by the ground valuation rather
than by any intrinsic quantity.
\end{warning}

\subsection{Weights as Lagrange multipliers}
\label{sec:ch31-multipliers}

\begin{example}[Weights as Lagrange multipliers]
\label{ex:ch31-lagrange}
The values \(t_i\) are constraints and the masses
\(a_i=\MA(\Phi_{S,t})(\{x_i\})\) are their Lagrange multipliers.  Adding a
common constant to \(t\) adds that constant to the envelope and leaves every
\(a_i\) unchanged.  On the quotient \(\mathbf R^N/\mathbf R(1,\ldots,1)\),
the energy is concave and its gradient is the mass vector.  This is the
finite-dimensional shadow of the envelope differentiability theorem.
\end{example}

The two-component curve realizes this statement with \(N=2\) and closed
formulae, and it is the smallest place in the book where a Lagrange
multiplier can be seen as an intersection number.

\begin{example}[The two-pole programme on the two-component curve]
\label{ex:ch31-two-pole-programme}
Keep the data of \Cref{ex:two-component-complete}: \(\Gamma=[x_0,x_1]\) of
length one, \(\mu_0=m_0\delta_{x_0}+m_1\delta_{x_1}\), \(V=m_0+m_1\).  Take
\(S=\{x_0,x_1\}\) and \(t=(t_0,t_1)\), and write \(\delta=t_1-t_0\).

\emph{The envelope.}  By \eqref{eq:ch31-segment-psh} an \(\omega\)-psh model
function is convex with \(u'(0^+)\geq-m_0\) and \(u'(1^-)\leq m_1\); among
those satisfying \(u(x_0)\leq t_0\) and \(u(x_1)\leq t_1\), the largest is
affine, and one finds
\begin{equation}
\label{eq:ch31-two-pole-envelope}
 \Phi_{S,t}=
 \begin{cases}
  \text{affine with values }(t_0,t_1), & -m_0\leq\delta\leq m_1,\\
  \text{affine with values }(t_0,t_0+m_1), & \delta>m_1,\\
  \text{affine with values }(t_1+m_0,t_1), & \delta<-m_0 .
 \end{cases}
\end{equation}
In the first regime both constraints are active; in the second only the
constraint at \(x_0\); in the third only that at \(x_1\).

\emph{The multipliers.}  In the first regime,
\eqref{eq:two-component-ma} gives
\begin{equation}
\label{eq:ch31-two-pole-masses}
 a_0=m_0+\delta,\qquad a_1=m_1-\delta,\qquad a_0+a_1=V ,
\end{equation}
and as \(\delta\) runs over \([-m_0,m_1]\) the pair \((a_0,a_1)\) runs over
the whole simplex \(\{a_i\geq0,\ a_0+a_1=V\}\).  That is the converse half
of \Cref{prop:finite-pole-envelope} in this case, made explicit: every
admissible atomic measure is realized, and the map \(\delta\mapsto(a_0,a_1)\)
is a bijection onto the simplex.  In the second regime
\((a_0,a_1)=(V,0)\) and in the third \((0,V)\); the constraint vector is
therefore \emph{not} unique, which is the last sentence of
\Cref{prop:finite-pole-envelope}.

\emph{The energy and its gradient.}  In the first regime, by
\eqref{eq:ch26-two-vertex-energy},
\begin{equation}
\label{eq:ch31-two-pole-energy}
 \mathcal F(t):=E_\omega(\Phi_{S,t})=m_0t_0+m_1t_1-\tfrac12(t_0-t_1)^2 ,
\end{equation}
and
\[
 \frac{\partial\mathcal F}{\partial t_0}=m_0+\delta=a_0,
 \qquad
 \frac{\partial\mathcal F}{\partial t_1}=m_1-\delta=a_1 .
\]
Thus the gradient of the energy in the constraint variables is the mass
vector, which is \Cref{ex:ch31-lagrange}.  The Hessian of \eqref{eq:ch31-two-pole-energy}
is
\(-\bigl(\begin{smallmatrix}1&-1\\-1&1\end{smallmatrix}\bigr)\), negative
semidefinite with kernel \(\mathbf R(1,1)\); so \(\mathcal F\) is concave and, along
the diagonal direction, affine with slope \(m_0+m_1=V\).  That is
\eqref{eq:bfj-energy-translation} in coordinates.

\emph{Numbers.}  With \(m_0=2\), \(m_1=3\) as in
\Cref{ex:ch31-two-component-numbers} and \(t=(0,0)\): \(\delta=0\),
\((a_0,a_1)=(2,3)\), \(\mathcal F=0\).  With \(t=(0,3)\): \(\delta=3=m_1\),
\((a_0,a_1)=(5,0)\), \(\mathcal F=9-\tfrac92=\tfrac92\); the envelope is the
affine function with values \((0,3)\), whose measure is \(5\delta_{x_0}\),
and which is the potential of \Cref{ex:two-component-complete} shifted by
\(+3\), so its energy is \(-\tfrac{21}2+3\cdot5=\tfrac92\), as it must be.
With \(t=(0,5)\): \(\delta=5>3\), the envelope is unchanged from the
previous case, and \(\mathcal F\) is likewise unchanged; the second constraint has
become inactive, and the multiplier \(a_1\) has stuck at \(0\).  Along the
segment \(t_1\in[3,5]\) the function \(\mathcal F\) is therefore constant, which is
where the concave function \(\mathcal F\) of \eqref{eq:ch31-two-pole-energy} ceases
to be given by that formula and becomes affine with zero slope: the
finite-dimensional shadow of the fact that the derivative of the energy is a
\emph{measure}, and that the measure has stopped charging \(x_1\).
\end{example}

\begin{remark}[The programme in general, and what it is a shadow of]
\label{rem:ch31-programme-general}
For general \(N\) the picture of \Cref{ex:ch31-two-pole-programme} persists.
The function
\(\mathcal F(t)=E_\omega(\Phi_{S,t})\) on \(\mathbf R^N\) is concave, satisfies
\(\mathcal F(t+c\mathbf 1)=\mathcal F(t)+cV\), and by
\Cref{thm:energy-envelope-differentiability} together with
\Cref{prop:ch27-supergradients-on-contact} its superdifferential at \(t\) is
the set of mass vectors of Monge--Amp\`ere measures of \(\Phi_{S,t}\)
carried by \(S\); when the envelope is differentiable in the model
directions, \(\nabla\mathcal F(t)=(a_1,\ldots,a_N)\).  Its Legendre transform is
the dual energy of \Cref{def:dual-energy-measure} restricted to measures
carried by \(S\), and the correspondence
\(t\leftrightarrow a\) is the finite-dimensional case of the
Calabi--Yau correspondence \Cref{thm:bfj-supported-cy}.  The
\(N\)-dimensional concave programme is genuinely a shadow and not a proof:
concavity of \(\mathcal F\) on \(\mathbf R^N\) does not by itself give the
existence of a maximizer in the infinite-dimensional problem, which is why
\Cref{chap:bfj-finite-energy-measures} needs coercivity.
\end{remark}

\subsection{The envelope need not be a model function}
\label{sec:ch31-not-model}

\Cref{prop:finite-pole-envelope} reduces a prescribed atomic measure to a
finite-dimensional convex problem, but it does not say that \(\Phi_{S,t}\)
is a model function.  Even for \(N=1\), its modelness is tied to finite
generation of a section algebra and may fail.

\begin{remark}[What modelness would mean, and what is known]
\label{rem:ch31-modelness}
By \Cref{thm:bfj-skeletal-characterization}, \(\Phi_{S,t}\) is a model
function precisely when it is \(\mathbf Q\)-piecewise linear on some dual
complex and factors through the corresponding retraction.  The examples of
\Cref{sec:ch31-one-pole} are all of this kind: \eqref{eq:ch31-p1-green-formula}
is piecewise linear with two pieces, and the toric envelope of
\Cref{ex:ch31-toric-green} is a support function.  That is a feature of
dimension one and of the invariant category, not a theorem.  In general the
question is whether the graded algebra
\(\bigoplus_m H^0(X,mL\otimes\cI_m)\) attached to the constraint is finitely
generated; when it is, the envelope is a model function, and when it is not,
it need not be.  Boucksom--Jonsson exhibit a divisorial Green function that
is not \(\mathbf Q\)-piecewise linear \cite[Theorem~A]{BJGreen24}; the
statement is proved there for the trivially valued theory and is discussed
in \Cref{part:bj}.  Nothing in the present chapter asserts modelness beyond
the cases in which it is computed.
\end{remark}

\begin{warning}{Continuity is a theorem; modelness is not}
The continuity of \(\Phi_{S,t}\) is
\Cref{thm:bfj-envelope-regularity} and is available under the standing
hypotheses.  Modelness is a separate and much stronger assertion, and no
theorem of \Cref{part:bfj} provides it.  A proof that silently replaces a
continuous envelope by a model function is using finite generation without
saying so.
\end{warning}

\section{Simple degenerations}
\label{sec:ch31-degenerations}

The last family concerns metrics that are \emph{canonical}: distinguished by
a symmetry or by a functional equation rather than by a choice.  In each of
the cases below the canonical metric is computed in closed form, its
Monge--Amp\`ere measure is identified, and the mass is checked against an
intersection number.

\subsection{Good reduction: the degenerate case of everything}
\label{sec:ch31-good-reduction}

\begin{example}[Good reduction]
\label{ex:ch31-good-reduction}
Suppose \(X\) has a smooth projective model \(\cX\) over \(R\) with
\(\cX_s\) irreducible, and let \(\cL\) be a model of \(L\) on \(\cX\), nef
relatively over \(R\).  The dual complex \(\Delta_{\cX}\) is a single point
\(x_0\), the divisorial point of \(\cX_s\), and
\begin{equation}
\label{eq:ch31-good-reduction-measure}
 \MA_\omega(0)=b_{\cX_s}\bigl(c_1(\cL)|_{\cX_s}^{\,n}\bigr)\delta_{x_0}
 =(L^n)\,\delta_{x_0}=V\delta_{x_0},
\end{equation}
because \(\cX_s\) is reduced and \((c_1(\cL)|_{\cX_s}^{\,n})=(L^n)\) by
flatness.  Every model function determined on \(\cX\) is constant, so
\(\MA_\omega\) is constant on that model; nonconstant model functions come
from blowups, as in \Cref{ex:ch31-p1-blowup}, and produce measures with
atoms at the divisorial points of the exceptional divisors.

The Green function of \Cref{prop:ch31-one-pole} at \(x=x_0\) is the constant
\(0\), and \(\MA_\omega(\varphi_{x_0})=V\delta_{x_0}\) trivially.  If
\(X=A\) is an abelian variety with good reduction and \(L\) is ample and
symmetric, the canonical metric is the model metric of the N\'eron model, so
\eqref{eq:ch31-good-reduction-measure} \emph{is} the canonical measure:
a single atom of mass \((L^g)\).  This is
item~\cref{itm:ch19-scope:good} of \Cref{sec:ch19-canonical}, and it is the
degenerate case \(r=0\) of \Cref{thm:ch19-canonical-ma}.
\end{example}

\begin{remark}[Why good reduction is not a special case worth much]
\label{rem:ch31-good-reduction-scope}
Nothing in \Cref{ex:ch31-good-reduction} tests a normalization beyond the
mass, because a measure supported at one point is determined by its mass.
It is included because it is the boundary case against which the others are
read: the skeleton has dimension zero, the Laplacian is the zero operator,
the energy of a model function determined on \(\cX\) is \(cV\) for the
constant \(c\), and every table entry degenerates correctly.  A
normalization that fails here fails everywhere.
\end{remark}

\subsection{A nodal degeneration of curves}
\label{sec:ch31-nodal}

Let \(\mathcal X\to\Spec K^\circ\) be regular semistable with special fibre
\(E_0+E_1\), where the components meet transversely at one node.  Let
\(\mathcal L\) be relatively nef with
\(\deg(\mathcal L|_{E_i})=m_i\).  The associated model metric has
\begin{equation}
                   c_1(L,\|\cdot\|_{\mathcal L})
                   =m_0\delta_{x_0}+m_1\delta_{x_1}.
\label{eq:nodal-reference-measure}
\end{equation}
Moving all mass to \(x_0\) changes the weight by the affine graph function
of \Cref{ex:two-component-complete}.  Algebraically, the vertical
divisor changes the degrees on \(E_0,E_1\) by \((+m_1,-m_1)\); analytically,
its slope changes the two atomic coefficients by the same amounts.  This is
the simplest complete instance of the equality between vertical
intersection and graph Laplacian.

Written out, the ledger is the following.  Let
\(D=m_1E_1\) be the vertical divisor with \(\varphi_D(x_0)=0\),
\(\varphi_D(x_1)=m_1\) (the sign convention being
\eqref{eq:appJ-model-divisor}: \(D\) effective gives
\(\varphi_D\geq0\)).  On the regular semistable surface \(\cX\) one has
\[
 (E_0\cdot E_1)=1,\qquad
 (E_0^2)=-1,\qquad (E_1^2)=-1,
\]
the last two from \((\cX_s\cdot E_i)=0\) with \(\cX_s=E_0+E_1\).  Hence
\begin{equation}
\label{eq:ch31-nodal-ledger}
 \begin{aligned}
 \deg\bigl((\cL\otimes\cO(D))|_{E_0}\bigr)
  &=m_0+m_1(E_1\cdot E_0)=m_0+m_1,\\
 \deg\bigl((\cL\otimes\cO(D))|_{E_1}\bigr)
  &=m_1+m_1(E_1^2)=0 .
 \end{aligned}
\end{equation}
The two numbers are exactly the coefficients of
\eqref{eq:two-component-ma} at \((a,b)=(0,m_1)\), which is
\(\varphi_D\) itself; and they are the masses of the envelope
\(\Phi_{\{x_0\},0}\) of \Cref{ex:ch31-two-pole-programme}.  So the same pair
\((V,0)\) has now been computed three ways: by graph slopes, by vertical
intersection numbers, and by a constrained supremum.  Their agreement is
\Cref{thm:ch20-ddc-exact} in the smallest case, and it is
check~\cref{itm:ch31-check:vertical}.

\begin{example}[A node of thickness \(a\)]
\label{ex:ch31-nodal-thickness}
Suppose instead that the node has completed local ring
\(R[[x,y]]/(xy-\pi^a)\) with \(a\in\mathbf Z_{>0}\), so that the edge of
\(\Gamma\) has length \(a\) (\Cref{ex:ch08-An-edge}).  Then \(\cX\) is
regular only after resolving into a chain of \(a+1\) components; on the
normal two-component model the Laplacian of an affine function
with endpoint values \(a_0,a_1\) is
\(\tfrac{a_1-a_0}{a}\delta_{x_0}+\tfrac{a_0-a_1}{a}\delta_{x_1}\), and
\eqref{eq:two-component-ma} becomes
\[
 \MA_\omega(u)=\Bigl(m_0+\frac{a_1-a_0}{a}\Bigr)\delta_{x_0}
              +\Bigl(m_1+\frac{a_0-a_1}{a}\Bigr)\delta_{x_1}.
\]
Moving all the mass to \(x_0\) now requires \(a_1-a_0=am_1\), and the energy
of the resulting potential, normalized by \(a_1=0\), is
\(-m_0m_1a-\tfrac12m_1^2a\) by \Cref{prop:ch26-graph-energy}.  Both
quantities are proportional to \(a\), the masses are not: the lengths carry
the valuation normalization and the intersection numbers do not.  That is
check~\cref{itm:ch31-check:normalization} in its simplest form, and it
anticipates \Cref{sec:ch31-base-change}.  For \(a=3\) the resolution is the
chain of \Cref{ex:ch20-chain-three}, and one may verify the same numbers on
that model after subdividing.
\end{example}

\subsection{Degenerating elliptic curves}
\label{sec:ch31-elliptic}

The Tate curve was computed in \Cref{ex:ch31-tate-bfj}.  We record here what
it says as a statement about degenerations, and how the two elliptic cases
sit relative to each other.

\begin{example}[The two elliptic regimes side by side]
\label{ex:ch31-elliptic-dichotomy}
Let \(E\) be an elliptic curve over \(K\) and \(L\) of degree \(d>0\); let
\(\|\cdot\|_{\mathrm{can}}\) be the canonical metric, characterized by
\([2]^*\|\cdot\|_{\mathrm{can}}=\|\cdot\|_{\mathrm{can}}^{\otimes4}\) for
symmetric \(L\).
\begin{center}
\small
\begin{tabular}{p{.20\textwidth}|p{.22\textwidth}|p{.24\textwidth}|p{.24\textwidth}}
\toprule
&\textbf{good reduction}&\textbf{multiplicative, \(v(q)=\ell\)}
 &\textbf{where}\\
\midrule
skeleton&a point&circle of circumference \(\ell\)
 &\Cref{ex:ch08-tate-curve}\\
canonical weight on the skeleton&constant
 &\(\tfrac{d}{2\ell}t(t-\ell)\)
 &\eqref{eq:ch31-tate-potential}\\
canonical measure&\(d\,\delta_{x_0}\)
 &\(\tfrac d\ell\,\rd t\)
 &\eqref{eq:ch19-tate-measure}\\
total mass&\(d\)&\(d\)&\eqref{eq:examples-global-normalization}\\
is it a model metric?&yes&no; only uniform piecewise affine model
 approximations&\Cref{ex:ch31-tate-bfj}\\
energy relative to the balanced model metric&\(0\)
 &\(-\tfrac{\ell d^2}{24}\)&\eqref{eq:ch31-tate-energy}\\
\bottomrule
\end{tabular}
\end{center}
The two columns are not deformations of one another: the good-reduction
measure is a single atom, the multiplicative one is diffuse, and the energy
in the second column tends to \(0\) as \(\ell\downarrow0\), which is the
correct limiting behaviour but is not a proof of anything about families,
since \(\ell\) is not a continuous parameter over a fixed \(K\).  What the
table records is that a single normalization
\eqref{eq:examples-global-normalization} produces the correct mass \(d\) in
both regimes.
\end{example}

\begin{remark}[The canonical metric is not a model metric in general]
\label{rem:ch31-tate-not-model}
By \eqref{eq:ch31-tate-potential} the canonical weight is a genuine
quadratic in \(t\), hence not piecewise affine; so the canonical metric of a
Tate curve is a continuous semipositive metric that is not a model metric.
It is a uniform decreasing limit of model metrics along the nested dyadic
subdivisions described after \eqref{eq:ch31-tate-energy-approx}; the
associated measures are those of \eqref{eq:ch19-tate-model-approx},
which is exactly the hypothesis under which
\Cref{thm:ch19-cld-bfj-comparison}\,(B) applies.  This is the second
occurrence in this chapter of a potential in the continuous row but not the
model row of \Cref{sec:ch24-four-products}, the first being
\Cref{ex:ch31-chain-lebesgue}, and it is the one that matters for
arithmetic applications.
\end{remark}

\subsection{Totally degenerate abelian varieties}
\label{sec:ch31-abelian}

Let \(A\) be a totally degenerate abelian variety over an algebraically
closed complete nontrivially valued field.  Raynaud uniformization produces
a canonical skeleton
\[
                         \Sigma_A=N_\mathbf R/\Lambda,
\]
a compact real torus.  Translation-invariant semipositive metrics correspond
to convex quasi-periodic functions on \(N_\mathbf R\).  Their
Chambert--Loir measures are the descended real Alexandrov measures.  The
canonical metric of an ample symmetric line bundle has constant real
Monge--Amp\`ere density, hence
\begin{equation}
               c_1(L,\|\cdot\|_{\mathrm{can}})^n
               =(L^n)\,m_{\mathrm{Haar}},
\label{eq:abelian-canonical-haar}
\end{equation}
where Haar measure has total mass one.  Prescribing a smooth positive density
on \(\Sigma_A\) becomes a periodic real Monge--Amp\`ere equation.  The
invariant tropical theory and its energy classes are developed uniformly in
\cite{BGJKK25}.

Unwinding \eqref{eq:abelian-canonical-haar} gives the numbers.  By
\Cref{thm:ch19-canonical-structure} the canonical weight is
\(\tfrac12q_L\) up to a linear term, where \(q_L\) is the quadratic form
attached to the polarization by \eqref{eq:ch19-qL}; by
\eqref{eq:real-ma-density} its real Monge--Amp\`ere measure is
\begin{equation}
\label{eq:ch31-abelian-density}
 \MA_{\mathbf R}\bigl(\tfrac12q_L\bigr)
 =\det\bigl(D^2(\tfrac12q_L)\bigr)\rd x=\det(B)\,\rd x,
\end{equation}
where \(B\) is the Gram matrix of \(q_L\) in a basis of
\(\tau(\Lambda)\), and the non-Archimedean measure is \(g!\det(B)\rd x\) on a
torus of lattice volume one, of total mass
\begin{equation}
\label{eq:ch31-abelian-mass}
 g!\det(B)=(L^g),
\end{equation}
the last equality being the Riemann relation \((L^g)=g!\chi(L)\) together
with \(\chi(L)=d_1\cdots d_g=\det(B)\), the \(d_i\) being the elementary
divisors of \(B\) and their product its determinant by Smith normal form.
\Cref{ex:ch19-abelian-surface-degenerate} performs this computation
for \(g=2\) with
\(B=\bigl(\begin{smallmatrix}2&1\\1&2\end{smallmatrix}\bigr)\), obtaining
density \(6\,\rd x\) and mass \(6=(L^2)\) for a polarization of type
\((1,3)\); the principally polarized case \(B=I\) gives density
\(2\,\rd x\) and mass \(2\).

\begin{example}[An abelian surface, transported to the BFJ side]
\label{ex:ch31-abelian-surface-bfj}
Keep the data of \Cref{ex:ch19-abelian-surface-degenerate}:
\(A=\mathbf G_{m,K}^{2,\An}/\Lambda\) with
\(\tau(\Lambda)=\mathbf Z^2\), skeleton
\(\Sigma_A=\mathbf R^2/\mathbf Z^2\) of lattice volume \(1\), and
\(q_L(x)=x^{\mathsf T}Bx\) with
\(B=\bigl(\begin{smallmatrix}2&1\\1&2\end{smallmatrix}\bigr)\).  Then
\[
 \MA_L(\|\cdot\|_{\mathrm{can}})=6\,\rd x
 \quad\text{on }\Sigma_A,
 \qquad\text{mass }6=(L^2),
\]
and by \Cref{thm:ch19-cld-bfj-comparison}\,(B), together with
\Cref{cor:ch19-transfer}\,\cref{itm:ch19-tr:abelian}, this is simultaneously
the BFJ measure, the Chambert--Loir measure and the CLD Bedford--Taylor
product.  Three checks are available and all three are performed in
\Cref{ex:ch19-abelian-surface-degenerate}: the polarization type
\((1,3)\) gives \(\chi(L)=3\) and \((L^2)=2!\cdot3=6\); the restriction to
the closed geodesic in the direction \(w=(1,0)\) has degree
\(w^{\mathsf T}Bw=2\); and in the direction \(w=(1,1)\) it has degree \(6\).

For the BFJ theory the transfer is worth stating explicitly.  The canonical
metric is a uniform limit of semipositive model metrics --- concretely, of
the piecewise affine interpolations of \(\tfrac12q_L\) on the subdivision of
\(\mathbf R^2/\mathbf Z^2\) by the lattice \(\tfrac1N\mathbf Z^2\) --- so
it lies in the second row of \Cref{sec:ch24-four-products}; its energy
relative to the model metric of the \(N=1\) interpolation is finite; and its
Monge--Amp\`ere measure has full mass \(6\), so it lies in
\(\cE^1(A,\omega)\) with the full-mass property of
\Cref{thm:ch26-e1-structure}.  Nothing beyond bounded potentials is
involved.
\end{example}

\begin{remark}
\label{rem:ch31-abelian-not-restriction}
The real torus is a skeleton, but the proof is not simply ``restrict BFJ to
the skeleton.''  Raynaud uniformization identifies invariant metrics and
their intersection products with convex data.  Without invariance, a metric
on \(A^{\An}\) contains information not determined by its restriction to
\(\Sigma_A\).

\Cref{ex:ch31-noninvariant} is the smallest illustration of the last
sentence, on \(\mathbf P^1\) rather than on an abelian variety; the same
construction applies verbatim to \(A\), by blowing up a non-torus-fixed
closed point of a special fibre.  The correct general statement of the
relationship between \(A^{\An}\) and \(\Sigma_A\) is
\Cref{thm:ch19-raynaud}\,\cref{itm:ch19-ray:skeleton}, a deformation
retraction of topological spaces, which is not a statement about psh
functions.
\end{remark}

\begin{warning}{Two meanings of ``canonical''}
On an abelian variety, \emph{canonical metric} means the metric fixed by the
multiplication maps, computed in \Cref{sec:ch19-canonical}.  In the
degeneration literature the same word sometimes denotes the limit of a
family of Kähler--Einstein metrics.  The two agree for abelian varieties and
do not agree in general.  \Cref{rem:ch19-canonical-two-meanings} separates
them; nothing in \Cref{part:bfj} uses the second.
\end{warning}

\subsection{Base change: what is invariant and what is not}
\label{sec:ch31-base-change}

Every length in this chapter was measured with the normalization
\(v(\pi)=1\).  The following example records precisely which of the computed
quantities depend on that choice.

\begin{example}[Ramified base change on the two-component curve]
\label{ex:ch31-base-change}
Let \(K'/K\) be totally ramified of degree \(e\), with uniformizer
\(\pi'\), \(\pi=(\pi')^e\) up to a unit.  Two normalizations are possible
for the absolute value on \(K'\), and they must not be mixed
(\Cref{sec:appJ-valuations}):
\begin{enumerate}
\item\label{itm:ch31-bc:extend} \emph{Extend} \(v\), so that
 \(v(\pi')=1/e\); or
\item\label{itm:ch31-bc:renormalize} \emph{renormalize}, so that
 \(v'(\pi')=1\), that is, \(v'=e\,v\) on \(K\).
\end{enumerate}
Take \(X\) with the two-component semistable model of
\Cref{sec:ch31-nodal}, node of thickness \(1\), and let
\(X'=X\otimes_KK'\).  The minimal semistable model of \(X'\) has special
fibre a chain of \(e+1\) components, and the skeleton is a segment.
\begin{center}
\small
\begin{tabular}{p{.30\textwidth}|p{.20\textwidth}|p{.20\textwidth}|p{.20\textwidth}}
\toprule
\textbf{Quantity}&\textbf{over \(K\)}
 &\textbf{over \(K'\), convention \cref{itm:ch31-bc:extend}}
 &\textbf{over \(K'\), convention \cref{itm:ch31-bc:renormalize}}\\
\midrule
edge length of \(\Gamma\)&\(1\)&\(1\)&\(e\)\\
number of vertices&\(2\)&\(e+1\)&\(e+1\)\\
\(\deg L\), hence \(V\)&\(m_0+m_1\)&\(m_0+m_1\)&\(m_0+m_1\)\\
masses of \(\MA_\omega(0)\)&\((m_0,m_1)\)&\((m_0,0,\ldots,0,m_1)\)
 &\((m_0,0,\ldots,0,m_1)\)\\
potential solving \(\MA=V\delta_{x_0}\)&\((-m_1,0)\)
 &\((-m_1,0)\) at the ends&\((-em_1,0)\) at the ends\\
energy \eqref{eq:two-component-energy}&\(-m_0m_1-\tfrac12m_1^2\)
 &unchanged&multiplied by \(e\)\\
\bottomrule
\end{tabular}
\end{center}
The general statements behind the table are
\Cref{prop:base-change-model-metrics} and \Cref{thm:base-change-ma}: the
Monge--Amp\`ere measure commutes with base change, and the total mass, being
an intersection number, is insensitive to it.  What changes is the metric on
the skeleton, and therefore every quantity built from slopes: potentials,
Dirichlet energies, capacities.

The rule to remember is that \emph{numbers computed by intersection theory
are invariant and numbers computed by integrating slopes are not}.  A
computation that produces a base-change-dependent intersection number, or a
base-change-independent Dirichlet energy, contains an error.
\end{example}

\begin{remark}[The Gauss extension]
\label{rem:ch31-gauss}
A different change of base, the Gauss extension \eqref{eq:gauss-extension},
turns a trivially valued field into a nontrivially valued one and is the
bridge used in \Cref{part:bj}.  It is recorded here only to say that it is
not a base change of the kind tabulated in \Cref{ex:ch31-base-change}: it
changes the ground field and the analytic space simultaneously, and the
compatibility statement is \Cref{prop:gauss-model-compatibility}.  The
examples of the present chapter are all over a fixed nontrivially valued
\(K\).
\end{remark}

\section{The BFJ--CLD compatibility table}
\label{sec:ch31-table}

\subsection{What a row asserts, and what it does not}
\label{sec:ch31-table-reading}

\Cref{thm:ch19-cld-bfj-comparison} compares two theories on two loci and
lists four statements it does not prove;
\Cref{sec:ch19-limitations-table} lists fourteen assertions that fail, are
undefined, or are open for singular potentials.  Neither is organized
theorem by theorem, and that is what is supplied here.  Each row of the
tables below names one statement of \Cref{part:bfj} and one statement of
\Cref{part:cld}, says on which locus they agree, records the hypotheses each
side needs, and classifies the relation.

The classification uses four symbols, and the whole value of the tables lies
in keeping them apart.

\begin{description}[style=nextline,leftmargin=0pt,labelindent=0pt,
  labelsep=0pt,itemsep=.3\baselineskip,font=\normalfont\bfseries]
\item[\(=\) equality]
The two statements are about the same objects on the stated locus and assert
the same thing, with all normalizations of \Cref{app:normalizations} taken
into account.  A proof is cited.
\item[\(\Rightarrow\) one-way implication]
The BFJ statement implies the CLD statement on the stated locus, and the
converse is not proved in this book.  The cell says where the converse
fails or is open.
\item[\(\Leftarrow\) one-way implication]
The CLD statement implies the BFJ statement, and not conversely.
\item[\(?\) unknown]
Both statements are meaningful on the stated locus and neither implication
is proved here.  A pointer to the limitations table or to an open problem is
given.
\item[\(\varnothing\) no counterpart]
One side has no statement at all: either the objects are not defined, or the
theory does not address the question.  This is not the same as \(?\), and
conflating the two produces false open problems, as
\Cref{sec:ch19-limitations-reading} insists.
\end{description}

Three standing conventions apply to every row.  First, \emph{hypotheses}
means the hypotheses actually used in the proof cited in that row, not the
standing hypotheses of the chapter in which the statement appears.  Second,
\emph{locus} always names a class of potentials or metrics, never a subset
of \(X^{\An}\).  Third, every row is consistent with
\Cref{thm:ch19-cld-bfj-comparison}\,(C) and with
\Cref{sec:ch19-limitations-table}: where a row would strengthen either, it
is marked \(?\) and the corresponding row \textup{(L\(k\))} of the
limitations table is cited.  No row of the present tables asserts anything
that \Cref{chap:cld-examples} declines to assert.

\begin{warning}{The tables compare statements, not theories}
A reader who reads down the \(=\) column and concludes that BFJ and CLD are
the same theory has misread the tables.  Nineteen of the rows below concern
model or continuous metrics, where the two constructions provably agree; the
rows that concern singular potentials, global positivity, or existence
theorems are \(?\) or \(\varnothing\) without exception.  The correct summary
is \Cref{sec:ch31-table-not}.
\end{warning}

\subsection{Table I: classes of positive objects}
\label{sec:ch31-table-classes}

\begin{center}
\footnotesize
\begin{tabular}{p{.18\textwidth}|p{.175\textwidth}|p{.155\textwidth}|p{.22\textwidth}|p{.10\textwidth}}
\toprule
\textbf{BFJ statement}&\textbf{CLD/GK statement}&\textbf{Locus of
 agreement}&\textbf{Hypotheses each side needs}&\textbf{Relation}\\
\midrule

(C1) \(\varphi\) is \(\theta\)-psh, i.e.\ \(\theta+\ddc\varphi\) relatively
 nef (\Cref{def:ch20-semipositive}, \Cref{thm:ch21-model-consistency})
 &semipositive formal metric \(=\) positive Chern current \(=\) positive
  Chern delta-form \(=\) functorially psh \(=\) curve test
  (\Cref{thm:common-model-core},
   \Cref{thm:gk-formal-positivity-comparison})
 &model functions on \(X^{\An}\)
 &BFJ: \eqref{eq:ch19-bfj-hypotheses}, for the existence of models over
  \(R\).  CLD: \(K\) algebraically closed and \(X\) proper reduced; the
  algebraic closure is removed by faithfully flat descent
 &\(=\)\\

(C2) \(\varphi\) continuous \(\theta\)-psh, i.e.\ a uniform limit of
 \(\theta\)-psh model functions (\Cref{thm:ch21-continuous-zhang})
 &\(\|\cdot\|_\varphi\) continuous semipositive in Zhang's sense
  (\Cref{sec:metrics-semipositive})
 &continuous potentials
 &BFJ: density of model functions,
  \Cref{prop:metrics-model-function-density}.  CLD: none beyond
  properness of \(X\)
 &\(=\)\\

(C3) same as (C2)
 &\(c_1(L,\|\cdot\|_\varphi)\geq0\) as a CLD current
  (\Cref{prop:bt-current-stability})
 &continuous potentials
 &BFJ: as in (C2).  CLD: weak closedness of positivity.  The converse fails
  to be provable because positivity of a current is a condition from above;
  row \textup{(L4)}
 &\(\Rightarrow\)\\

(C4) \(\varphi\) bounded \(\theta\)-psh (\Cref{def:ch21-bfj-psh})
 &\(\varphi\) locally CLD-psh (\Cref{def:ch15-hierarchy})
 &none proved
 &Both sides are defined; neither implication is proved.  Forward:
  \Cref{thm:ch19-cld-bfj-comparison}\,\cref{itm:ch19-cb:scope1}.  Backward:
  \cref{itm:ch19-cb:scope2}; row \textup{(L10)},
  \Cref{prob:frontier-cld-global-metric}
 &\(?\)\\

\bottomrule
\end{tabular}
\end{center}

\noindent
The table continues with three families in which the comparison extends
beyond the model locus.  Equality of the two psh \emph{classes} is claimed
only in rows whose hypotheses include the required regularization theorem.

\begin{center}
\footnotesize
\begin{tabular}{p{.18\textwidth}|p{.175\textwidth}|p{.155\textwidth}|p{.22\textwidth}|p{.10\textwidth}}
\toprule
\textbf{BFJ statement}&\textbf{CLD/GK statement}&\textbf{Locus of
 agreement}&\textbf{Hypotheses each side needs}&\textbf{Relation}\\
\midrule

(C5) \(\varphi\) continuous \(\theta\)-psh on a \emph{curve}
 &the associated local metric weights have positive current and hence are
  Thuillier-subharmonic
  (\Cref{thm:wanner-curve-comparison},
   \Cref{thm:ch19-cld-bfj-comparison})
 &continuous semipositive metrics on \(C^{\An}\), \(\dim C=1\)
 &BFJ: \(\{\theta\}\) ample of positive degree.  Current/Thuillier:
  \(K\) algebraically closed, complete, nontrivially valued
  \cite[Corollary~4.8]{Wanner19}.  In this same scope, membership in the
  first ordinary-smooth decreasing-limit stage follows from
  \Cref{thm:p063-current-psh-ordinary-regularization}
 &\(\Rightarrow\)\\

(C6) invariant \(\theta\)-psh toric potential, i.e.\ \(F\) convex with
 recession \(\check h_P\) (\Cref{thm:toric-bfj-alexandrov})
 &psh \(=\) functorially psh \(=\) functorially delta-psh \(=\) convexity of
  \(F\) (\Cref{thm:toric-positivity},
   \Cref{cor:ch19-toric-convexity-test})
 &torus-invariant piecewise smooth metrics
 &Both sides need \emph{invariance}; it cannot be dropped, by
  \Cref{ex:ch31-noninvariant}.  CLD additionally needs a complete fan for
  the last equivalence
 &\(=\)\\

(C7) \(\varphi\) an \(\mathbf R\)-PL model function
 &residue-nef \(=\) classically psh \(=\) CLD-v2 psh on the relative
  interior (\Cref{thm:verified-pl-psh-equivalence},
  \cite[Theorem~4.3]{GR25Classical})
 &\(\mathbf R\)-PL functions on a good strict analytic space
 &The reverse CLD-v2 implication at a relative-boundary point is not
  proved; see the boundary clause of
  \Cref{thm:verified-pl-psh-equivalence}
 &\(=\) interior, \(\Rightarrow\) boundary\\
\bottomrule
\end{tabular}
\end{center}

\subsection{Table II: measures and operators}
\label{sec:ch31-table-main}

\begin{center}
\footnotesize
\begin{tabular}{p{.185\textwidth}|p{.175\textwidth}|p{.155\textwidth}|p{.215\textwidth}|p{.10\textwidth}}
\toprule
\textbf{BFJ statement}&\textbf{CLD/GK statement}&\textbf{Locus of
 agreement}&\textbf{Hypotheses each side needs}&\textbf{Relation}\\
\midrule

(M1) mixed model measure
 \(\theta_1\wedge\cdots\wedge\theta_n
   =\sum_Eb_E(\cdots)\delta_{x_E}\),
 \eqref{eq:model-mixed-measure}, \Cref{def:mixed-model-measure}
 &CLD Bedford--Taylor product of formal metrics
  (\Cref{thm:cld-formal-intersection}), equal to the Chambert--Loir
  measure (\Cref{thm:cld-recovers-cl}) and to the delta-form product
  (\Cref{thm:gk-delta-cl-measure})
 &\(n\)-tuples of model functions
 &BFJ: \(K\) discretely valued, so that the defect factor of
  \Cref{thm:cld-formal-intersection} is the multiplicity \(b_E\).  CLD:
  \(\cX\) relatively normal, \(X\) proper reduced
 &\(=\)\\

(M2) total mass \(=V=(\{\theta\}^n)\)
 (\Cref{thm:ch24-mass}, \eqref{eq:mixed-total-mass})
 &degree formula (\Cref{prop:ch16-degree-formula})
 &model and continuous potentials
 &BFJ: \(\{\theta\}\) ample.  CLD: \(L\) ample and the product defined.
  For arbitrary unbounded potentials the mass may drop; finite-energy
  potentials retain full mass.  See row \textup{(L2)} and \Cref{def:e1}
 &\(=\)\\

(M3) \(\MA_\theta\) on continuous potentials
 (\Cref{prop:ch24-continuous-uniform})
 &weak continuity of the CLD product under uniform convergence
  (\Cref{thm:cld-bt-continuity})
 &continuous potentials, \eqref{eq:ch19-cld-bfj-continuous}
 &BFJ: uniform approximation by model functions.  CLD: local boundedness
  and uniform convergence.  Both are supplied by
  \Cref{thm:ch19-cld-bfj-comparison}\,(B), Step~3
 &\(=\)\\

(M4) monotone extension to bounded \(\theta\)-psh potentials
 (\Cref{thm:bounded-ma-extension})
 &Bedford--Taylor product of locally bounded CLD-psh functions
  (\Cref{thm:bt-symmetry-locality})
 &only where the potential is also continuous, i.e.\ row (M3)
 &The two mechanisms differ: BFJ uses decreasing \emph{nets} and
  \Cref{prop:dini-nets}; CLD needs uniform convergence
  (\Cref{cor:ch16-decreasing-nets}).  Membership of a bounded
  \(\theta\)-psh function in the CLD class is (C4)
 &\(?\)\\

\bottomrule
\end{tabular}
\end{center}

\noindent
The table continues.

\begin{center}
\footnotesize
\begin{tabular}{p{.185\textwidth}|p{.175\textwidth}|p{.155\textwidth}|p{.215\textwidth}|p{.10\textwidth}}
\toprule
\textbf{BFJ statement}&\textbf{CLD/GK statement}&\textbf{Locus of
 agreement}&\textbf{Hypotheses each side needs}&\textbf{Relation}\\
\midrule

(M5) nonpluripolar product on \(\cE^1\)
 (\Cref{def:nonpluripolar-ma}, \Cref{thm:e1-mixed-products})
 &no general canonical CLD product is constructed on the entire
  \(\cE^1\) class (\Cref{lem:continuous-times-positive-current}
  fails for unbounded \(u\))
 &special subclasses only
 &Row \textup{(L1)}: the recursion needs \(uT\), undefined for unbounded
  \(u\).  Where a canonical CLD extension does exist it may charge a polar
  set, whereas the nonpluripolar product discards that mass;
  \Cref{ex:ch16-polar-failure},
  \Cref{thm:ch19-cld-bfj-comparison}\,\cref{itm:ch19-cb:scope3}
 &no general equality\\

(M6) \(\ddc\) of a model function (\Cref{def:ch20-ddc}); on a curve,
 the graph Laplacian (\Cref{prop:ch20-curve-laplacian},
 \eqref{eq:ch31-graph-laplacian})
 &Poincar\'e--Lelong (\Cref{thm:cld-poincare-lelong}) and the slope
  formula (\Cref{thm:ch08-slope-formula},
  \Cref{thm:ch19-curve-slope-formula})
 &model functions; \(\mathbf R\)-PL weights
 &BFJ: the sign convention \eqref{eq:ch20-sign}.  CLD: the lattice
  orientation \eqref{eq:appJ-lagerberg-orientation}.  Both are fixed in
  \Cref{app:normalizations}
 &\(=\)\\

(M7) \(\MA_\omega(u)=\mu_0+\Delta_\Gamma u\) on a curve
 (\Cref{lem:ch31-curve-ma-affine})
 &\(c_1(\overline L)=\mu_0+\Delta_\Gamma F\)
  (\Cref{prop:ch19-curve-model-curvature})
 &continuous semipositive metrics on a curve
 &BFJ: a regular semistable model.  CLD: a formal model, no
  discreteness.  Transfer:
  \Cref{cor:ch19-transfer}\,\cref{itm:ch19-tr:curve}
 &\(=\)\\

\bottomrule
\end{tabular}
\end{center}

\noindent
The table continues.

\begin{center}
\footnotesize
\begin{tabular}{p{.185\textwidth}|p{.175\textwidth}|p{.155\textwidth}|p{.215\textwidth}|p{.10\textwidth}}
\toprule
\textbf{BFJ statement}&\textbf{CLD/GK statement}&\textbf{Locus of
 agreement}&\textbf{Hypotheses each side needs}&\textbf{Relation}\\
\midrule

(M8) \(\trop_*\MA_L(F)=n!\MA_{\mathbf R}(F)\)
 (\Cref{thm:toric-bfj-alexandrov}, \eqref{eq:toric-bfj-alexandrov})
 &\(\trop_*c_1(L,\|\cdot\|_\phi)^n=n!\MA_{\mathbf R}(\phi)\)
  (\Cref{thm:ch19-toric-ma}, \eqref{eq:cld-toric-ma-comparison})
 &continuous invariant toric metrics
 &BFJ: discrete \(v\) and a toric model for the polyhedral case.  CLD: the
  locally psh-approximable class of \cite[Theorem~6.1.1]{BGJKK25}.  Both
  need invariance and the factor \(n!\)
 &\(=\)\\

(M9) canonical measure on a totally degenerate abelian variety,
 \eqref{eq:abelian-canonical-haar}
 &\(c_1(\overline L)^g\) is Haar of mass \((L^g)\)
  (\Cref{thm:ch19-canonical-ma})
 &translation-invariant metrics with split Raynaud data
 &BFJ: \eqref{eq:ch19-bfj-hypotheses}.  CLD: \(K\) algebraically closed for
  splitness.  Transfer:
  \Cref{cor:ch19-transfer}\,\cref{itm:ch19-tr:abelian}
 &\(=\)\\

(M10) no mass on \(\{\varphi=-\infty\}\) for bounded entries
 (\Cref{prop:bounded-ma-no-polar-mass})
 &negligibility of the polar locus
  (\Cref{thm:cld-polar-negligibility})
 &bounded, respectively locally approachable and locally bounded off the
  polar set
 &BFJ: the entries are bounded and \(\varphi\) is only integrated.  CLD:
  the weighted CLN estimate \eqref{eq:ch16-weighted-cln-main}.  The two
  answers differ for unbounded entries; row \textup{(L3)},
  \Cref{thm:ch19-cld-bfj-comparison}\,\cref{itm:ch19-cb:scope4}
 &\(=\)\\

(M11) comparison principle for bounded potentials
 (\Cref{cor:bounded-comparison-principle})
 &chartwise comparison principle
  (\Cref{thm:ch16-comparison-chart})
 &continuous potentials
 &BFJ: capacity and quasicontinuity,
  \Cref{thm:bounded-quasicontinuity}.  CLD: the subgradient inclusion
  \Cref{lem:ch16-subgradient-inclusion}, a statement about finite convex
  functions.  Row \textup{(L7)}
 &\(\Rightarrow\)\\

(M12) maximum principle (\Cref{prop:model-maximum-principle},
 \Cref{thm:bounded-maximum-principle})
 &maximum principle in dimension one only
  (\Cref{thm:ch15-max-onedim}, \Cref{cor:ch15-cld-curve-max})
 &continuous potentials on a curve
 &BFJ: global, all dimensions, for its own class.  CLD: open in dimension
  \(\geq2\); row \textup{(L6)}
 &\(\Rightarrow\)\\

(M13) uniform model estimate
 \eqref{eq:ch24-model-uniform-estimate}
 &Chern--Levine--Nirenberg (\Cref{thm:cld-cln}) and its weighted form
  (\Cref{thm:ch16-weighted-cln})
 &model and continuous potentials
 &Different mechanisms: BFJ bounds a global mass by \(2V\sum_i\|\cdot\|_
  \infty\); CLD bounds a local mass by a product of sup-norms, infinite for
  unbounded potentials.  Row \textup{(L14)}
 &\(?\)\\
\bottomrule
\end{tabular}
\end{center}

\subsection{Table III: global apparatus with no local counterpart}
\label{sec:ch31-table-global}

The rows of this table are the ones a reader is most likely to
misremember, because they are places where \Cref{part:bfj} proves something
substantial and \Cref{part:cld} says nothing.  The correct reading of a
\(\varnothing\) is that the question is not a CLD question, not that it is
open.

\begin{center}
\footnotesize
\begin{tabular}{p{.185\textwidth}|p{.175\textwidth}|p{.145\textwidth}|p{.225\textwidth}|p{.10\textwidth}}
\toprule
\textbf{BFJ statement}&\textbf{CLD/GK statement}&\textbf{Locus}
 &\textbf{Why there is no counterpart}&\textbf{Relation}\\
\midrule

(G1) compactness modulo constants
 (\Cref{thm:bfj-psh-compactness})
 &none
 &---
 &Compactness is a statement about a \emph{global} class attached to a
  fixed \(\theta\); the CLD class is a sheaf with no global reference
 &\(\varnothing\)\\

(G2) decreasing regularization by model functions
 (\Cref{thm:bfj-decreasing-regularization})
 &collapse of the transfinite psh hierarchy
  (\Cref{prop:ch15-hierarchy-stabilizes})
 &---
 &Different questions: BFJ approximates from above by model functions on a
  net; CLD asks whether the hierarchy generated from below by smooth psh
  functions collapses at stage one.  Row \textup{(L5)},
  \Cref{prob:frontier-cld-hierarchy}
 &\(?\)\\

(G3) continuity of envelopes
 (\Cref{thm:bfj-envelope-regularity},
  \Cref{thm:ch23-gjkm-positive-characteristic})
 &none
 &---
 &An envelope is defined by a supremum over a global class; there is no
  local formulation
 &\(\varnothing\)\\

(G4) capacity and quasicontinuity
 (\Cref{def:bfj-capacity},
  \Cref{thm:bounded-quasicontinuity},
  \Cref{thm:e1-quasicontinuity})
 &none
 &---
 &The capacity \eqref{eq:capacity-definition} is a supremum of global
  masses; no local analogue is constructed in \Cref{part:cld}
 &\(\varnothing\)\\

(G5) energy and the class \(\cE^1\)
 (\Cref{def:model-energy}, \Cref{def:e1},
  \Cref{thm:ch26-e1-structure})
 &none
 &---
 &Energy is an integral over all of \(X^{\An}\) against a measure of total
  mass \(V\); it is meaningless chartwise
 &\(\varnothing\)\\

\bottomrule
\end{tabular}
\end{center}

\noindent
The table continues.

\begin{center}
\footnotesize
\begin{tabular}{p{.185\textwidth}|p{.175\textwidth}|p{.145\textwidth}|p{.225\textwidth}|p{.10\textwidth}}
\toprule
\textbf{BFJ statement}&\textbf{CLD/GK statement}&\textbf{Locus}
 &\textbf{Why there is no counterpart}&\textbf{Relation}\\
\midrule

(G6) orthogonality and differentiability
 (\Cref{def:orthogonality-property},
  \Cref{thm:energy-envelope-differentiability})
 &none
 &---
 &Both statements are variational and concern the global functional of
  (G5)
 &\(\varnothing\)\\

(G7) solvability of \(\MA_\omega(\varphi)=\mu\)
 (\Cref{thm:bfj-supported-cy},
  \Cref{cor:bfj-original-continuous-solution})
 &none
 &---
 &Row \textup{(L11)}: solvability is proved by global variational methods
  requiring energy, ampleness and, originally, discreteness and residue
  characteristic zero.  No local CLD existence theorem is available;
  \Cref{prob:frontier-general-ma}
 &\(\varnothing\)\\

(G8) functoriality: finite pushforward and base change
 (\Cref{thm:finite-ma-projection},
  \Cref{thm:base-change-ma})
 &pullback of psh functions and proper pushforward of currents
  (\Cref{cor:ch15-pullback-cases})
 &model and continuous metrics
 &Agreement on the common locus; but pullback of an \emph{arbitrary}
  current through a tropicalization is not asserted on either side,
  \Cref{sec:ch72-resolved-current-residue}
 &\(=\) on (M1), (M3); \(\varnothing\) beyond\\

(G9) a global \(\theta\) exists for a given local psh datum
 &---
 &none
 &This is the converse direction of (C4) and is
  \Cref{prob:frontier-cld-global-metric}; row \textup{(L10)}
 &\(?\)\\
\bottomrule
\end{tabular}
\end{center}

\subsection{Consistency with \texorpdfstring{\Cref{chap:cld-examples}}{Chapter 19}}
\label{sec:ch31-table-consistency}

The tables are a refinement of two statements already proved, and it is
worth checking that they refine and do not contradict.

\emph{The four unproved statements.}
\Cref{thm:ch19-cld-bfj-comparison}\,(C) declines to assert four things.
Each appears above, marked \(?\) or \(\varnothing\), and nowhere else:
\cref{itm:ch19-cb:scope1} and \cref{itm:ch19-cb:scope2} are the two halves of
row (C4), and the second reappears as (G9);
\cref{itm:ch19-cb:scope3} is row (M5); and \cref{itm:ch19-cb:scope4} is the
last clause of row (M10).  In particular no row of
\Cref{sec:ch31-table-classes} asserts an equality of psh \emph{classes}
outside (C1), (C2), (C6) and the interior half of (C7), which are
exactly the loci on which \Cref{chap:cld-examples} proves one.

\emph{The fourteen limitations.}  Every row of
\Cref{sec:ch19-limitations-table} that concerns an object of
\Cref{part:bfj} is cited above: \textup{(L1)} in (M5), \textup{(L2)} in
(M2), \textup{(L3)} in (M10), \textup{(L4)} in (C3), \textup{(L5)} in (G2),
\textup{(L6)} in (M12), \textup{(L7)} in (M11), \textup{(L10)} in (C4) and
(G9), \textup{(L11)} in (G7), \textup{(L14)} in (M13).  The remaining rows
\textup{(L8)}, \textup{(L9)}, \textup{(L12)} and \textup{(L13)} concern
restriction to a subvariety, ordinary delta-metricity, and local
regularization; they are internal to \Cref{part:cld} and have no BFJ
counterpart to compare, which is why they are absent above rather than
marked \(\varnothing\).

\emph{A census.}  The three tables have \(7+13+9=29\) rows.  Eleven are
marked \(=\) outright: (C1), (C2), (C6) and
(M1), (M2), (M3), (M6), (M7), (M8), (M9), (M10).  Two more are equalities
only on a sub-locus: (C7) on the relative interior and (G8) on the model and
continuous loci.  Four are one-way implications, five are unknown, and
seven have no counterpart.  Nineteen of the twenty-nine name a model or a
continuous class as their locus; the ten that do not are exactly the rows
where one of the two theories has nothing to say or where the question is
open.

\emph{The direction of the arrows.}  Four rows carry
\(\Rightarrow\) rather than \(=\): (C3), (C5), (M11), (M12).  In each the
BFJ statement is the stronger one.  For (C3), (M11), and (M12), the
obstruction to the converse is identified in
\Cref{sec:ch19-limitations-table} --- rows \textup{(L4)}, \textup{(L7)},
and \textup{(L6)}, respectively.  In (C5),
\cref{thm:p063-current-psh-ordinary-regularization} supplies the local
current-positive-to-first-stage implication on continuous smooth proper
curves, but not the reverse local-to-global BFJ numerical implication.  No
row carries
\(\Leftarrow\).  That asymmetry is not an artefact of exposition: the BFJ
theory is global and its statements quantify over all of \(X^{\An}\), so
where both theories speak, the global statement implies the local one and
not conversely.

\begin{remark}[How to use a row]
\label{rem:ch31-how-to-use}
Suppose one wishes to transport a statement proved in \Cref{part:cld} to the
BFJ setting.  Find the row whose CLD column contains it.  If the relation is
\(=\), the transport is legitimate on the stated locus once the hypotheses
in the fourth column are verified.  If it is \(\Rightarrow\), the transport
runs the wrong way and one must find a different argument.  If it is
\(?\) or \(\varnothing\), the transport is not available at all, and the
correct response is to add a hypothesis placing the object on a locus where
some row applies --- most often continuity, which activates (C2), (M2),
(M3), (M7).  This is the discipline that \Cref{chap:master-dictionary} will
apply to six theories at once.
\end{remark}

\subsection{The synthesis theorem}
\label{sec:ch31-synthesis-theorem}

The rows marked \(=\) in the model and continuous loci are all instances of
one statement, which we now prove in the form in which it will be quoted.

\begin{theorem}
\label{thm:bfj-cld-common-locus-synthesis}
Let \(L_1,\ldots,L_n\) carry semipositive model metrics.  The BFJ mixed
measure, the Chambert--Loir mixed measure, and the CLD/GK product of the
corresponding formal first Chern currents coincide.  The equality persists
for continuous semipositive metrics obtained as uniform limits whenever the
CLD product is defined by local uniform Bedford--Taylor continuity.
\end{theorem}

\begin{proof}
\emph{Step 1: a common model.}  Finitely many model metrics are determined
on finitely many models, and any two models of \((X,L)\) are dominated by a
third (\Cref{cor:ch20-limit-structure}).  Replacing all data by their
pullbacks to a common model \(\cX\), we may assume that every
\(\|\cdot\|_i\) is the model metric of a relatively nef model
\(\cL_i\) on \(\cX\), and that \(\cX\) is normal, flat and projective over
\(R\).  Nefness is preserved by pullback (\Cref{prop:nef-pullback}), and the
measures are unchanged (\Cref{prop:model-ma-basic}), so nothing is lost.

\emph{Step 2: the three coefficients.}  Write \(\cX_s=\sum_Eb_EE\).  The BFJ
mixed measure is, by \eqref{eq:model-mixed-measure},
\[
 \sum_Eb_E\bigl(c_1(\cL_1)|_E\cdots c_1(\cL_n)|_E\bigr)\,\delta_{x_E},
\]
and each coefficient is an intersection number on the projective
\(k\)-variety \(E\).  The Chambert--Loir measure is defined by the same
formula (\Cref{sec:metrics-cl-measure}, \eqref{eq:appJ-cl-measure}), the two
divisorial normalizations agreeing by \Cref{rem:divisorial-normalization};
this is \Cref{cor:ch24-cl-identification}.  The CLD Bedford--Taylor product
of the corresponding formal metrics is computed by
\Cref{thm:cld-formal-intersection} and identified with the same measure by
\Cref{thm:cld-recovers-cl}, whose hypotheses --- \(K\) discretely valued,
\(X\) proper reduced, \(\cX\) relatively normal --- hold by Step~1 and
\eqref{eq:ch19-bfj-hypotheses}.  Finally, that the delta-form product
\(c_1^\delta(\cL_1)\wedge\cdots\wedge c_1^\delta(\cL_n)\) is again this
measure is \Cref{thm:gk-delta-cl-measure}, proved in \cite{GK17,GK19}; the
comparison of formal metrics with the CLD calculus of \cite{CLD25} is the
content of \Cref{thm:common-model-core}.  Hence all four coefficients agree
at every \(E\), and the four measures are equal.  Their common total mass is
\((\{\theta_1\}\cdots\{\theta_n\})\) by \eqref{eq:ch24-mass}.

\emph{Step 3: continuous metrics.}  Let \(\|\cdot\|_i\) be continuous
semipositive, i.e.\ uniform limits of semipositive model metrics.  Choose,
for each \(i\), a sequence \(\varphi_{i,m}\) of \(\theta_i\)-psh model
functions with \(\|\varphi_{i,m}-\varphi_i\|_\infty\to0\); after the shift
of Step~3 of the proof of \Cref{thm:ch19-cld-bfj-comparison} we may take the
convergence to be uniform and decreasing.  On the BFJ side,
\Cref{prop:ch24-continuous-uniform} gives
\[
 \bigl|\MA(\varphi_{1,m},\ldots,\varphi_{n,m})(f)
      -\MA(\varphi_1,\ldots,\varphi_n)(f)\bigr|
 \longrightarrow0
\]
for every \(f\in C^0(X^{\An})\), the quantitative form being
\eqref{eq:ch24-model-uniform-estimate}.  On the CLD side,
\Cref{thm:cld-bt-continuity} gives weak convergence of the products, its
hypotheses --- local boundedness of the potentials and local uniform
convergence --- being satisfied because each \(\varphi_{i,m}\) is a
continuous psh function and the convergence is uniform on all of
\(X^{\An}\).  By Step~2 the two sequences of measures agree term by term,
hence so do their limits.  This is \eqref{eq:ch19-cld-bfj-continuous} in the
mixed form, and it is \Cref{thm:bfj-cld-ma-common-locus}.
\end{proof}

\begin{corollary}
\label{cor:ch31-four-families-transfer}
Under \eqref{eq:ch19-bfj-hypotheses}, in each of the following situations
the BFJ measure, the Chambert--Loir measure and the CLD Bedford--Taylor
product coincide, and the stated value is their common value.  When the
metric is formal/model, or more generally is a delta-metric for which the
delta-form product is defined, that product has the same value as well.
\begin{enumerate}
\item\label{itm:ch31-tr:curve} \(\dim X=1\) and the metric is continuous
 semipositive: the value is \(\mu_0+\Delta_\Gamma u\), of mass
 \(\deg L\) (\Cref{lem:ch31-curve-ma-affine},
 \Cref{prop:graph-ma-solution}).
\item\label{itm:ch31-tr:toric} \(X\) toric and the metric invariant
 continuous semipositive: the value pushes forward to
 \(n!\MA_{\mathbf R}(F)\), of mass \(n!\vol_M(P)=(L^n)\)
 (\Cref{thm:toric-bfj-alexandrov}).
\item\label{itm:ch31-tr:divisorial} the potential is the finite-pole
 envelope \(\Phi_{S,t}\) of \eqref{eq:finite-pole-envelope}: the value is
 \(\sum_ia_i\delta_{x_i}\) with \(\sum_ia_i=V\)
 (\Cref{prop:finite-pole-envelope}).
\item\label{itm:ch31-tr:abelian} \(X=A\) totally degenerate abelian with
 split Raynaud data and the metric canonical: the value is
 \((L^g)\,m_{\mathrm{Haar}}\) on \(\Sigma_A\)
 (\eqref{eq:abelian-canonical-haar},
 \Cref{thm:ch19-canonical-ma}).
\end{enumerate}
\end{corollary}

\begin{proof}
In each case the metric is continuous semipositive, so
\Cref{thm:bfj-cld-common-locus-synthesis} applies once that is checked.
For \cref{itm:ch31-tr:curve} it is
\Cref{prop:graph-ma-solution}; for \cref{itm:ch31-tr:toric} it is Step~4 of
the proof of \Cref{thm:toric-bfj-alexandrov}; for
\cref{itm:ch31-tr:divisorial} it is the continuity clause of
\Cref{prop:finite-pole-envelope}, which rests on
\Cref{thm:bfj-envelope-regularity}; and for \cref{itm:ch31-tr:abelian} it is
\eqref{eq:ch19-tate-model-approx} in dimension one and the piecewise affine
interpolation of \(\tfrac12q_L\) in general, as in
\Cref{cor:ch19-transfer}\,\cref{itm:ch19-tr:abelian}.  The stated values are
computed in the cited places.
\end{proof}

\begin{remark}[An equality of constructions, not of categories]
\label{rem:ch31-not-an-equivalence}
The theorem is an equality of constructions on an overlap.  It is not an
equivalence of categories.  BFJ begins with models and global positivity;
CLD begins with tropical charts and a local differential calculus; the
delta-form theory enlarges the latter so that model curvature and tropical
cycles are represented internally.  Their agreement on formal metrics is
the hinge connecting the approaches, and every \(=\) in
\Cref{sec:ch31-table-main} is downstream of that hinge.  The rows marked
\(\varnothing\) in \Cref{sec:ch31-table-global} are not gaps in the hinge;
they are questions that only one of the two theories poses.
\end{remark}

\subsection{The descriptive ledger}
\label{sec:ch31-ledger}

For readers who prefer prose to a grid, the following ledger states the same
information one class of object at a time.  ``Agreement'' includes the
multiplicity, lattice, and total-mass normalizations fixed in
Appendix~J.

\begin{description}[style=nextline,leftmargin=0pt,labelindent=0pt,
  labelsep=0pt,itemsep=.8\baselineskip,font=\normalfont\bfseries]
\item[Semipositive model metrics]
\emph{Agreement.} The BFJ intersection measure, Chambert--Loir measure, and
the formal Chern delta-current product of CLD/GK agree coefficient by
coefficient.
\par\emph{Boundary.} Equality of the ambient psh classes is not asserted.

\item[Continuous semipositive metrics]
\emph{Agreement.} BFJ and Chambert--Loir measures agree by uniform model
approximation.  If the metric is locally CLD-psh and the CLD Bedford--Taylor
product is defined through the same approximation, that product agrees as
well.
\par\emph{Boundary.} It is not asserted that every continuous BFJ-psh metric
is locally approachable in an arbitrary CLD chart.

\item[Rational polyhedral tropical charts]
\emph{Agreement.} Vertical intersections equal calibrated tropical
intersections, and Poincar\'e--Lelong agrees.
\par\emph{Boundary.} Pullback of arbitrary currents through a
tropicalization is not asserted.

\item[Smooth convex potential on an open skeleton face]
\emph{Agreement.} The non-Archimedean measure equals the real Alexandrov
measure with the residue-degree and multiplicity factor of \cite{Vil20}.
\par\emph{Boundary.} No global formula across all boundary faces is asserted
without balancing data.

\item[Toric invariant metrics]
\emph{Agreement.} BFJ, CLD or tropical, and real Alexandrov measures agree as
in \eqref{eq:toric-bfj-alexandrov}.
\par\emph{Boundary.} Noninvariant metrics are not determined by
tropicalization.

\item[Bounded global BFJ psh potentials]
\emph{Agreement.} The BFJ bounded Monge--Amp\`ere product is defined.
\par\emph{Boundary.} Automatic identification with the local CLD psh class
or its Bedford--Taylor product is not asserted.

\item[Finite-energy BFJ potentials]
\emph{Agreement.} The global nonpluripolar product is defined and has full
mass.
\par\emph{Boundary.} No general CLD current wedge for unbounded potentials is
asserted.

\item[CLD smooth forms and currents]
\emph{Agreement.} Local differential calculus, Stokes,
Poincar\'e--Lelong, and proper pushforward are available.
\par\emph{Boundary.} This does not yield a general global BFJ potential or a
pullback of arbitrary currents.

\item[Harmonic tropicalization forms]
\emph{Agreement.} These are functorial local forms in the sense of
\cite{GJR26}.
\par\emph{Boundary.} Not every such form is identified with the curvature of
a global metric.
\end{description}

Each entry corresponds to one or more rows of the tables: the first to (C1)
and (M1); the second to (C2), (C3), (M2), (M3); the third to (M6); the
fourth and fifth to (C6) and (M8); the sixth to (C4) and (M4); the seventh
to (M5) and (G5); the eighth to (G8); and the ninth to row \textup{(L9)} of
\Cref{sec:ch19-limitations-table}, which has no BFJ counterpart.

\subsection{What the tables do not say}
\label{sec:ch31-table-not}

Four summary statements are \emph{not} consequences of anything above, and
each has been asserted somewhere in the literature or in conversation.

\begin{enumerate}
\item\label{itm:ch31-not:classes}
 \emph{That the BFJ and CLD psh classes coincide.}  They are compared only
 on the model and continuous loci, rows (C1), (C2), (C5), (C6).  Row (C4)
 is \(?\) and is \Cref{prob:frontier-cld-global-metric}.
\item\label{itm:ch31-not:operators}
 \emph{That the two Monge--Amp\`ere operators coincide on their whole
 domains.}  They coincide on (M1) and (M3).  Row (M4) is \(?\) and row
 (M5) is \(\varnothing\).
\item\label{itm:ch31-not:local}
 \emph{That a local CLD statement can be globalized by patching.}  Nothing
 in \Cref{part:cld} produces a global class; that is precisely the content
 of (G9).
\item\label{itm:ch31-not:skeleton}
 \emph{That a metric is determined by its restriction to a skeleton.}
 False already for a model metric on \(\mathbf P^1\), by
 \Cref{ex:ch31-noninvariant}; true after \(\trop_*\) only in dimension one,
 by \Cref{rem:ch31-retraction}.
\end{enumerate}

\begin{warning}{The tables will grow, and they will not shrink}
\Cref{chap:master-dictionary} extends this comparison to the trivially
valued theory, the complex theory, the Monge--Amp\`ere-psh theory of
\Cref{part:complex}, and the model theory of \Cref{part:foundations}.  Rows
may be added, and a \(?\) may become an \(=\) if a theorem is proved.  What
may not happen is that an \(=\) be asserted on a larger locus without a
proof: the discipline of \Cref{sec:ch19-limitations} is that the boundary of
each statement is part of the statement.
\end{warning}

\section{A final diagnostic list}
\label{sec:ch31-diagnostic}

Before using a Monge--Amp\`ere identity, check the following in order:
\begin{enumerate}
\item\label{itm:ch31-diag:class} Is the class ample, or merely nef and big
      after a birational pullback?
\item\label{itm:ch31-diag:potential} Is the potential model, continuous
      semipositive, bounded BFJ-psh, or finite energy?
\item\label{itm:ch31-diag:mass} Is the measure normalized to mass \(V\) or to
      probability?
\item\label{itm:ch31-diag:cld} Does a local CLD product exist for the chosen
      potentials?
\item\label{itm:ch31-diag:skeleton} If a skeleton appears, is the metric
      genuinely determined by its restriction, or only its prescribed measure
      supported there?
\item\label{itm:ch31-diag:field} Under field extension, has the absolute
      value been extended or renormalized?
\end{enumerate}

On curves all six questions collapse to graph potential theory.  On toric
varieties they collapse, in the invariant category, to convex analysis.  In
general they do not collapse, which is precisely why the model, CLD, BFJ,
and finite-energy layers were constructed separately.

Each question has an address in this chapter, and each has an example in
which the wrong answer produces a false statement.

\begin{center}
\small
\begin{tabular}{p{.06\textwidth}|p{.30\textwidth}|p{.26\textwidth}|p{.28\textwidth}}
\toprule
&\textbf{Where the answer is decided}&\textbf{What goes wrong}
 &\textbf{Example}\\
\midrule
\cref{itm:ch31-diag:class}
 &\eqref{eq:ch19-bfj-hypotheses}; the volume
  \eqref{eq:ch24-volume}
 &\(V=0\), so the normalization is meaningless
 &\Cref{ex:ch31-nef-not-big}\\
\cref{itm:ch31-diag:potential}
 &the four rows of \Cref{sec:ch24-four-products}
 &a continuity theorem is applied to a potential in the wrong row
 &\Cref{ex:ch31-chain-lebesgue},
  \Cref{rem:ch31-tate-not-model}\\
\cref{itm:ch31-diag:mass}
 &\eqref{eq:examples-global-normalization},
  \eqref{eq:appJ-prob-ma}
 &a factor of \(V\) or of \(n!\) is lost
 &\Cref{ex:ch31-toric-atom-pn},
  \Cref{ex:ch31-toric-logsumexp}\\
\cref{itm:ch31-diag:cld}
 &\Cref{sec:ch31-table-main}, rows (M4), (M5)
 &an equality is asserted where the table says \(?\) or
  \(\varnothing\)
 &\Cref{sec:ch31-table-not}\,\cref{itm:ch31-not:operators}\\
\cref{itm:ch31-diag:skeleton}
 &\Cref{rem:ch31-retraction}
 &a metric is reconstructed from its restriction
 &\Cref{ex:ch31-noninvariant}\\
\cref{itm:ch31-diag:field}
 &\Cref{sec:ch31-base-change}
 &an energy is quoted as base-change invariant
 &\Cref{ex:ch31-base-change}\\
\bottomrule
\end{tabular}
\end{center}

\begin{warning}{The commonest single error}
It is question~\cref{itm:ch31-diag:mass}.  Of the four families of this
chapter, three produce measures whose total mass is an intersection number
greater than one --- \(\deg L\), \((L^n)\), \((L^g)\) --- and the fourth,
the abelian case, is usually written with a probability Haar measure and an
explicit factor \((L^g)\) in front.  A statement quoted from one source and
combined with a statement quoted from another will silently mix the two
normalizations unless the masses are checked.  That is why every example
above ends with a mass.
\end{warning}

\section{Synthesis of Part III}
\label{sec:ch31-synthesis}

\subsection{The dependency chain, and where it can be entered}
\label{sec:ch31-dependency}

\Cref{part:bfj} is a linear construction with three genuinely
non-formal steps.  The following is not a summary of contents but a record
of what depends on what, so that a reader who wants only one theorem knows
what must be believed.

\emph{The linear layer.}  \Cref{chap:bfj-model-forms} constructs
\(\mathcal Z^{1,1}(X)\) and \(\ddc\), and proves the exactness
\Cref{thm:ch20-ddc-exact}; its algebraic input is relative intersection
theory and the relative Hodge index theorem
\Cref{thm:ch20-mixed-hodge}.  \Cref{chap:bfj-psh} defines
\(\PSH(X,\theta)\) and proves the closure properties
\Cref{thm:ch21-stability}; nothing here is hard.  This layer suffices for
the model computations of \Cref{sec:ch31-curves} and
\Cref{sec:ch31-toric-polyhedral}.

\emph{The compactness layer.}  \Cref{chap:bfj-compactness} proves
equicontinuity on dual complexes
(\Cref{thm:bfj-skeletal-equicontinuity}), compactness modulo constants
(\Cref{thm:bfj-psh-compactness}), and decreasing regularization
(\Cref{thm:bfj-decreasing-regularization}).  The last uses multiplier
ideals and is where residue characteristic zero first bites.  Everything
about bounded potentials rests on it.

\emph{The envelope layer.}  \Cref{chap:bfj-envelopes} proves continuity of
envelopes (\Cref{thm:bfj-envelope-regularity}), the second place where
residue characteristic zero is used; \Cref{thm:ch23-gjkm-positive-characteristic}
replaces it in positive characteristic at the cost of test ideals.  Every
statement about \(\Phi_{S,t}\) in \Cref{sec:ch31-divisorial} rests on this,
as item \cref{itm:ch31-pole:envelope} records.

\emph{The measure layer.}  \Cref{chap:bfj-ma} builds \(\MA_\theta\) in three
incarnations; \Cref{chap:bfj-capacity} supplies the capacity and
quasicontinuity that make locality statements true for bounded potentials;
\Cref{chap:bfj-energy} supplies \(E\), \(\cE^1\), and the nonpluripolar
product.  \Cref{chap:bfj-differentiability} couples orthogonality to
differentiability, and \Cref{chap:bfj-finite-energy-measures} solves the
equation variationally.  \Cref{chap:bfj-continuity} upgrades the solution to
a continuous one under support hypotheses;
\Cref{chap:bfj-functoriality} makes the whole package functorial.

\emph{Three non-formal steps.}  They deserve to be named, because a reader
who accepts them may take the rest on faith.
\begin{enumerate}
\item\label{itm:ch31-hard:hodge}
 The relative Hodge index theorem \Cref{thm:ch20-mixed-hodge}, which makes
 the model Dirichlet form negative
 (\eqref{eq:model-dirichlet-negativity}) and the energy concave.  In this
 chapter it is visible as the sign of \(-\tfrac12\int|u'|^2\rd s\) in
 \Cref{prop:ch26-graph-energy}, verified in every curve example.
\item\label{itm:ch31-hard:envelope}
 Continuity of envelopes, \Cref{thm:bfj-envelope-regularity}.  In this
 chapter it is what makes \(\Phi_{S,t}\) a legitimate object.
\item\label{itm:ch31-hard:monotone}
 The monotone extension \Cref{thm:bounded-ma-extension}, which is the only
 construction in \Cref{part:bfj} that produces a measure not already given
 by an intersection number.  No example of this chapter requires it: all
 four families are continuous, and continuity suffices by
 \Cref{prop:ch24-continuous-uniform}.
\end{enumerate}
The third item is worth dwelling on.  The examples of a chapter like this
one can never exercise the hardest theorem of a part, because computable
objects are regular.  The reader should not conclude that
\Cref{thm:bounded-ma-extension} is decorative; it is what makes
\Cref{chap:bfj-energy} and hence \Cref{chap:bfj-finite-energy-measures}
possible, and \Cref{ex:ch24-branch-sum} is the example that exercises it.

\subsection{Where each hypothesis was used}
\label{sec:ch31-hypothesis-use}

\begin{center}
\small
\begin{tabular}{p{.20\textwidth}|p{.34\textwidth}|p{.36\textwidth}}
\toprule
\textbf{Hypothesis}&\textbf{What it buys}&\textbf{What happens without it}\\
\midrule
\(K\) discretely valued, \(v(\pi)=1\)
 &models over \(R\), the multiplicity \(b_E\), integral skeletal lengths
 &the CLD side survives with a defect factor
  (\Cref{rem:ch19-cld-bfj-hypotheses}); the BFJ side has no model theory in
  the form used here\\
residue characteristic \(0\)
 &decreasing regularization
  (\Cref{thm:bfj-decreasing-regularization}) and envelope continuity
  (\Cref{thm:bfj-envelope-regularity})
 &replaced by test ideals,
  \Cref{thm:ch23-gjkm-positive-characteristic},
  \Cref{thm:gjkm-positive-characteristic-cy}; mixed characteristic is
  \Cref{thm:fgk-mixed-characteristic-cy}\\
\(X\) smooth projective
 &resolution of singularities, SNC models
 &statement (A) of \Cref{thm:ch19-cld-bfj-comparison} survives for proper
  reduced \(X\); the envelope machinery does not\\
\(\{\theta\}\) ample
 &\(V>0\), a nondegenerate capacity, a nontrivial variational problem
 &\Cref{ex:ch31-nef-not-big}: everything degenerates\\
torus invariance (toric family)
 &positivity \(=\) convexity, \(\MA=n!\MA_{\mathbf R}\)
 &\Cref{ex:ch31-noninvariant}: positivity is not detected by the
  restriction, and the measure is not determined by it\\
\(x_i\) divisorial on a dual complex
 &the compactness estimate
  \Cref{thm:ch22-proper-normalization}
 &\Cref{ex:ch31-type-one-fails}: the envelope is \(+\infty\)\\
\(\mu\) supported on a dual complex
 &continuity of the solution
  (\Cref{thm:bfj-supported-regularity})
 &solvability remains, in finite energy
  (\Cref{thm:bfj-supported-cy}); continuity is not asserted without
  the support hypothesis, including on curves\\
\bottomrule
\end{tabular}
\end{center}

\subsection{A mass ledger for the four families}
\label{sec:ch31-mass-ledger}

Every measure computed in this chapter, with its total mass and the
intersection number the mass is checked against.  This table is the
\Cref{part:bfj} counterpart of the ledger \Cref{sec:ch08-mass-ledger} and of
\Cref{sec:ch19-mass-ledger}.

\begin{center}
\footnotesize
\begin{tabular}{p{.225\textwidth}|p{.225\textwidth}|p{.13\textwidth}|p{.13\textwidth}|p{.19\textwidth}}
\toprule
\textbf{Object}&\textbf{Measure}&\textbf{Mass}
 &\textbf{Checked against}&\textbf{Where}\\
\midrule
two-component curve, reference
 &\(m_0\delta_{x_0}+m_1\delta_{x_1}\)&\(m_0+m_1\)&\(\deg L\)
 &\eqref{eq:nodal-reference-measure}\\
two-component curve, \(a=-m_1\)
 &\(V\delta_{x_0}\)&\(V=m_0+m_1\)&\(\deg L\)
 &\Cref{ex:two-component-complete}\\
the same with \(m_0=2,m_1=3\)
 &\(5\delta_{x_0}\)&\(5\)&\(\deg L=5\)
 &\Cref{ex:ch31-two-component-numbers}\\
chain of four, \(\mu=2\delta_1\)
 &\(2\delta_{x_1}\)&\(2\)&\(\deg L=2\)
 &\Cref{ex:ch31-chain-solve}\\
chain of four, diffuse
 &\(\tfrac23\,\rd s\) on \([0,3]\)&\(2\)&\(\deg L=2\)
 &\Cref{ex:ch31-chain-lebesgue}\\
Tate curve, canonical
 &\(\tfrac d\ell\,\rd t\) on \(\mathbf R/\ell\mathbf Z\)&\(d\)
 &\(\deg L=d\)&\eqref{eq:ch31-tate-potential}\\
Tate curve, \(N\)-th model
 &\(\tfrac dN\sum_{i}\delta_{t_i}\)&\(d\)&\(\deg L=d\)
 &\eqref{eq:ch19-tate-model-approx}\\
\(\mathbf P^1\), blowup at a non-toric point
 &\(\delta_{x_E}\)&\(1\)&\(\deg\cO(1)=1\)
 &\eqref{eq:ch31-p1-two-measures}\\
\bottomrule
\end{tabular}
\end{center}

\begin{center}
\footnotesize
\begin{tabular}{p{.225\textwidth}|p{.225\textwidth}|p{.13\textwidth}|p{.13\textwidth}|p{.19\textwidth}}
\toprule
\textbf{Object}&\textbf{Measure}&\textbf{Mass}
 &\textbf{Checked against}&\textbf{Where}\\
\midrule
\(\mathbf P^1\), \(\cO(2)\), polyhedral
 &\(\delta_{\iota(0)}+\delta_{\iota(1)}\)&\(2\)&\(\deg\cO(2)=2\)
 &\Cref{ex:ch31-toric-p1-two-atoms}\\
\(\mathbf P^1\), logistic weight
 &\(p(1-p)\,\rd u\)&\(1\)&\(\deg\cO(1)=1\)
 &\eqref{eq:ch31-logistic-measure}\\
\(\mathbf P^n\), \(\cO(d)\), reflected support function
 &\(d^n\delta_{\xi_{\mathrm G}}\)&\(d^n\)&\((\cO(d)^n)=d^n\)
 &\Cref{ex:ch31-toric-atom-pn}\\
\(\mathbf P^2\), \(\cO(2)\), two atoms
 &\(2\delta_{\iota(0,1)}+2\delta_{\iota(1,0)}\)&\(4\)&\((\cO(2)^2)=4\)
 &\eqref{eq:ch31-p2-two-atoms}\\
\(\mathbf P^n\), \(\cO(d)\), log-sum-exp
 &\(n!d^{2n}p_0\cdots p_n\,\rd u\)&\(d^n\)&\((\cO(d)^n)=d^n\)
 &\eqref{eq:ch31-logsumexp-mass}\\
\(\mathbf P^1\!\times\!\mathbf P^1\), mixed
 &\(\delta_{x_{\mathrm G}}\)&\(1\)&\((H_1\cdot H_2)=1\)
 &\Cref{ex:ch31-toric-mixed-numbers}\\
\(\mathbf P^1\!\times\!\mathbf P^1\), \(\cO(1,0)\)
 &\(0\)&\(0\)&\((L^2)=0\)
 &\Cref{ex:ch31-nef-not-big}\\
\bottomrule
\end{tabular}
\end{center}

\begin{center}
\footnotesize
\begin{tabular}{p{.225\textwidth}|p{.225\textwidth}|p{.13\textwidth}|p{.13\textwidth}|p{.19\textwidth}}
\toprule
\textbf{Object}&\textbf{Measure}&\textbf{Mass}
 &\textbf{Checked against}&\textbf{Where}\\
\midrule
divisorial Green function on \(\mathbf P^1\)
 &\(\delta_{x_E}\)&\(1\)&\(V=1\)
 &\eqref{eq:ch31-p1-green-formula}\\
finite-pole envelope, general
 &\(\sum_ia_i\delta_{x_i}\)&\(V\)&\((\{\omega\}^n)\)
 &\Cref{prop:finite-pole-envelope}\\
two-pole programme, \(m_0=2,m_1=3\)
 &\((2+\delta)\delta_{x_0}+(3-\delta)\delta_{x_1}\)&\(5\)&\(\deg L=5\)
 &\eqref{eq:ch31-two-pole-masses}\\
good reduction
 &\(V\delta_{x_0}\)&\(V\)&\((L^n)\)
 &\eqref{eq:ch31-good-reduction-measure}\\
abelian, canonical
 &\((L^g)\,m_{\mathrm{Haar}}\)&\((L^g)\)&\(g!\det B\)
 &\eqref{eq:abelian-canonical-haar},
  \eqref{eq:ch31-abelian-mass}\\
abelian surface, \(B=\bigl(\begin{smallmatrix}2&1\\1&2\end{smallmatrix}\bigr)\)
 &\(6\,\rd x\) on \(\mathbf R^2/\mathbf Z^2\)&\(6\)&\((L^2)=6\)
 &\Cref{ex:ch31-abelian-surface-bfj}\\
\bottomrule
\end{tabular}
\end{center}

Twenty rows, twenty masses, twenty intersection numbers.  A reader who
reproduces the table has verified the normalization
\eqref{eq:examples-global-normalization} in every regime in which it is used
in this book.

\subsection{Numbers to remember}
\label{sec:ch31-numbers}

\begin{center}
\small
\begin{tabular}{p{.40\textwidth}|p{.20\textwidth}|p{.32\textwidth}}
\toprule
\textbf{Quantity}&\textbf{Value}&\textbf{Where}\\
\midrule
energy of the two-component solution
 &\(-m_0m_1-\tfrac12m_1^2\)
 &\eqref{eq:two-component-energy}\\
the same with \(m_0=2,m_1=3\)
 &\(-\tfrac{21}2\)
 &\Cref{ex:ch31-two-component-numbers}\\
energy of the chain solution for \(2\delta_0\)
 &\(-\tfrac92\)
 &\Cref{ex:ch31-chain-solve}\\
energy of the diffuse chain solution
 &\(-\tfrac12\)
 &\Cref{ex:ch31-chain-lebesgue}\\
canonical weight of a Tate curve
 &\(\tfrac{d}{2\ell}t(t-\ell)\)
 &\eqref{eq:ch31-tate-potential}\\
its minimum
 &\(-\tfrac{d\ell}{8}\)
 &\eqref{eq:ch31-tate-potential}\\
its energy
 &\(-\tfrac{\ell d^2}{24}\)
 &\eqref{eq:ch31-tate-energy}\\
energy of the \(N\)-th model approximation
 &\(-\tfrac{\ell d^2}{24}(1-N^{-2})\)
 &\eqref{eq:ch31-tate-energy-approx}\\
uniform error of that approximation
 &\(\tfrac{d\ell}{8N^2}\)
 &\Cref{ex:ch31-tate-bfj}\\
\bottomrule
\end{tabular}
\end{center}

\begin{center}
\small
\begin{tabular}{p{.40\textwidth}|p{.20\textwidth}|p{.32\textwidth}}
\toprule
\textbf{Quantity}&\textbf{Value}&\textbf{Where}\\
\midrule
energy of \(-\varphi_E\) on \(\mathbf P^1\)
 &\(-\tfrac12\)
 &\Cref{ex:ch31-p1-blowup}\\
energy of the divisorial Green function
 &\(+\tfrac12\)
 &\Cref{ex:ch31-p1-green}\\
sup of the Green function at distance \(m\)
 &\(m\)
 &\Cref{ex:ch31-type-one-fails}\\
its energy
 &\(\tfrac m2\)
 &\Cref{ex:ch31-type-one-fails}\\
lattice area of each subgradient triangle on \(\mathbf P^2\)
 &\(1\)
 &\Cref{ex:ch31-toric-p2-two-atoms}\\
mixed volume \(V(P_1,P_2)\) on \(\mathbf P^1\!\times\!\mathbf P^1\)
 &\(\tfrac12\)
 &\Cref{ex:ch31-toric-mixed-numbers}\\
the factor relating the two Monge--Amp\`ere operators
 &\(n!\)
 &\eqref{eq:toric-bfj-alexandrov}\\
number of rows of the compatibility tables
 &\(7+13+9=29\)
 &\Cref{sec:ch31-table}\\
of which marked \(=\) outright
 &\(12\)
 &\Cref{sec:ch31-table-consistency}\\
of which equal only on a sub-locus
 &\(2\), namely (C7), (G8)
 &\Cref{sec:ch31-table-consistency}\\
of which marked \(\Rightarrow\), \(?\), \(\varnothing\)
 &\(3\), \(5\), \(7\)
 &\Cref{sec:ch31-table-consistency}\\
\bottomrule
\end{tabular}
\end{center}

\subsection{What \texorpdfstring{\Cref{part:bfj}}{Part III} hands on}
\label{sec:ch31-handoff}

Three things leave this part.

\emph{To \Cref{part:bj}.}  The variational method --- energy, dual energy,
coercivity, orthogonality, differentiability --- is transplanted almost
verbatim to the trivially valued setting, with a fixed ground valuation
replaced by a homogeneous valuation cone.  What does \emph{not} transplant is
the model theory: over a trivially valued field the models are test
configurations and the scaling action is an extra structure, so
\Cref{chap:bfj-model-forms} has no direct analogue.  The envelope theorem of
\Cref{part:bj} is proved by different means.

\emph{To \Cref{part:comparisons}.}  The compatibility tables of
\Cref{sec:ch31-table} become four columns of the master dictionary
\Cref{thm:dict-master-diagram}, and the synthesis theorem
\Cref{thm:bfj-cld-common-locus-synthesis} becomes the model core
\Cref{prop:dict-model-core} of that diagram.  The continuous common locus
reappears as \Cref{cor:dict-continuous-common-locus} and the mixed
comparison as \Cref{thm:cma-continuous-comparison}.

\emph{To the reader.}  A working knowledge of four measures --- an atomic
measure on a metric graph, an Alexandrov measure on \(N_\mathbf R\), a
single atom of mass \(V\) at a divisorial point, and a Haar measure on a real
torus --- together with the discipline of checking the mass.  Everything
else in \Cref{part:bfj} is machinery for producing measures that are not on
that list.

\section*{Source notes}

\emph{Provenance of the examples.}  The curve and atomic calculations are
\cite[\S\S8--9]{BFJ15}; \Cref{prop:graph-ma-solution} and
\Cref{prop:finite-pole-envelope} are \cite[Lemma~8.5 and
Proposition~8.6]{BFJ15}, and the divisorial Green function
\eqref{eq:green-envelope-divisorial} is the object studied in
\cite{BJGreen24}.  Toric and invariant tropical pluripotential theory is
treated in \cite{BGJKK25}, whose Theorem~6.1.1 is the tropical form of
\eqref{eq:toric-bfj-alexandrov} and whose Theorems~6.3.13 and~8.3.3 solve
the invariant Monge--Amp\`ere equation; the local real/non-Archimedean
comparison is \cite{Vil20}.  The arithmetic toric formalism, in which
\(n!\vol_M(P)=(L^n)\) first appears in the form used here, is
\cite{BGPS14}.  Formal Chern currents and delta-forms are developed in
\cite{GK17,GK19}; harmonic tropicalizations are \cite{GJR26}; the local
calculus is \cite{CLD25}.  The Tate curve computation in the shape of
\eqref{eq:ch31-tate-potential} goes back to the classical theory of
non-Archimedean N\'eron functions and is recorded in the present
normalization in \Cref{ex:ch19-tate-canonical}.

\emph{Provenance of the general statements quoted.}  The monotone extension
\Cref{thm:bounded-ma-extension} is \cite[Theorem~3.1]{BFJ15}; the psh
formalism it uses is \cite{BFJ16}, with the survey \cite{BFJ16b} as a short
entry point.  Envelope continuity in residue characteristic zero is
\cite{BFJ15}; algebraizability is removed in \cite{BGGJKM20}, the
equal-characteristic test-ideal extension is \cite{GJKM19}, and the
conditional mixed-characteristic extension is \cite{FGK25}.  Continuity of
solutions for positive skeletal-supported measures of matching mass on
smooth projective geometrically connected curves with an ample line bundle
over arbitrary complete nontrivially valued fields is
\Cref{thm:curves-arbitrary-field-continuity}, proved in
\Cref{chap:bfj-continuity}, where its provenance is recorded.  The
comparison of algebraic, formal, and
piecewise linear semipositivity over an arbitrary complete field is
\cite{GM19}, and the original Chambert--Loir measure is \cite{CL06}, with
Zhang's semipositive metrics \cite{Zhang1995} as the source of the
continuous locus.  Wanner's theorem, used in row (C5), is
\cite{Wanner19}; the \(\mathbf R\)-PL comparison of row (C7) is
\cite{GR25Classical,GR25PL}.  Broader comparisons under semiampleness are
\cite{BE21}.  The trivially valued counterpart of the entire part is
\cite{BJ22}, and the synthetic treatment of finite-energy measures is
\cite{BJ25FiniteEnergy}.

\emph{The compatibility table.}  Tables comparing non-Archimedean
pluripotential theories exist in the literature only in prose, typically as
a paragraph of scope remarks attached to a comparison theorem.  The grid of
\Cref{sec:ch31-table} is assembled here from
\Cref{thm:ch19-cld-bfj-comparison} and
\Cref{sec:ch19-limitations-table}, and its arrangement --- one row per pair
of statements, with the locus, the hypotheses of each side, and the type of
relation as separate columns --- is a device of this book.  It contains no
mathematical content beyond those two sources; that is the point of it.  The
\(\varnothing\) rows of \Cref{sec:ch31-table-global} are the ones most worth
the reader's attention, since they record questions that one of the two
theories does not pose, and it is the confusion of \(\varnothing\) with
\(?\) that manufactures spurious open problems, as
\Cref{sec:ch19-limitations-reading} warns.

\emph{What is organized differently here.}  Three choices depart from the
sources.  First, the four families are carried to explicit numbers and each
is closed by a mass check against an intersection number; the sources
usually verify the mass once and thereafter suppress it, which is efficient
for a paper and fatal for a reader assembling statements from several
papers.  Second, the invariance hypothesis of the toric comparison is tested
by an explicit counterexample, \Cref{ex:ch31-noninvariant}, built from a
single blowup of \(\mathbf P^1_R\); we know of no place in the literature
where the hypothesis is shown to be necessary rather than merely convenient.
Third, the finite-pole envelope of \Cref{prop:finite-pole-envelope} is
presented as a finite-dimensional concave programme with Lagrange
multipliers, \Cref{ex:ch31-two-pole-programme}, so that the differentiability
theorem \Cref{thm:energy-envelope-differentiability} acquires a
two-dimensional shadow that can be checked by hand.

\emph{What is new here, and what is not.}  No theorem of this chapter is
new.  \Cref{lem:ch31-curve-ma-affine}, \Cref{lem:ch31-green-kernel},
\Cref{prop:ch31-one-pole} and \Cref{prop:ch31-toric-mixed} are elementary
consequences of statements proved earlier in the book, written out because
each is used in a form that is easy to misremember.
\Cref{thm:toric-bfj-alexandrov} is given a model-side proof, complementing
the tropical proof of \Cref{thm:ch19-toric-ma}; the two are genuinely
different arguments for the same identity, and their agreement is a
normalization check rather than a new theorem.  Among the examples,
\Cref{ex:ch31-p1-blowup}, \Cref{ex:ch31-noninvariant},
\Cref{ex:ch31-type-one-fails} and \Cref{ex:ch31-two-pole-programme} have no
counterpart in the sources known to us; the first three exist to show that a
hypothesis cannot be dropped, and the fourth to make a variational statement
finite-dimensional.  The remaining examples are computations of statements
in \cite{BFJ15,BGJKK25,GJR26} carried to numbers.
